\newMonth{December}
\setrunningtitles{December}{December}

% ---- martyrology/mart12/mart1201.htm
\needspace{10\baselineskip}
\begin{paracol}{2}
\selectlanguage{latin}
\begin{center}{\color{gregoriocolor} Kaléndis Decémbris. 
 Luna\dots\ }\end{center}
\switchcolumn
\selectlanguage{english}
\begin{center}{\color{gregoriocolor} The 
 1\textsuperscript{st} Day of 
 December. The\dots\ Day of the Moon.}\end{center}
\end{paracol}

\noindent\begin{tabularx}{\linewidth}{*{19}{>{\centering\arraybackslash}X}}
 \textcolor{gregoriocolor}{a} & \textcolor{gregoriocolor}{b} & \textcolor{gregoriocolor}{c} & \textcolor{gregoriocolor}{d} & \textcolor{gregoriocolor}{e} & \textcolor{gregoriocolor}{f} & \textcolor{gregoriocolor}{g} & \textcolor{gregoriocolor}{h} & \textcolor{gregoriocolor}{i} & \textcolor{gregoriocolor}{k} & \textcolor{gregoriocolor}{l} & \textcolor{gregoriocolor}{m} & \textcolor{gregoriocolor}{n} & \textcolor{gregoriocolor}{p} & \textcolor{gregoriocolor}{q} & \textcolor{gregoriocolor}{r} & \textcolor{gregoriocolor}{s} & \textcolor{gregoriocolor}{t} & \textcolor{gregoriocolor}{u} \\
 11 & 12 & 13 & 14 & 15 & 16 & 17 & 18 & 19 & 20 & 21 & 22 & 23 & 24 & 25 & 26 & 27 & 28 & 29 \\
\end{tabularx}
\vspace{0.5\baselineskip}
\noindent\begin{tabularx}{\linewidth}{*{12}{>{\centering\arraybackslash}X}}
 \textcolor{gregoriocolor}{A} & \textcolor{gregoriocolor}{B} & \textcolor{gregoriocolor}{C} & \textcolor{gregoriocolor}{D} & \textcolor{gregoriocolor}{E} & F & \textcolor{gregoriocolor}{F} & \textcolor{gregoriocolor}{G} & \textcolor{gregoriocolor}{H} & \textcolor{gregoriocolor}{M} & \textcolor{gregoriocolor}{N} & \textcolor{gregoriocolor}{P} \\
 1 & 2 & 3 & 4 & 5 & 6 & 5 & 6 & 7 & 8 & 9 & 10 \\
\end{tabularx}

\begin{paracol}{2}
\selectlanguage{latin}
\lettrine[lines=1]{S}{ancti} Nahum Prophétæ, in Bégabar quiescéntis.
\switchcolumn
\selectlanguage{english}
\lettrine[lines=1]{T}{he} prophet Nahum, who was buried in Bagabar.
\switchcolumn*
\selectlanguage{latin}
Romæ sanctórum Mártyrum 
 Diodóri Presbyteri, et Mariáni Diáconi, cum áliis plúribus, qui, sub 
 Numeriáno Príncipe, cum in Arenário natalítia Mártyrum ágerent, illic, 
 obstrúcta a persecutóribus jánua cryptæ ac díruta désuper mole, martyrii 
 glóriam meruérunt.
\switchcolumn
\selectlanguage{english}
At Rome, the holy martyrs Diodorus, 
 a priest, and Marian, a deacon, with many others, while they were observing 
 the birthdays of the martyrs in the catacombs. They were made 
 partakers in the glory of martyrdom when the persecutors, by order of 
 Emperor Numerian, walled up the door of the oratory and piled up a great 
 mass of stones against it.
\switchcolumn*
\selectlanguage{latin}
Item Romæ pássio 
 sanctórum Lúcii, Rogáti, Cassiáni et Cándidæ.
\switchcolumn
\selectlanguage{english}
Also in Rome, the martyrdom of the 
 Saints Lucius, Rogatus, Cassian, and Candida.
\switchcolumn*
\selectlanguage{latin}
Nárniæ sancti Próculi, 
 Epíscopi et Mártyris; qui, post multa egrégia ópera, a Rege Gothórum Tótila 
 jussus est decollári.
\switchcolumn
\selectlanguage{english}
At Narni, St. Proculus, bishop and 
 martyr, who, after performing many good works, was beheaded by order of 
 Totila, king of the Goths.
\switchcolumn*
\selectlanguage{latin}
In civitáte Casalénsi 
 sancti Evásii, Epíscopi et Mártyris.
\switchcolumn
\selectlanguage{english}
At Casale, St. Evasius, bishop and 
 martyr.
\switchcolumn*
\selectlanguage{latin}
Eódem die sancti Ansáni 
 Mártyris, qui, sub Diocletiáno Imperatóre, Romæ conféssus Christum et in cárcerem trusus, deínde Senas, in Túscia, perdúctus, ibídem cápitis 
 obtruncatióne cursum martyrii perfécit.
\switchcolumn
\selectlanguage{english}
The same day, St. Ansanus, martyr, 
 who confessed Christ at Rome, and was cast into prison in the time of 
 Emperor Diocletian. Afterwards he was taken to Siena in Tuscany, where 
 he ended the course of his martyrdom by beheading.
\switchcolumn*
\selectlanguage{latin}
Amériæ, in Umbria, 
 sancti Olympíadis, viri Consuláris, qui a beáta Firmína ad fidem est 
 convérsus, et sub Diocletiáno, in equúleo tortus, martyrium consummávit.
\switchcolumn
\selectlanguage{english}
At Amelia in Umbria, St. Olympias, 
 ex-consul, who was converted to the faith by blessed Firmina, was tortured 
 on the rack, and under Diocletian achieved martyrdom.
\switchcolumn*
\selectlanguage{latin}
Arbéle, in Pérside, 
 sancti Anániæ Mártyris.
\switchcolumn
\selectlanguage{english}
At Arbela in Persia, St. Ananias, 
 martyr.
\switchcolumn*
\selectlanguage{latin}
Medioláni sancti 
 Castritiáni Epíscopi, qui, in máxima Ecclésiæ perturbatióne, virtútum 
 méritis ac rerum pie religioséque gestárum laude enítuit.
\switchcolumn
\selectlanguage{english}
At Milan, St. Castritian, bishop, 
 who was eminent for virtues and the practice of pious and religious deeds 
 during the greatest troubles of the Church.
\switchcolumn*
\selectlanguage{latin}
Bríxiæ sancti Ursicíni 
 Epíscopi.
\switchcolumn
\selectlanguage{english}
At Brescia, St. Ursicinus, bishop.
\switchcolumn*
\selectlanguage{latin}
Noviómi, in Bélgio, 
 sancti Elígii Epíscopi, cujus vitam admirándam múltiplex signórum númerus 
 comméndat.
\switchcolumn
\selectlanguage{english}
At Noyon in Belgium, St. Eligius, 
 bishop, whose life is rendered illustrious by a considerable number of 
 miracles.
\switchcolumn*
\selectlanguage{latin}
Apud Virodúnum, in 
 Gállia, sancti Ageríci Epíscopi.
\switchcolumn
\selectlanguage{english}
At Verdun in France, St. Agericus, 
 bishop.
\switchcolumn*
\selectlanguage{latin}
Eódem die sanctæ 
 Natalítiæ, uxóris beáti Hadriáni Mártyris, quæ, sub Diocletiáno Imperatóre, 
 sanctis Martyribus, Nicomedíæ in cárcere deténtis, multo témpore ministrávit; 
 impletóque eórum certámine, Constantinópolim est profécta, et ibídem in pace 
 quiévit.
\switchcolumn
\selectlanguage{english}
The same day, St. Natalia, wife of 
 the blessed martyr Adrian, in the time of Emperor Diocletian. She long 
 served the holy martyrs imprisoned at Nicomedia, and when their trials were 
 over, went to Constantinople where she peacefully went to her rest in the 
 Lord.
\switchcolumn*
\selectlanguage{latin}
\end{paracol}

% \continueMonth{December}


% ---- martyrology/mart12/mart1202.htm
\needspace{10\baselineskip}
\begin{paracol}{2}
\selectlanguage{latin}
\begin{center}{\color{gregoriocolor} Quarto Nonas Decémbris. 
 Luna\dots\ }\end{center}
\switchcolumn
\selectlanguage{english}
\begin{center}{\color{gregoriocolor} The 
 2\textsuperscript{nd} Day of 
 December. The\dots\ Day of the Moon.}\end{center}
\end{paracol}

\noindent\begin{tabularx}{\linewidth}{*{19}{>{\centering\arraybackslash}X}}
 \textcolor{gregoriocolor}{a} & \textcolor{gregoriocolor}{b} & \textcolor{gregoriocolor}{c} & \textcolor{gregoriocolor}{d} & \textcolor{gregoriocolor}{e} & \textcolor{gregoriocolor}{f} & \textcolor{gregoriocolor}{g} & \textcolor{gregoriocolor}{h} & \textcolor{gregoriocolor}{i} & \textcolor{gregoriocolor}{k} & \textcolor{gregoriocolor}{l} & \textcolor{gregoriocolor}{m} & \textcolor{gregoriocolor}{n} & \textcolor{gregoriocolor}{p} & \textcolor{gregoriocolor}{q} & \textcolor{gregoriocolor}{r} & \textcolor{gregoriocolor}{s} & \textcolor{gregoriocolor}{t} & \textcolor{gregoriocolor}{u} \\
 12 & 13 & 14 & 15 & 16 & 17 & 18 & 19 & 20 & 21 & 22 & 23 & 24 & 25 & 26 & 27 & 28 & 29 & 1 \\
\end{tabularx}
\vspace{0.5\baselineskip}
\noindent\begin{tabularx}{\linewidth}{*{12}{>{\centering\arraybackslash}X}}
 \textcolor{gregoriocolor}{A} & \textcolor{gregoriocolor}{B} & \textcolor{gregoriocolor}{C} & \textcolor{gregoriocolor}{D} & \textcolor{gregoriocolor}{E} & F & \textcolor{gregoriocolor}{F} & \textcolor{gregoriocolor}{G} & \textcolor{gregoriocolor}{H} & \textcolor{gregoriocolor}{M} & \textcolor{gregoriocolor}{N} & \textcolor{gregoriocolor}{P} \\
 2 & 3 & 4 & 5 & 6 & 7 & 6 & 7 & 8 & 9 & 10 & 11 \\
\end{tabularx}

\begin{paracol}{2}
\selectlanguage{latin}
\lettrine[lines=2]{R}{omæ} pássio sanctæ 
 Bibiánæ, Vírginis et Mártyris, quæ, sub Juliáno Imperatóre sacrílego, ob 
 Christum támdiu plumbátis cæsa est, donec rédderet spíritum.
\switchcolumn
\selectlanguage{english}
\lettrine[lines=2]{A}{t} Rome, the martyrdom of the 
 saintly virgin Bibiana, under the sacrilegious Emperor Julian. For the 
 sake of our Lord she was scourged with leaded whips until she expired.
\switchcolumn*
\selectlanguage{latin}
Apud Forum Cornélii, in 
 Æmília, natális sancti Petri, Epíscopi Ravennátis, Confessóris et Ecclésiæ 
 Doctóris, cognoménto Chrysólogi, doctrína et sanctitáte célebris. 
 Ipsíus tamen festum prídie Nonas hujus mensis recólitur.
\switchcolumn
\selectlanguage{english}
At Imola, St. Peter Chrysologus, 
 bishop of Ravenna, confessor and doctor of the Church, celebrated for his 
 learning and sanctity. His feast is celebrated on the 4th of this 
 month.
\switchcolumn*
\selectlanguage{latin}
In Sanciáno, Sinárum ínsula, item natális sancti Francísci Xavérii, Sacerdótis e Societáte Jesu 
 et Confessóris, Indiárum Apóstoli, géntium conversióne, donis et miráculis 
 clari; qui plenus méritis et labóribus obdormívit in Dómino. Ipsum 
 beátum virum Pius Décimus, Póntifex Máximus, cæléstem sodalitáti et óperi 
 Propagándæ Fídei Protectórem elégit atque constítuit; Pius vero Papa 
 Undécimus peculiárem ómnibus Missiónibus Patrónum dedit et confirmávit. 
 Ejus autem festívitas, jussu Alexándri Papæ Séptimi, sequénti die celebrátur.
\switchcolumn
\selectlanguage{english}
In Sanchan, an island of China, the 
 birthday of St. Francis Xavier, priest of the Society of Jesus, confessor 
 and Apostle of the Indies. He was renowned for his conversion of the 
 heathen, his gifts and miracles, and he was filled with merits and good 
 works when he fell asleep in the Lord. Pope Pius X chose and appointed 
 him the heavenly protector of the Society for the Propagation of the Faith 
 and of the work for the same object. Pope Pius XI confirmed this and 
 appointed him the special patron of all the Foreign Missions. His 
 feast, by decree of Pope Alexander VII, is kept on the following day.
\switchcolumn*
\selectlanguage{latin}
Romæ sanctórum Mártyrum 
 Eusébii Presbyteri, Marcélli Diáconi, Hippólyti, Máximi, Adriæ, Paulínæ, 
 Neónis, Maríæ, Martánæ et Auréliæ; qui omnes in persecutióne Valeriáni, sub 
 Secundiáno Júdice, martyrium complevérunt.
\switchcolumn
\selectlanguage{english}
At Rome, the holy martyrs Eusebius, 
 a priest, Marcellus, a deacon, Hippolytus, Maximus, Adria, Paulina, Neon, 
 Mary, Martana, and Aurelia, who fulfilled their martyrdoms under the judge 
 Secundian in the persecution of Valerian.
\switchcolumn*
\selectlanguage{latin}
Item Romæ sancti 
 Pontiáni Mártyris, cum áliis quátuor.
\switchcolumn
\selectlanguage{english}
Also at Rome, St. Pontian, martyr, 
 with four others.
\switchcolumn*
\selectlanguage{latin}
In Africa natális 
 sanctórum Mártyrum Sevéri, Secúri, Januárii et Victoríni; qui ibídem 
 martyrio coronáti sunt.
\switchcolumn
\selectlanguage{english}
In Africa, the birthday of the holy 
 martyrs Severus, Securus, Januarius, and Victorinus, who were there crowned 
 with martyrdom.
\switchcolumn*
\selectlanguage{latin}
Aquiléjæ sancti 
 Chromátii, Epíscopi et Confessóris.
\switchcolumn
\selectlanguage{english}
At Aquileia, St. Chromatius, bishop 
 and confessor.
\switchcolumn*
\selectlanguage{latin}
Verónæ sancti Lupi, 
 Epíscopi et Confessóris.
\switchcolumn
\selectlanguage{english}
At Verona, St. Lupus, bishop and 
 confessor.
\switchcolumn*
\selectlanguage{latin}
Edéssæ, in Syria, 
 sancti Nonni Epíscopi, cujus précibus Pelágia pænitens ad Christum convérsa 
 est.
\switchcolumn
\selectlanguage{english}
At Edessa in Syria, St. Nonnus, 
 bishop, by whose prayers Pelagia the penitent was converted to Christ.
\switchcolumn*
\selectlanguage{latin}
Tróade, in Phrygia, 
 sancti Silváni Epíscopi, miráculis clari.
\switchcolumn
\selectlanguage{english}
At Troas in Phrygia, St. Silvanus, 
 bishop, renowned for miracles.
\switchcolumn*
\selectlanguage{latin}
Bríxiæ sancti Evásii 
 Epíscopi.
\switchcolumn
\selectlanguage{english}
At Brescia, St. Evasius, bishop.
\switchcolumn*
\selectlanguage{latin}
\end{paracol}


% ---- martyrology/mart12/mart1203.htm
\needspace{10\baselineskip}
\begin{paracol}{2}
\selectlanguage{latin}
\begin{center}{\color{gregoriocolor} Tértio Nonas Decémbris. 
 Luna\dots\ }\end{center}
\switchcolumn
\selectlanguage{english}
\begin{center}{\color{gregoriocolor} The 
 3\textsuperscript{rd} Day of 
 December. The\dots\ Day of the Moon.}\end{center}
\end{paracol}

\noindent\begin{tabularx}{\linewidth}{*{19}{>{\centering\arraybackslash}X}}
 \textcolor{gregoriocolor}{a} & \textcolor{gregoriocolor}{b} & \textcolor{gregoriocolor}{c} & \textcolor{gregoriocolor}{d} & \textcolor{gregoriocolor}{e} & \textcolor{gregoriocolor}{f} & \textcolor{gregoriocolor}{g} & \textcolor{gregoriocolor}{h} & \textcolor{gregoriocolor}{i} & \textcolor{gregoriocolor}{k} & \textcolor{gregoriocolor}{l} & \textcolor{gregoriocolor}{m} & \textcolor{gregoriocolor}{n} & \textcolor{gregoriocolor}{p} & \textcolor{gregoriocolor}{q} & \textcolor{gregoriocolor}{r} & \textcolor{gregoriocolor}{s} & \textcolor{gregoriocolor}{t} & \textcolor{gregoriocolor}{u} \\
 13 & 14 & 15 & 16 & 17 & 18 & 19 & 20 & 21 & 22 & 23 & 24 & 25 & 26 & 27 & 28 & 29 & 1 & 2 \\
\end{tabularx}
\vspace{0.5\baselineskip}
\noindent\begin{tabularx}{\linewidth}{*{12}{>{\centering\arraybackslash}X}}
 \textcolor{gregoriocolor}{A} & \textcolor{gregoriocolor}{B} & \textcolor{gregoriocolor}{C} & \textcolor{gregoriocolor}{D} & \textcolor{gregoriocolor}{E} & F & \textcolor{gregoriocolor}{F} & \textcolor{gregoriocolor}{G} & \textcolor{gregoriocolor}{H} & \textcolor{gregoriocolor}{M} & \textcolor{gregoriocolor}{N} & \textcolor{gregoriocolor}{P} \\
 3 & 4 & 5 & 6 & 7 & 8 & 7 & 8 & 9 & 10 & 11 & 12 \\
\end{tabularx}

\begin{paracol}{2}
\selectlanguage{latin}
\lettrine[lines=2]{S}{ancti} Francísci 
 Xavérii, Sacerdótis e Societáte Jesu et Confessóris, Indiárum Apóstoli, 
 sodalitátis et óperis Propagándæ Fídei atque Missiónum ómnium Patróni 
 cæléstis; qui prídie hujus diéi quiévit in pace.
\switchcolumn
\selectlanguage{english}
\lettrine[lines=2]{S}{t.} Francis Xavier, priest of the 
 Society of Jesus, confessor, Apostle of the Indies, and heavenly patron of 
 the Society for the Propagation of the Faith, and also of all the Missions, 
 who died on the day previous.
\switchcolumn*
\selectlanguage{latin}
In Judæa sancti Sophoníæ Prophétæ.
\switchcolumn
\selectlanguage{english}
In Judea, the holy prophet Zephaniah.
\switchcolumn*
\selectlanguage{latin}
Romæ sanctórum Mártyrum 
 Cláudii Tribúni, et uxóris Hiláriæ, ac filiórum Jásonis et Mauri, cum 
 septuagínta milítibus. Ex eis Cláudium jussit Numeriánus Imperátor, 
 ingénti saxo alligátum, in flumen præcípitem dari; mílites
 vero et ipsíus 
 Cláudii fílios capitáli senténtia puníri. Beáta autem Hilária, cum 
 filiórum córpora sepelísset, paulo post, orans ad eórum sepúlcrum, tenta est 
 a Pagánis, et, in cárcerem trusa, migrávit ad Dóminum.
\switchcolumn
\selectlanguage{english}
At Rome, the holy martyrs Claudius, 
 a tribune, and Hilaria, his wife, with Jason and Maur, their sons, and 
 seventy soldiers. By the command of Emperor Numerian, Claudius was 
 fastened to a large stone and thrown into the river, the soldiers and the 
 sons of Claudius were condemned to capital punishment. But blessed 
 Hilaria, after having buried the bodies of her sons, and while praying at 
 their tomb, was arrested by the pagans, and shortly after departed for 
 heaven.
\switchcolumn*
\selectlanguage{latin}
Tingi, in Mauritánia, pássio sancti Cassiáni Mártyris, qui, cum exceptóris diu gessísset offícium, 
 tandem, admirátus intrépida beáti Marcélli Centuriónis respónsa et immóbilem 
 in Christi fide constántiam, atque cælitus inspirátus, exsecrábile duxit 
 Christianórum neci deservíre; ideóque, cum renuntiásset eídem offício, et 
 ipse, sub Christiána professióne, cápite abscíssus, triúmphum méruit 
 obtinére martyrii.
\switchcolumn
\selectlanguage{english}
At Tangier in Morocco, St. Cassian, 
 martyr. After having been a recorder for a long time, at length, by an 
 inspiration from heaven, he deemed it a hateful thing to contribute to the 
 massacre of the Christians, and therefore abandoned his office, and making a 
 profession of Christianity, he deserved to obtain the triumph of martyrdom.
\switchcolumn*
\selectlanguage{latin}
Item, in Africa 
 sanctórum Mártyrum Cláudii, Crispíni, Magínæ, Joánnis et Stéphani.
\switchcolumn
\selectlanguage{english}
Also in Africa, the holy martyrs 
 Claudius, Crispin, Magina, John, and Stephen.
\switchcolumn*
\selectlanguage{latin}
In Pannónia sancti 
 Agrícolæ Mártyris.
\switchcolumn
\selectlanguage{english}
In Hungary, St. Agricola, martyr.
\switchcolumn*
\selectlanguage{latin}
Nicomedíæ pássio 
 sanctórum Ambici, Victóris et Júlii.
\switchcolumn
\selectlanguage{english}
At Nicomedia, the martyrdom of the 
 Saints Ambicus, Victor, and Julius.
\switchcolumn*
\selectlanguage{latin}
Medioláni sancti 
 Miroclétis, Epíscopi et Confessóris; cujus aliquándo sanctus Ambrósius 
 méminit.
\switchcolumn
\selectlanguage{english}
At Milan, St. Mirocles, bishop and 
 confessor, sometimes mentioned by St. Ambrose.
\switchcolumn*
\selectlanguage{latin}
Dorcéstriæ, in Anglia, 
 sancti Biríni, qui fuit primus ejúsdem civitátis Epíscopus.
\switchcolumn
\selectlanguage{english}
At Dorchester in England, St. 
 Birinus, who was the first bishop of that city.
\switchcolumn*
\selectlanguage{latin}
Cúriæ, in Germánia, sancti Lúcii, Britannórum Regis, qui primus ex iis Régibus fidem Christi 
 suscépit, témpore sancti Eleuthérii Papæ.
\switchcolumn
\selectlanguage{english}
At Chur in Germany, St. Lucius, king 
 of the Britons, who in the time of Pope Eleutherius, was the first of their 
 kings to receive the faith of Christ.
\switchcolumn*
\selectlanguage{latin}
Senis, in Túscia, 
 sancti Galgáni Eremítæ.
\switchcolumn
\selectlanguage{english}
At Siena in Tuscany, St. Galganus, 
 hermit.
\switchcolumn*
\selectlanguage{latin}
\end{paracol}


% ---- martyrology/mart12/mart1204.htm
\needspace{10\baselineskip}
\begin{paracol}{2}
\selectlanguage{latin}
\begin{center}{\color{gregoriocolor} Prídie Nonas Decémbris. 
 Luna\dots\ }\end{center}
\switchcolumn
\selectlanguage{english}
\begin{center}{\color{gregoriocolor} The 
 4\textsuperscript{th} Day of 
 December. The\dots\ Day of the Moon.}\end{center}
\end{paracol}

\noindent\begin{tabularx}{\linewidth}{*{19}{>{\centering\arraybackslash}X}}
 \textcolor{gregoriocolor}{a} & \textcolor{gregoriocolor}{b} & \textcolor{gregoriocolor}{c} & \textcolor{gregoriocolor}{d} & \textcolor{gregoriocolor}{e} & \textcolor{gregoriocolor}{f} & \textcolor{gregoriocolor}{g} & \textcolor{gregoriocolor}{h} & \textcolor{gregoriocolor}{i} & \textcolor{gregoriocolor}{k} & \textcolor{gregoriocolor}{l} & \textcolor{gregoriocolor}{m} & \textcolor{gregoriocolor}{n} & \textcolor{gregoriocolor}{p} & \textcolor{gregoriocolor}{q} & \textcolor{gregoriocolor}{r} & \textcolor{gregoriocolor}{s} & \textcolor{gregoriocolor}{t} & \textcolor{gregoriocolor}{u} \\
 14 & 15 & 16 & 17 & 18 & 19 & 20 & 21 & 22 & 23 & 24 & 25 & 26 & 27 & 28 & 29 & 1 & 2 & 3 \\
\end{tabularx}
\vspace{0.5\baselineskip}
\noindent\begin{tabularx}{\linewidth}{*{12}{>{\centering\arraybackslash}X}}
 \textcolor{gregoriocolor}{A} & \textcolor{gregoriocolor}{B} & \textcolor{gregoriocolor}{C} & \textcolor{gregoriocolor}{D} & \textcolor{gregoriocolor}{E} & F & \textcolor{gregoriocolor}{F} & \textcolor{gregoriocolor}{G} & \textcolor{gregoriocolor}{H} & \textcolor{gregoriocolor}{M} & \textcolor{gregoriocolor}{N} & \textcolor{gregoriocolor}{P} \\
 4 & 5 & 6 & 7 & 8 & 9 & 8 & 9 & 10 & 11 & 12 & 13 \\
\end{tabularx}

\begin{paracol}{2}
\selectlanguage{latin}
\lettrine[lines=2]{S}{ancti} Petri Chrysólogi, 
 Epíscopi Ravennátis, Confessóris et Ecclésiæ Doctóris, cujus memória quarto 
 Nonas hujus mensis recensétur.
\switchcolumn
\selectlanguage{english}
\lettrine[lines=2]{S}{t.} Peter Chrysologus, bishop of 
 Ravenna, confessor, and doctor of the Church, whose birthday is kept on the 
 2nd of December.
\switchcolumn*
\selectlanguage{latin}
Nicomedíæ pássio sanctæ 
 Bárbaræ, Vírginis et Mártyris; quæ, in persecutióne Maximíni, post diram cárceris maceratiónem, lampadárum adustiónem, mamillárum præcisiónem atque 
 ália torménta, gládio martyrium consummávit.
\switchcolumn
\selectlanguage{english}
At Nicomedia, the passion of St. 
 Barbara, virgin and martyr, in the persecution of Maximinus. After a 
 series of sufferings, a long imprisonment, the burning with torches, and the 
 cutting away of her breasts, her martyrdom was fulfilled by the sword.
\switchcolumn*
\selectlanguage{latin}
Constantinópoli 
 sanctórum Theóphanis et Sociórum.
\switchcolumn
\selectlanguage{english}
At Constantinople, St. Theophanes 
 and his companions.
\switchcolumn*
\selectlanguage{latin}
In Ponto beáti Melétii, 
 Epíscopi et Confessóris; qui, cum esset ob eruditiónis prærogatívam 
 præcípuus, ob virtútem tamen ánimi et vitæ sinceritátem longe magnificéntior 
 éxstitit.
\switchcolumn
\selectlanguage{english}
In Pontus, blessed Meletius, bishop 
 and confessor, who joined to an eminent gift of knowledge the more 
 distinguished glory of fortitude and integrity of life.
\switchcolumn*
\selectlanguage{latin}
Bonóniæ sancti Felícis Epíscopi; qui, ántea Mediolanénsis Ecclésiæ, sub sancto
  Ambrósio, Diáconus fúerat.
\switchcolumn
\selectlanguage{english}
At Bologna, St. Felix, bishop was 
 one time deacon of the Milanese Church under St. Ambrose.
\switchcolumn*
\selectlanguage{latin}
In Anglia, sancti 
 Osmúndi, Epíscopi et Confessóris.
\switchcolumn
\selectlanguage{english}
In England, St. Osmund, bishop and 
 confessor.
\switchcolumn*
\selectlanguage{latin}
Colóniæ Agrippínæ 
 sancti Annónis Epíscopi.
\switchcolumn
\selectlanguage{english}
At Cologne, St. Anno, bishop.
\switchcolumn*
\selectlanguage{latin}
In Mesopotámia sancti 
 Marúthæ Epíscopi, qui Dei Ecclésias, ob persecutiónem Isdegérdis Regis 
 collápsas, in Pérside reparávit, multísque miráculis clarus, apud hostes 
 étiam méruit honorári.
\switchcolumn
\selectlanguage{english}
In Mesopotamia, St. Maruthas, 
 bishop, who restored the churches of God that had been ruined in Persia by 
 the persecution of King Isdegerd. Being renowned for many miracles, he 
 merited to be honoured even by his enemies.
\switchcolumn*
\selectlanguage{latin}
Parmæ sancti Bernárdi, 
 Cardinális et ejúsdem civitátis Epíscopi, ex Ordine Vallis Umbrósæ.
\switchcolumn
\selectlanguage{english}
At Parma, St. Bernard, cardinal and 
 bishop of that city, of the Congregation of Vallombrosa of the Order of St. 
 Benedict.
\switchcolumn*
\selectlanguage{latin}
\end{paracol}


% ---- martyrology/mart12/mart1205.htm
\needspace{10\baselineskip}
\begin{paracol}{2}
\selectlanguage{latin}
\begin{center}{\color{gregoriocolor} Nonis Decémbris. 
 Luna\dots\ }\end{center}
\switchcolumn
\selectlanguage{english}
\begin{center}{\color{gregoriocolor} The 
 5\textsuperscript{th} Day of 
 December. The\dots\ Day of the Moon.}\end{center}
\end{paracol}

\noindent\begin{tabularx}{\linewidth}{*{19}{>{\centering\arraybackslash}X}}
 \textcolor{gregoriocolor}{a} & \textcolor{gregoriocolor}{b} & \textcolor{gregoriocolor}{c} & \textcolor{gregoriocolor}{d} & \textcolor{gregoriocolor}{e} & \textcolor{gregoriocolor}{f} & \textcolor{gregoriocolor}{g} & \textcolor{gregoriocolor}{h} & \textcolor{gregoriocolor}{i} & \textcolor{gregoriocolor}{k} & \textcolor{gregoriocolor}{l} & \textcolor{gregoriocolor}{m} & \textcolor{gregoriocolor}{n} & \textcolor{gregoriocolor}{p} & \textcolor{gregoriocolor}{q} & \textcolor{gregoriocolor}{r} & \textcolor{gregoriocolor}{s} & \textcolor{gregoriocolor}{t} & \textcolor{gregoriocolor}{u} \\
 15 & 16 & 17 & 18 & 19 & 20 & 21 & 22 & 23 & 24 & 25 & 26 & 27 & 28 & 29 & 1 & 2 & 3 & 4 \\
\end{tabularx}
\vspace{0.5\baselineskip}
\noindent\begin{tabularx}{\linewidth}{*{12}{>{\centering\arraybackslash}X}}
 \textcolor{gregoriocolor}{A} & \textcolor{gregoriocolor}{B} & \textcolor{gregoriocolor}{C} & \textcolor{gregoriocolor}{D} & \textcolor{gregoriocolor}{E} & F & \textcolor{gregoriocolor}{F} & \textcolor{gregoriocolor}{G} & \textcolor{gregoriocolor}{H} & \textcolor{gregoriocolor}{M} & \textcolor{gregoriocolor}{N} & \textcolor{gregoriocolor}{P} \\
 5 & 6 & 7 & 8 & 9 & 10 & 9 & 10 & 11 & 12 & 13 & 14 \\
\end{tabularx}

\begin{paracol}{2}
\selectlanguage{latin}
\lettrine[lines=2]{I}{n} Judæa sancti Sabbæ 
 Abbátis, in óppido Cappadóciæ Mútala orti, qui miro sanctitátis exémplo 
 refúlsit, et pro fide cathólica, advérsus impugnántes sanctam Synodum 
 Chalcedonénsem, strénue laborávit, ac tandem in ea diœcésis Hierosolymitánæ 
 laura, quæ ipsíus sancti Sabbæ nómine póstmodum est insigníta, requiévit in 
 pace.
\switchcolumn
\selectlanguage{english}
\lettrine[lines=2]{I}{n} Judea, St. Sabbas, abbot, who was 
 born in the town of Mutala in Cappadocia. He gave a wondrous example 
 of holiness and laboured most zealously for the Catholic faith against those 
 who attacked the holy Council of Chalcedon. He rested in peace in the 
 monastery later named for him in the diocese of Jerusalem.
\switchcolumn*
\selectlanguage{latin}
Níciæ, apud Varum 
 flúvium, sancti Bassi Epíscopi, qui, in persecutióne Décii et Valeriáni, a 
 Perénnio Præside, ob Christi fidem, equúleo tortus, láminis candéntibus 
 ustus, fústibus et scorpiónibus cæsus, in ignem missus, et, cum inde 
 evasísset illæsus, duóbus clavis confíxus, illústre martyrium consummávit.
\switchcolumn
\selectlanguage{english}
At Nice, near the river Var, St. 
 Bassus, bishop. In the persecution of Decius and Valerian, he was 
 tortured by the governor Perennius for the faith of Christ, burned with hot 
 plates of metal, beaten with rods and whips garnished with pieces of iron, 
 and thrown into the fire. When he came out of it unhurt, he was 
 pierced with two spikes, and thus completed an illustrious martyrdom.
\switchcolumn*
\selectlanguage{latin}
Papíæ sancti Dalmátii, 
 Epíscopi et Mártyris; qui in persecutióne Maximiáni passus est.
\switchcolumn
\selectlanguage{english}
At Pavia, St. Dalmatius, bishop and 
 martyr, who suffered in the persecution of Maximian.
\switchcolumn*
\selectlanguage{latin}
Corfínii, in Pelígnis, 
 sancti Pelíni, Epíscopi Brundusíni, qui, cum ob ejus oratiónem, sub Juliáno Apóstata, templum Martis corruísset, a templórum Pontifícibus diríssime 
 cæsus est, atque, octogínta et quinque vulnéribus confóssus, martyrii 
 corónam proméruit.
\switchcolumn
\selectlanguage{english}
At Corfinio in Peligno, St. Pelinus, 
 bishop of Brindisi, at the time of Julian the Apostate. When the 
 temple of Mars fell to the ground at his prayer, he was severely scourged by 
 the priests of the temple, and being pierced with eighty-five wounds, he 
 merited the crown of martyrdom.
\switchcolumn*
\selectlanguage{latin}
Item sancti Anastásii 
 Mártyris, qui, præ ardóre martyrii, sponte se persecutóribus óbtulit.
\switchcolumn
\selectlanguage{english}
Also, St. Anastasius, martyr, who in 
 his ardent desire for martyrdom gave himself up voluntarily to the 
 persecutors.
\switchcolumn*
\selectlanguage{latin}
Thagúræ, in Africa, 
 sanctórum Mártyrum Júlii, Potámiæ, Crispíni, Felícis, Grati et aliórum 
 septem.
\switchcolumn
\selectlanguage{english}
At Thagura in Africa, the holy 
 martyrs Julius, Potamias, Crispin, Felix, Gratus, and seven others.
\switchcolumn*
\selectlanguage{latin}
Thebéste, in Numídia, sanctæ Crispínæ, nobilíssimæ féminæ, quæ, tempóribus Diocletiáni et 
 Maximiáni, cum sacrificáre nollet, jussu Anolíni Procónsulis decolláta est; 
 quam sanctus Augustínus sæpe láudibus célebrat.
\switchcolumn
\selectlanguage{english}
At Thebaste in Africa, St. Crispina, 
 a woman of the highest nobility who refused to sacrifice to idols during the 
 reign of Diocletian and Maximian, and was beheaded by order of the proconsul 
 Anolinus. Her praises are often celebrated by St. Augustine.
\switchcolumn*
\selectlanguage{latin}
Tréviris sancti Nicétii 
 Epíscopi, miræ sanctitátis viri.
\switchcolumn
\selectlanguage{english}
At Treves, St. Nicetius, bishop, a 
 man of great sanctity.
\switchcolumn*
\selectlanguage{latin}
Polyboti, in Asia, 
 sancti Joánnis Epíscopi, cognoménto Thaumatúrgi.
\switchcolumn
\selectlanguage{english}
At Polybotum in Asia, St. John, 
 bishop, surnamed the Wonderworker.
\switchcolumn*
\selectlanguage{latin}
\end{paracol}


% ---- martyrology/mart12/mart1206.htm
\needspace{10\baselineskip}
\begin{paracol}{2}
\selectlanguage{latin}
\begin{center}{\color{gregoriocolor} Octávo Idus Decémbris. 
 Luna\dots\ }\end{center}
\switchcolumn
\selectlanguage{english}
\begin{center}{\color{gregoriocolor} The 
 6\textsuperscript{th} Day of 
 December. The\dots\ Day of the Moon.}\end{center}
\end{paracol}

\noindent\begin{tabularx}{\linewidth}{*{19}{>{\centering\arraybackslash}X}}
 \textcolor{gregoriocolor}{a} & \textcolor{gregoriocolor}{b} & \textcolor{gregoriocolor}{c} & \textcolor{gregoriocolor}{d} & \textcolor{gregoriocolor}{e} & \textcolor{gregoriocolor}{f} & \textcolor{gregoriocolor}{g} & \textcolor{gregoriocolor}{h} & \textcolor{gregoriocolor}{i} & \textcolor{gregoriocolor}{k} & \textcolor{gregoriocolor}{l} & \textcolor{gregoriocolor}{m} & \textcolor{gregoriocolor}{n} & \textcolor{gregoriocolor}{p} & \textcolor{gregoriocolor}{q} & \textcolor{gregoriocolor}{r} & \textcolor{gregoriocolor}{s} & \textcolor{gregoriocolor}{t} & \textcolor{gregoriocolor}{u} \\
 16 & 17 & 18 & 19 & 20 & 21 & 22 & 23 & 24 & 25 & 26 & 27 & 28 & 29 & 1 & 2 & 3 & 4 & 5 \\
\end{tabularx}
\vspace{0.5\baselineskip}
\noindent\begin{tabularx}{\linewidth}{*{12}{>{\centering\arraybackslash}X}}
 \textcolor{gregoriocolor}{A} & \textcolor{gregoriocolor}{B} & \textcolor{gregoriocolor}{C} & \textcolor{gregoriocolor}{D} & \textcolor{gregoriocolor}{E} & F & \textcolor{gregoriocolor}{F} & \textcolor{gregoriocolor}{G} & \textcolor{gregoriocolor}{H} & \textcolor{gregoriocolor}{M} & \textcolor{gregoriocolor}{N} & \textcolor{gregoriocolor}{P} \\
 6 & 7 & 8 & 9 & 10 & 11 & 10 & 11 & 12 & 13 & 14 & 15 \\
\end{tabularx}

\begin{paracol}{2}
\selectlanguage{latin}
\lettrine[lines=2]{M}{yræ,} quæ est 
 metrópolis Lyciæ, natális sancti Nicolái, Epíscopi et Confessóris, de quo, 
 inter plura miraculórum insígnia, illud memorábile fertur, quod Imperatórem 
 Constantínum ab intéritu quorúmdam se invocántium, longe constitútus, ad 
 misericórdiam per visum mónitis defléxit et minis.
\switchcolumn
\selectlanguage{english}
\lettrine[lines=2]{A}{t} Myra, which is the metropolis of 
 Lycia, the birthday of St. Nicholas, bishop and confessor, of whom it is 
 related, among other miracles, that, while at a great distance from Emperor 
 Constantine, he appeared to him in a vision and moved him to mercy so as to 
 deter him from putting to death some persons who had implored his 
 assistance.
\switchcolumn*
\selectlanguage{latin}
Eódem die sancti 
 Polychrónii Presbyteri, qui, témpore Constántii Imperatóris, cum ad altáre 
 Missas ágeret, invásus est ab Ariánis et jugulátus.
\switchcolumn
\selectlanguage{english}
On the same day, St. Polychronius, 
 priest, who was surprised while offering Mass at the altar and slain by the 
 Arians, in the reign of Emperor Constantius.
\switchcolumn*
\selectlanguage{latin}
In Africa sancti 
 Majórici, fílii sanctæ Dionysiæ, qui, cum esset adolescéntulus ac torménta 
 pavésceret, matris obtútibus verbísque corroborátus est, et, céteris fórtior 
 factus, in torméntis ánimam réddidit; quem amplexáta mater domi sepelívit, 
 et ad ejus sepúlcrum assídue oráre consuévit.
\switchcolumn
\selectlanguage{english}
In Africa, St. Majoricus, son of St. 
 Dionysia, who, being quite young and dreading the torments, was strengthened 
 by the looks and words of his mother, and becoming stronger than the rest, 
 expired in torments. His mother took him in her arms, and having 
 buried him in her own home, was wont to pray diligently at his tomb.
\switchcolumn*
\selectlanguage{latin}
Ibídem sanctárum 
 mulíerum Dionysiæ, quæ sancti Majórici Mártyris éxstitit mater, Datívæ ac 
 Leóntiæ; itémque religiósi viri, nómine Tértii, Æmiliáni médici, et 
 Bonifátii, cum áliis tribus. Hi omnes, in persecutióne Wandálica, sub 
 Ariáno Rege Hunneríco, gravíssimis et innúmeris supplíciis pro cathólicæ 
 fídei defensióne cruciáti, sanctórum Christi Confessórum número sociári 
 meruérunt.
\switchcolumn
\selectlanguage{english}
In the same place, the holy women 
 Dionysia, who was the mother of St. Majoricus the martyr, Dativa, and Leontia; 
 also a pious man named Tertius, Emilian a physician, Boniface, and three 
 others. In the persecution of the Vandals, under the Arian king 
 Hunneric, they were subjected to numberless most painful tortures for the 
 Catholic faith, and thus merited to rank among the confessors of Christ.
\switchcolumn*
\selectlanguage{latin}
Romæ sanctæ Aséllæ 
 Vírginis, quæ (ut beátus Hierónymus scribit), ex útero matris benedícta, 
 vitam in jejúniis et oratiónibus usque ad senéctam prodúxit.
\switchcolumn
\selectlanguage{english}
At Rome, St. Asella, virgin, who 
 according to the words of St. Jerome, being blessed from her mother's womb, 
 lived to old age in fasting and prayer.
\switchcolumn*
\selectlanguage{latin}
Granátæ, in Hispánia, 
 pássio beáti Petri Paschásii, Epíscopi Giennénsis et Mártyris, ex Ordine 
 beátæ Maríæ de Mercéde redemptiónis captivórum.
\switchcolumn
\selectlanguage{english}
At Granada in Spain, the passion of 
 blessed Peter Paschasius, bishop of Jaen and martyr, a member of the Order 
 of our Lady of Ransom for the Redemption of Captives.
\switchcolumn*
\selectlanguage{latin}
\end{paracol}


% ---- martyrology/mart12/mart1207.htm
\needspace{10\baselineskip}
\begin{paracol}{2}
\selectlanguage{latin}
\begin{center}{\color{gregoriocolor} Séptimo Idus Decémbris. 
 Luna\dots\ }\end{center}
\switchcolumn
\selectlanguage{english}
\begin{center}{\color{gregoriocolor} The 
 7\textsuperscript{th} Day of 
 December. The\dots\ Day of the Moon.}\end{center}
\end{paracol}

\noindent\begin{tabularx}{\linewidth}{*{19}{>{\centering\arraybackslash}X}}
 \textcolor{gregoriocolor}{a} & \textcolor{gregoriocolor}{b} & \textcolor{gregoriocolor}{c} & \textcolor{gregoriocolor}{d} & \textcolor{gregoriocolor}{e} & \textcolor{gregoriocolor}{f} & \textcolor{gregoriocolor}{g} & \textcolor{gregoriocolor}{h} & \textcolor{gregoriocolor}{i} & \textcolor{gregoriocolor}{k} & \textcolor{gregoriocolor}{l} & \textcolor{gregoriocolor}{m} & \textcolor{gregoriocolor}{n} & \textcolor{gregoriocolor}{p} & \textcolor{gregoriocolor}{q} & \textcolor{gregoriocolor}{r} & \textcolor{gregoriocolor}{s} & \textcolor{gregoriocolor}{t} & \textcolor{gregoriocolor}{u} \\
 17 & 18 & 19 & 20 & 21 & 22 & 23 & 24 & 25 & 26 & 27 & 28 & 29 & 1 & 2 & 3 & 4 & 5 & 6 \\
\end{tabularx}
\vspace{0.5\baselineskip}
\noindent\begin{tabularx}{\linewidth}{*{12}{>{\centering\arraybackslash}X}}
 \textcolor{gregoriocolor}{A} & \textcolor{gregoriocolor}{B} & \textcolor{gregoriocolor}{C} & \textcolor{gregoriocolor}{D} & \textcolor{gregoriocolor}{E} & F & \textcolor{gregoriocolor}{F} & \textcolor{gregoriocolor}{G} & \textcolor{gregoriocolor}{H} & \textcolor{gregoriocolor}{M} & \textcolor{gregoriocolor}{N} & \textcolor{gregoriocolor}{P} \\
 7 & 8 & 9 & 10 & 11 & 12 & 11 & 12 & 13 & 14 & 15 & 16 \\
\end{tabularx}

\begin{paracol}{2}
\selectlanguage{latin}
\lettrine[lines=2]{V}{igília} Conceptiónis 
 Immaculátæ beátæ Maríæ Vírginis.
\switchcolumn
\selectlanguage{english}
\lettrine[lines=2]{T}{he} Vigil of the Immaculate 
 Conception of the Blessed Virgin Mary.
\switchcolumn*
\selectlanguage{latin}
Sancti Ambrósii 
 Epíscopi, Confessóris et Ecclésiæ Doctóris, qui prídie Nonas Aprílis 
 obdormívit in Dómino, sed hac die potíssimum cólitur, qua Mediolanénsem 
 Ecclésiam gubernándam suscépit.
\switchcolumn
\selectlanguage{english}
St. Ambrose, bishop and doctor of 
 the Church, who fell asleep in the Lord on the 4th of April; his feast is 
 kept on this day, the day on which he assumed the government of the Church 
 of Milan.
\switchcolumn*
\selectlanguage{latin}
Romæ beáti Eutychiáni 
 Papæ, qui per divérsa loca trecéntos quadragínta duos Mártyres manu sua 
 sepelívit; quibus et ipse deínde sociátus, sub Numeriáno Imperatóre, 
 martyrio coronátus est, et in cœmetério Callísti sepúltus.
\switchcolumn
\selectlanguage{english}
At Rome, blessed Eutychian, pope, 
 who with his own hand buried three hundred and forty-two martyrs in various 
 places. He himself was joined with them, crowned with martyrdom under 
 Emperor Numerian, and was buried in the cemetery of Callistus.
\switchcolumn*
\selectlanguage{latin}
Alexandríæ natális 
 beáti Agathónis militáris, qui, in persecutióne Décii, cum prohibéret 
 quosdam voléntes illúdere cadavéribus Mártyrum, clamor repénte totíus vulgi 
 advérsus eum extóllitur; oblátus autem Júdici, et in Christi confessióne 
 persístens, cápite pro pietáte damnátus est.
\switchcolumn
\selectlanguage{english}
At Alexandria, the birthday of 
 blessed Agatho, soldier. In the persecution of Decius, because he 
 prevented some people from mocking the bodies of the martyrs, a sudden 
 clamour was raised against him by the crowd. Being brought before the 
 judge, and persisting in his confession of Christ, he was sentenced to death 
 for his reverence.
\switchcolumn*
\selectlanguage{latin}
Antiochíæ sanctórum 
 Mártyrum Polycárpi et Theodóri.
\switchcolumn
\selectlanguage{english}
At Antioch, the holy martyrs 
 Polycarp and Theodore.
\switchcolumn*
\selectlanguage{latin}
Tubúrbi, in Africa, 
 sancti Servi Mártyris, qui, in persecutióne Wandálica, sub Ariáno Rege 
 Hunneríco, fústibus diutíssime cæsus, tróchleis frequénter in sublíme elevátus atque ictu céleri super sílices póndere córporis dimíssus, et 
 lapídibus acutíssimis perfricátus, martyrii palmam adéptus est.
\switchcolumn
\selectlanguage{english}
At Tuburbum in Africa, during the 
 persecution of the Vandals, under the Arian king Hunneric, St. Servus, 
 martyr, who, being for a very long time beaten with rods, lifted up on high 
 with pulleys, and suddenly dropped on flint-stones with his whole weight, 
 and rubbed over with sharp stones, obtained the palm of martyrdom.
\switchcolumn*
\selectlanguage{latin}
Theáni, in Campánia, sancti Urbáni, Epíscopi et Confessóris.
\switchcolumn
\selectlanguage{english}
At Teano in Campania, St. Urban, 
 bishop and confessor.
\switchcolumn*
\selectlanguage{latin}
Apud Sántonas, in 
 Gállia, sancti Martíni Abbátis, ad cujus túmulum crebérrima divínitus fiunt 
 mirácula.
\switchcolumn
\selectlanguage{english}
At Saintes in France, St. Martin, 
 abbot, at whose tomb frequent miracles have been worked through the power of 
 God.
\switchcolumn*
\selectlanguage{latin}
Eboríaci, in território 
 Meldénsi, commemorátio sanctæ Faræ, étiam Burgundofáræ nómine appellátæ, 
 Abbatíssæ et Vírginis, cujus dies natális tértio Nonas Aprilis recensétur.
\switchcolumn
\selectlanguage{english}
At Faremoutiers, in the diocese of 
 Meaux, the commemoration of St. Fara, who is also called Burgundofara, 
 abbess and virgin. Her birthday is on the 3rd of April.
\switchcolumn*
\selectlanguage{latin}
\end{paracol}


% ---- martyrology/mart12/mart1208.htm
\needspace{10\baselineskip}
\begin{paracol}{2}
\selectlanguage{latin}
\begin{center}{\color{gregoriocolor} Sexto Idus Decémbris. 
 Luna\dots\ }\end{center}
\switchcolumn
\selectlanguage{english}
\begin{center}{\color{gregoriocolor} The 
 8\textsuperscript{th} Day of 
 December. The\dots\ Day of the Moon.}\end{center}
\end{paracol}

\noindent\begin{tabularx}{\linewidth}{*{19}{>{\centering\arraybackslash}X}}
 \textcolor{gregoriocolor}{a} & \textcolor{gregoriocolor}{b} & \textcolor{gregoriocolor}{c} & \textcolor{gregoriocolor}{d} & \textcolor{gregoriocolor}{e} & \textcolor{gregoriocolor}{f} & \textcolor{gregoriocolor}{g} & \textcolor{gregoriocolor}{h} & \textcolor{gregoriocolor}{i} & \textcolor{gregoriocolor}{k} & \textcolor{gregoriocolor}{l} & \textcolor{gregoriocolor}{m} & \textcolor{gregoriocolor}{n} & \textcolor{gregoriocolor}{p} & \textcolor{gregoriocolor}{q} & \textcolor{gregoriocolor}{r} & \textcolor{gregoriocolor}{s} & \textcolor{gregoriocolor}{t} & \textcolor{gregoriocolor}{u} \\
 18 & 19 & 20 & 21 & 22 & 23 & 24 & 25 & 26 & 27 & 28 & 29 & 1 & 2 & 3 & 4 & 5 & 6 & 7 \\
\end{tabularx}
\vspace{0.5\baselineskip}
\noindent\begin{tabularx}{\linewidth}{*{12}{>{\centering\arraybackslash}X}}
 \textcolor{gregoriocolor}{A} & \textcolor{gregoriocolor}{B} & \textcolor{gregoriocolor}{C} & \textcolor{gregoriocolor}{D} & \textcolor{gregoriocolor}{E} & F & \textcolor{gregoriocolor}{F} & \textcolor{gregoriocolor}{G} & \textcolor{gregoriocolor}{H} & \textcolor{gregoriocolor}{M} & \textcolor{gregoriocolor}{N} & \textcolor{gregoriocolor}{P} \\
 8 & 9 & 10 & 11 & 12 & 13 & 12 & 13 & 14 & 15 & 16 & 17 \\
\end{tabularx}

\begin{paracol}{2}
\selectlanguage{latin}
\lettrine[lines=2]{C}{oncéptio} Immaculáta gloriósæ semper Vírginis Genitrícis Dei Maríæ, quam fuísse 
 præservátam, singulári Dei privilégio, ab omni originális culpæ labe immúnem, 
 Pius Nonus, Póntifex Máximus, hac ipsa recurrénte die, solémniter definívit.
\switchcolumn
\selectlanguage{english}
\lettrine[lines=2]{T}{he} Immaculate 
 Conception of the glorious and ever Virgin Mary, Mother of God. On 
 this day, Pius IX solemnly declared her to have been by a singular privilege 
 of God preserved from all stain of original sin.
\switchcolumn*
\selectlanguage{latin}
Tréviris sancti 
 Euchárii, qui fuit discípulus beáti Petri Apóstoli et primus ejúsdem 
 civitátis Epíscopus.
\switchcolumn
\selectlanguage{english}
At Treves, St. Eucharius, a disciple 
 of blessed Peter the Apostle, first bishop of that city.
\switchcolumn*
\selectlanguage{latin}
Alexandríæ sancti 
 Macárii Mártyris, qui, témpore Décii, cum a Júdice multis verbis ad negándum 
 Christum suaderétur, et eo majóri constántia suam profiterétur fidem, vivus 
 ad últimum exúri jubétur.
\switchcolumn
\selectlanguage{english}
At Alexandria, St. Macarius, martyr, 
 whose constancy in professing the faith increased with the efforts made by 
 the judge to persuade him to deny Christ. He was finally condemned to 
 be burned alive.
\switchcolumn*
\selectlanguage{latin}
In Cypro sancti 
 Sophrónii Epíscopi, qui pupíllórum, orphanórum ac viduárum defénsor 
 miríficus, et páuperum atque oppressórum ómnium adjútor fuit.
\switchcolumn
\selectlanguage{english}
In Cyprus, the holy bishop 
 Sophronius, who was a devoted protector of orphans and widows, and a helper 
 of the poor and oppressed.
\switchcolumn*
\selectlanguage{latin}
In monastério 
 Luxoviénsi, in Gállia, sancti Romárici Abbátis, qui, cum in aula Theodobérti 
 Regis primus esset, renuntiávit sæculo, et monásticæ étiam observántiæ laude 
 céteris antecélluit.
\switchcolumn
\selectlanguage{english}
In the monastery of Luxeuil in 
 France, St. Romaricus, abbot, who left the highest station at the court of 
 King Theodobert, renounced the world, and surpassed others in the observance 
 of monastic discipline.
\switchcolumn*
\selectlanguage{latin}
Constantinópoli sancti 
 Patápii Solitárii, virtútibus et miráculis clari.
\switchcolumn
\selectlanguage{english}
At Constantinople, St. Patapius, 
 solitary, renowned for virtues and miracles.
\switchcolumn*
\selectlanguage{latin}
Romæ Invéntio sanctórum 
 Mártyrum Nemésii Diáconi, ejúsque fíliæ Lucíllæ Vírginis, Symphrónii, 
 Olympii Tribúni, hujúsque uxóris Exsupériæ et Theodúli fílii; quorum memória 
 octávo Kaléndas Septémbris recensétur.
\switchcolumn
\selectlanguage{english}
At Rome, the finding of the holy 
 martyrs Nemesis, a deacon, his daughter Lucina, a virgin, Symphronius, 
 Olympius the tribune and his wife Exuperia and his son Theodulus, whose 
 commemoration is made on the 25th of August.
\switchcolumn*
\selectlanguage{latin}
Verónæ Ordinátio sancti 
 Zenónis Epíscopi.
\switchcolumn
\selectlanguage{english}
At Verona, the ordination of St. 
 Zeno, bishop.
\switchcolumn*
\selectlanguage{latin}
\end{paracol}


% ---- martyrology/mart12/mart1209.htm
\needspace{10\baselineskip}
\begin{paracol}{2}
\selectlanguage{latin}
\begin{center}{\color{gregoriocolor} Quinto Idus Decémbris. 
 Luna\dots\ }\end{center}
\switchcolumn
\selectlanguage{english}
\begin{center}{\color{gregoriocolor} The 
 9\textsuperscript{th} Day of 
 December. The\dots\ Day of the Moon.}\end{center}
\end{paracol}

\noindent\begin{tabularx}{\linewidth}{*{19}{>{\centering\arraybackslash}X}}
 \textcolor{gregoriocolor}{a} & \textcolor{gregoriocolor}{b} & \textcolor{gregoriocolor}{c} & \textcolor{gregoriocolor}{d} & \textcolor{gregoriocolor}{e} & \textcolor{gregoriocolor}{f} & \textcolor{gregoriocolor}{g} & \textcolor{gregoriocolor}{h} & \textcolor{gregoriocolor}{i} & \textcolor{gregoriocolor}{k} & \textcolor{gregoriocolor}{l} & \textcolor{gregoriocolor}{m} & \textcolor{gregoriocolor}{n} & \textcolor{gregoriocolor}{p} & \textcolor{gregoriocolor}{q} & \textcolor{gregoriocolor}{r} & \textcolor{gregoriocolor}{s} & \textcolor{gregoriocolor}{t} & \textcolor{gregoriocolor}{u} \\
 19 & 20 & 21 & 22 & 23 & 24 & 25 & 26 & 27 & 28 & 29 & 1 & 2 & 3 & 4 & 5 & 6 & 7 & 8 \\
\end{tabularx}
\vspace{0.5\baselineskip}
\noindent\begin{tabularx}{\linewidth}{*{12}{>{\centering\arraybackslash}X}}
 \textcolor{gregoriocolor}{A} & \textcolor{gregoriocolor}{B} & \textcolor{gregoriocolor}{C} & \textcolor{gregoriocolor}{D} & \textcolor{gregoriocolor}{E} & F & \textcolor{gregoriocolor}{F} & \textcolor{gregoriocolor}{G} & \textcolor{gregoriocolor}{H} & \textcolor{gregoriocolor}{M} & \textcolor{gregoriocolor}{N} & \textcolor{gregoriocolor}{P} \\
 9 & 10 & 11 & 12 & 13 & 14 & 13 & 14 & 15 & 16 & 17 & 18 \\
\end{tabularx}

\begin{paracol}{2}
\selectlanguage{latin}
\lettrine[lines=2]{C}{arthágine} sancti 
 Restitúti, Epíscopi et Mártyris, in cujus solemnitáte sanctus Augustínus de 
 ipso ad pópulum sermónem hábuit.
\switchcolumn
\selectlanguage{english}
\lettrine[lines=2]{A}{t} Carthage, St. Restitutus, bishop 
 and martyr, on whose feast St. Augustine delivered a discourse to the people 
 in which he set forth his praises.
\switchcolumn*
\selectlanguage{latin}
Item in Africa 
 sanctórum Mártyrum Petri, Succéssi, Bassiáni, Primitívi et aliórum vigínti.
\switchcolumn
\selectlanguage{english}
Also in Africa, the holy martyrs 
 Peter, Successus, Bassian, Primitivus, and twenty others.
\switchcolumn*
\selectlanguage{latin}
Toléti, in Hispánia, 
 natális sanctæ Leocádiæ, Vírginis et Mártyris; quæ, in persecutióne 
 Diocletiáni Imperatóris, a Præfécto Hispaniárum Daciáno inclúsa cárcere ac 
 dire maceráta, in eo tandem, cum gravíssimos beátæ Euláliæ et reliquórum 
 Mártyrum cruciátus audísset, impollútum spíritum, génibus in oratióne 
 pósitis, Christo réddidit.
\switchcolumn
\selectlanguage{english}
At Toledo in Spain, the birthday of 
 the holy virgin Leocadia, a martyr in the persecution of Emperor Diocletian. 
 She was condemned to a cruel imprisonment by Dacian, prefect of Spain, and 
 was pining away when, hearing of the barbarous tortures of blessed Eulalia 
 and the other martyrs, she knelt down to pray and yielded up her undefiled 
 spirit to Christ.
\switchcolumn*
\selectlanguage{latin}
Lemóvicis, in Aquitánia, sanctæ Valériæ, Vírginis et Mártyris.
\switchcolumn
\selectlanguage{english}
At Limoges in Aquitaine, St. 
 Valeria, virgin and martyr.
\switchcolumn*
\selectlanguage{latin}
Verónæ sancti Próculi 
 Epíscopi, qui, in persecutióne Diocletiáni, cólaphis ac fústibus cæsus, e 
 civitáte pulsus est, ac tandem in Ecclésiam suam restitútus, quiévit in 
 pace.
\switchcolumn
\selectlanguage{english}
At Verona, during the persecution of 
 Diocletian, St. Proculus, bishop, who was buffeted, scourged with rods, and 
 driven out of the city. Being at length restored to his church, he 
 died in peace.
\switchcolumn*
\selectlanguage{latin}
Papíæ sancti Syri, qui 
 fuit primus ejúsdem civitátis Epíscopus, atque apostólicis signis et 
 virtútibus cláruit.
\switchcolumn
\selectlanguage{english}
At Pavia, St. Syrus, first bishop of 
 that city, who was renowned for apostolic signs and virtues.
\switchcolumn*
\selectlanguage{latin}
Apaméæ, in Syria, beáti 
 Juliáni Epíscopi, qui, témpore Sevéri, sanctitáte refúlsit.
\switchcolumn
\selectlanguage{english}
At Apamea in Syria, blessed Julian, 
 bishop, who flourished in holiness in the time of Severus.
\switchcolumn*
\selectlanguage{latin}
Graji, in Burgúndia, 
 sancti Petri Fourier qui Canónicus Reguláris fuit Salvatóris Nostri, et 
 Canonissárum Regulárium Dóminæ Nostræ edocéndis puéllis Institútor; atque, 
 virtútibus ac miráculis clarus, a Leóne Décimo tértio, Pontífice Máximo, 
 Sanctórum catálogo adjúnctus est.
\switchcolumn
\selectlanguage{english}
At Gray in Burgundy, St. Peter 
 Fourier, Canon Regular of Our Saviour and the founder of the Canonesses 
 Regular of Our Lady for the education of children. Because of his 
 brilliant virtues and miracles, Leo XIII placed him the catalogue of the 
 Saints.
\switchcolumn*
\selectlanguage{latin}
Petragóricis, in 
 Gállia, sancti Cypriáni Abbátis, magnæ sanctitátis viri.
\switchcolumn
\selectlanguage{english}
At Perigueux in France, St. Cyprian, 
 abbot, a man of great sanctity.
\switchcolumn*
\selectlanguage{latin}
Naziánzi, in Cappadócia, 
 sanctæ Gorgóniæ, quæ fuit beátæ Nonnæ fília, atque beatórum Gregórii 
 Theólogi et Cæsárii soror, cujus ipse Gregórius virtútes et mirácula 
 conscrípsit.
\switchcolumn
\selectlanguage{english}
At Nazianzum in Cappadocia, St. 
 Gorgonia, of whose virtues and miracles St. Gregory has written. She 
 was the daughter of blessed Nonna and the sister of St. Gregory the 
 Theologian and St. Caesarius.
\switchcolumn*
\selectlanguage{latin}
\end{paracol}


% ---- martyrology/mart12/mart1210.htm
\needspace{10\baselineskip}
\begin{paracol}{2}
\selectlanguage{latin}
\begin{center}{\color{gregoriocolor} Quarto Idus Decémbris. 
 Luna\dots\ }\end{center}
\switchcolumn
\selectlanguage{english}
\begin{center}{\color{gregoriocolor} The 
 10\textsuperscript{th} Day of 
 December. The\dots\ Day of the Moon.}\end{center}
\end{paracol}

\noindent\begin{tabularx}{\linewidth}{*{19}{>{\centering\arraybackslash}X}}
 \textcolor{gregoriocolor}{a} & \textcolor{gregoriocolor}{b} & \textcolor{gregoriocolor}{c} & \textcolor{gregoriocolor}{d} & \textcolor{gregoriocolor}{e} & \textcolor{gregoriocolor}{f} & \textcolor{gregoriocolor}{g} & \textcolor{gregoriocolor}{h} & \textcolor{gregoriocolor}{i} & \textcolor{gregoriocolor}{k} & \textcolor{gregoriocolor}{l} & \textcolor{gregoriocolor}{m} & \textcolor{gregoriocolor}{n} & \textcolor{gregoriocolor}{p} & \textcolor{gregoriocolor}{q} & \textcolor{gregoriocolor}{r} & \textcolor{gregoriocolor}{s} & \textcolor{gregoriocolor}{t} & \textcolor{gregoriocolor}{u} \\
 20 & 21 & 22 & 23 & 24 & 25 & 26 & 27 & 28 & 29 & 1 & 2 & 3 & 4 & 5 & 6 & 7 & 8 & 9 \\
\end{tabularx}
\vspace{0.5\baselineskip}
\noindent\begin{tabularx}{\linewidth}{*{12}{>{\centering\arraybackslash}X}}
 \textcolor{gregoriocolor}{A} & \textcolor{gregoriocolor}{B} & \textcolor{gregoriocolor}{C} & \textcolor{gregoriocolor}{D} & \textcolor{gregoriocolor}{E} & F & \textcolor{gregoriocolor}{F} & \textcolor{gregoriocolor}{G} & \textcolor{gregoriocolor}{H} & \textcolor{gregoriocolor}{M} & \textcolor{gregoriocolor}{N} & \textcolor{gregoriocolor}{P} \\
 10 & 11 & 12 & 13 & 14 & 15 & 14 & 15 & 16 & 17 & 18 & 19 \\
\end{tabularx}

\begin{paracol}{2}
\selectlanguage{latin}
\lettrine[lines=2]{S}{ancti} Melchíadis, Papæ 
 et Mártyris, cujus dies natális recensétur tértio Idus Januárii.
\switchcolumn
\selectlanguage{english}
\lettrine[lines=2]{S}{t.} Melchiades, pope and martyr, 
 whose birthday is mentioned on the 11th of January.
\switchcolumn*
\selectlanguage{latin}
Romæ, via Ostiénsi, 
 Dedicátio Basílicæ sancti Pauli Apóstoli; quæ, simul cum Dedicatióne 
 Basílicæ sancti Petri, Apostolórum Príncipis, ánnua celebritáte recólitur 
 quartodécimo Kaléndas Decémbris.
\switchcolumn
\selectlanguage{english}
At Rome, on the Ostian Way, the 
 Dedication of the Basilica of St. Paul the Apostle. The yearly 
 commemoration of this Dedication, together with that of St. Peter, prince of 
 the apostles, is observed on the 18th of November.
\switchcolumn*
\selectlanguage{latin}
Eódem die sanctórum 
 Mártyrum Carpóphori Presbyteri, et Abúndii Diáconi; qui, in Diocletiáni 
 persecutióne, primo fústibus crudelíssime cæsi, deínde in cárcerem, negáto 
 cibo et potu, retrúsi, et rursum in equúleo torti, et post hæc diu in cárcere maceráti, novíssime gládio percússi sunt.
\switchcolumn
\selectlanguage{english}
Also, the holy martyrs Carpophorus, 
 a priest, and Abundius, a deacon, in the persecution of Diocletian. 
 They were first cruelly beaten with rods, then imprisoned and denied food 
 and drink; being placed on the rack a second time and again thrown into 
 prison, they were finally beheaded.
\switchcolumn*
\selectlanguage{latin}
Alexandríæ sanctórum 
 Mártyrum Mennæ, Hermógenis et Eugraphi; qui sub Galério Maximiáno passi sunt.
\switchcolumn
\selectlanguage{english}
At Alexandria, the holy martyrs 
 Mennas, Hermogenes, and Eugraphus, who suffered under Galerius Maximian.
\switchcolumn*
\selectlanguage{latin}
Apud Leontínos, in 
 Sicília, sanctórum Mártyrum Mercúrii et Sociórum mílitum; qui, sub Tertyllo 
 Præside, témpore Licínii Imperatóris, gládio cæsi sunt.
\switchcolumn
\selectlanguage{english}
At Lentini in Sicily, the holy 
 martyrs Mercurius and his soldier companions, who were slain by the sword 
 under the governor Tertyllus, in the reign of Emperor Licinius.
\switchcolumn*
\selectlanguage{latin}
Ancyræ, in Galátia, sancti Gemélli Mártyris, qui, post dira torménta, sub Juliáno Apóstata, 
 crucis supplício martyrium consummávit.
\switchcolumn
\selectlanguage{english}
At Ancyra in Galatia, St. Gemellus, 
 martyr, who, after severe torments, fulfilled his martyrdom by being 
 crucified in the time of Julian the Apostate.
\switchcolumn*
\selectlanguage{latin}
Eméritæ, in Hispánia, 
 pássio sanctæ Euláliæ Vírginis, quæ, sub Maximiáno Imperatóre, cum esset 
 annórum duódecim, ibi, jussu Daciáni Præsidis, pro confessióne Christi, 
 plúrima torménta est perpéssa; novíssime, in equúleo suspénsa et exunguláta, 
 fáculis ardéntibus ex utróque látere appósitis, hausto igne, spíritum 
 réddidit.
\switchcolumn
\selectlanguage{english}
At Merida in Spain, in the time of 
 Maximian, the martyrdom of the holy virgin Eulalia, who at twelve years of 
 age suffered many torments for the confession of Christ by order of the 
 governor Dacian. She was stretched on the rack, torn with iron claws, 
 had her sides burned with flaming torches, and swallowing the fire she 
 expired.
\switchcolumn*
\selectlanguage{latin}
Item ibídem sanctæ 
 Júliæ, Vírginis et Mártyris; quæ beátæ Euláliæ sócia fuit, et illi ad 
 passiónem properánti indivídua comes adhæsit.
\switchcolumn
\selectlanguage{english}
Also, in the same city, St. Julia, 
 virgin and martyr, the companion of the blessed Eulalia, who would not be 
 separated from her when the latter went to suffer.
\switchcolumn*
\selectlanguage{latin}
Romæ beáti Gregórii 
 Papæ Tértii, qui sanctitáte meritísque præclárus migrávit in cælum.
\switchcolumn
\selectlanguage{english}
At Rome, Pope St. Gregory III, who 
 departed for heaven renowned for his sanctity and good works.
\switchcolumn*
\selectlanguage{latin}
Viénnæ, in Gállia, 
 sancti Sindúlphi, Epíscopi et Confessóris.
\switchcolumn
\selectlanguage{english}
At Vienne in France, St. Sindulph, 
 bishop and confessor.
\switchcolumn*
\selectlanguage{latin}
Bríxiæ sancti Deúsdedit 
 Epíscopi.
\switchcolumn
\selectlanguage{english}
At Brescia, St. Deusdedit, bishop.
\switchcolumn*
\selectlanguage{latin}
Lauréti, in Picéno, Translátio sacræ Domus Genitrícis Dei Maríæ, qua in domo Verbum caro factum 
 est. Ipsam vero beatíssimam Vírginem, Lauretánæ título nuncupátam, 
 Benedíctus Papa Décimus quintus ómnibus aereonáutis præcípuam apud Deum 
 Patrónam attríbuit.
\switchcolumn
\selectlanguage{english}
At Loreto in Piceno, the 
 Translation of the Holy House of Mary the Mother of God, wherein the Word 
 was made flesh. Pope Benedict XV declared the same Blessed Virgin 
 Mary, under the title of Loreto, to be the chief Patroness before God of 
 all airmen.
\switchcolumn*
\selectlanguage{latin}
\end{paracol}


% ---- martyrology/mart12/mart1211.htm
\needspace{10\baselineskip}
\begin{paracol}{2}
\selectlanguage{latin}
\begin{center}{\color{gregoriocolor} Tértio Idus Decémbris. 
 Luna\dots\ }\end{center}
\switchcolumn
\selectlanguage{english}
\begin{center}{\color{gregoriocolor} The 
 11\textsuperscript{th} Day of 
 December. The\dots\ Day of the Moon.}\end{center}
\end{paracol}

\noindent\begin{tabularx}{\linewidth}{*{19}{>{\centering\arraybackslash}X}}
 \textcolor{gregoriocolor}{a} & \textcolor{gregoriocolor}{b} & \textcolor{gregoriocolor}{c} & \textcolor{gregoriocolor}{d} & \textcolor{gregoriocolor}{e} & \textcolor{gregoriocolor}{f} & \textcolor{gregoriocolor}{g} & \textcolor{gregoriocolor}{h} & \textcolor{gregoriocolor}{i} & \textcolor{gregoriocolor}{k} & \textcolor{gregoriocolor}{l} & \textcolor{gregoriocolor}{m} & \textcolor{gregoriocolor}{n} & \textcolor{gregoriocolor}{p} & \textcolor{gregoriocolor}{q} & \textcolor{gregoriocolor}{r} & \textcolor{gregoriocolor}{s} & \textcolor{gregoriocolor}{t} & \textcolor{gregoriocolor}{u} \\
 21 & 22 & 23 & 24 & 25 & 26 & 27 & 28 & 29 & 1 & 2 & 3 & 4 & 5 & 6 & 7 & 8 & 9 & 10 \\
\end{tabularx}
\vspace{0.5\baselineskip}
\noindent\begin{tabularx}{\linewidth}{*{12}{>{\centering\arraybackslash}X}}
 \textcolor{gregoriocolor}{A} & \textcolor{gregoriocolor}{B} & \textcolor{gregoriocolor}{C} & \textcolor{gregoriocolor}{D} & \textcolor{gregoriocolor}{E} & F & \textcolor{gregoriocolor}{F} & \textcolor{gregoriocolor}{G} & \textcolor{gregoriocolor}{H} & \textcolor{gregoriocolor}{M} & \textcolor{gregoriocolor}{N} & \textcolor{gregoriocolor}{P} \\
 11 & 12 & 13 & 14 & 15 & 16 & 15 & 16 & 17 & 18 & 19 & 20 \\
\end{tabularx}

\begin{paracol}{2}
\selectlanguage{latin}
\lettrine[lines=2]{R}{omæ} sancti Dámasi 
 Primi, Papæ et Confessóris; qui Apollinárem hæresiárcham damnávit, et Petrum, 
 Episcopum Alexandrínum, fugátum restítuit; multa étiam sanctórum Mártyrum 
 córpora invénit, eorúmque memórias vérsibus exornávit.
\switchcolumn
\selectlanguage{english}
\lettrine[lines=2]{A}{t} Rome, St. Damasus, pope and 
 confessor, who condemned the heresiarch Apollinaris, and restored to his See 
 Peter, bishop of Alexandria, who had been driven from it. He also 
 discovered the bodies of many holy martyrs and composed verses in their 
 honour.
\switchcolumn*
\selectlanguage{latin}
Item Romæ pássio sancti 
 Trasónis, qui, cum Christiános laborántes in thermis, aliísque opéribus 
 públicis fatigátos, et in cárcere pósitos, de suis facultátibus áleret, 
 jussu Maximiáni tentus est, et cum áliis duóbus, id est Pontiáno et 
 Prætextáto, martyrio coronátus.
\switchcolumn
\selectlanguage{english}
Also at Rome, St. Thrason. He 
 was arrested by order of Maximian for supporting with his goods the 
 Christians who laboured in the baths and at other public works, and those 
 confined in jail. He was crowned with martyrdom with two others, Pontian and 
 Prætextatus.
\switchcolumn*
\selectlanguage{latin}
Ambiáni, in Gállia, 
 sanctórum Mártyrum Victórici et Fusciáni, sub eódem Imperatóre, in quorum 
 náribus et áuribus jussit Rictiovárus Præses immítti tarínchas, et clavis 
 ardéntibus témpora transfígi, deínde óculos evélli, ac póstmodum eórum 
 córpora jaculári; sicque, una cum sancto Gentiáno, eórum hóspite, capítibus 
 amputátis, migravérunt ad Dóminum.
\switchcolumn
\selectlanguage{english}
At Amiens in France, the holy 
 martyrs Victoricus and Fuscian, under the same emperor. By order of 
 Governor Rictiovarus, they had iron pins driven into their ears and 
 nostrils, heated nails into their temples, and arrows into their bodies and 
 their eyes torn out. They were beheaded with St. Gentian, their guest, 
 and they passed to the Lord.
\switchcolumn*
\selectlanguage{latin}
In Pérside sancti 
 Bársabæ Mártyris.
\switchcolumn
\selectlanguage{english}
In Persia, St. Barbabas, martyr.
\switchcolumn*
\selectlanguage{latin}
In Hispánia sancti 
 Eutychii Mártyris.
\switchcolumn
\selectlanguage{english}
In Spain, St. Eutychius, martyr.
\switchcolumn*
\selectlanguage{latin}
Placéntiæ sancti Sabíni 
 Epíscopi, miráculis clari.
\switchcolumn
\selectlanguage{english}
At Piacenza, St. Sabinus, bishop, 
 renowned for miracles.
\switchcolumn*
\selectlanguage{latin}
Constantinópoli sancti 
 Daniélis Stylítæ.
\switchcolumn
\selectlanguage{english}
At Constantinople, St. Daniel 
 Stylites.
\switchcolumn*
\selectlanguage{latin}
\end{paracol}


% ---- martyrology/mart12/mart1212.htm
\needspace{10\baselineskip}
\begin{paracol}{2}
\selectlanguage{latin}
\begin{center}{\color{gregoriocolor} Prídie Idus Decémbris. 
 Luna\dots\ }\end{center}
\switchcolumn
\selectlanguage{english}
\begin{center}{\color{gregoriocolor} The 
 12\textsuperscript{th} Day of 
 December. The\dots\ Day of the Moon.}\end{center}
\end{paracol}

\noindent\begin{tabularx}{\linewidth}{*{19}{>{\centering\arraybackslash}X}}
 \textcolor{gregoriocolor}{a} & \textcolor{gregoriocolor}{b} & \textcolor{gregoriocolor}{c} & \textcolor{gregoriocolor}{d} & \textcolor{gregoriocolor}{e} & \textcolor{gregoriocolor}{f} & \textcolor{gregoriocolor}{g} & \textcolor{gregoriocolor}{h} & \textcolor{gregoriocolor}{i} & \textcolor{gregoriocolor}{k} & \textcolor{gregoriocolor}{l} & \textcolor{gregoriocolor}{m} & \textcolor{gregoriocolor}{n} & \textcolor{gregoriocolor}{p} & \textcolor{gregoriocolor}{q} & \textcolor{gregoriocolor}{r} & \textcolor{gregoriocolor}{s} & \textcolor{gregoriocolor}{t} & \textcolor{gregoriocolor}{u} \\
 22 & 23 & 24 & 25 & 26 & 27 & 28 & 29 & 1 & 2 & 3 & 4 & 5 & 6 & 7 & 8 & 9 & 10 & 11 \\
\end{tabularx}
\vspace{0.5\baselineskip}
\noindent\begin{tabularx}{\linewidth}{*{12}{>{\centering\arraybackslash}X}}
 \textcolor{gregoriocolor}{A} & \textcolor{gregoriocolor}{B} & \textcolor{gregoriocolor}{C} & \textcolor{gregoriocolor}{D} & \textcolor{gregoriocolor}{E} & F & \textcolor{gregoriocolor}{F} & \textcolor{gregoriocolor}{G} & \textcolor{gregoriocolor}{H} & \textcolor{gregoriocolor}{M} & \textcolor{gregoriocolor}{N} & \textcolor{gregoriocolor}{P} \\
 12 & 13 & 14 & 15 & 16 & 17 & 16 & 17 & 18 & 19 & 20 & 21 \\
\end{tabularx}

\begin{paracol}{2}
\selectlanguage{latin}
\lettrine[lines=2]{A}{lexandríæ} sanctórum Epímachi et Alexándri, qui, sub
  Décio Imperatóre, cum fuíssent multo témpore in vínculis, atque, divérsis
  supplíciis affécti, perdurássent in fide, ígnibus tandem consúmpti sunt.
  Sanctus vero Epímachus, simul cum sancto Gordiáno Mártyre, festíva
  celebritáte cólitur sexto Idus Maji.
\switchcolumn
\selectlanguage{english}
\lettrine[lines=2]{A}{t} Alexandria, in the time of Decius, 
 the holy martyrs Epimachus and Alexander, who were kept in chains a long 
 time and subjected to various torments, but as they persevered in the faith, 
 they were finally consumed by fire. The feast of St. Epimachus 
 together with that of St. Gordian the martyr is observed on the 10th of May.
\switchcolumn*
\selectlanguage{latin}
Romæ sancti Synésii 
 Mártyris, qui, beáti Xysti Papæ Secúndi témpore ordinátus Lector, et, cum 
 multos convertísset ad Christum, apud Aureliánum Imperatórem accusátus, 
 martyrii corónam gládio percússus accépit.
\switchcolumn
\selectlanguage{english}
At Rome, the holy martyr Synesius, 
 who was ordained lector in the time of blessed Pope Sixtus. Having 
 converted many to Christ, he was accused before Emperor Aurelian, and being 
 put to the sword, received the crown of martyrdom.
\switchcolumn*
\selectlanguage{latin}
Eódem die sanctórum 
 Mártyrum Hermógenis, Donáti et aliórum vigínti duórum.
\switchcolumn
\selectlanguage{english}
On the same day, the holy martyrs 
 Hermogenes, Donatus, and twenty-two others.
\switchcolumn*
\selectlanguage{latin}
Tréviris sanctórum 
 Mártyrum Maxéntii, Constántii, Crescéntii, Justíni et Sociórum; qui in 
 persecutióne Diocletiáni, sub Rictiováro Præside, passi sunt.
\switchcolumn
\selectlanguage{english}
At Treves, the holy martyrs 
 Maxentius, Constantius, Crescentius, Justinus, and their companions, who 
 suffered in the persecution of Diocletian, under the governor Rictiovarus.
\switchcolumn*
\selectlanguage{latin}
Alexandríæ sanctárum Ammonáriæ Vírginis, Mercúriæ, Dionysiæ et altérius
  Ammonáriæ. Harum prima, in persecutióne Décii, inaudítis tormentórum
  genéribus superátis, beátum vitæ finem, ferro cædénte, percépit; tres vero
  áliæ, cum Judex a féminis superári erubésceret, ac dubitáret, ne, éadem in
  illas exércens cruciaménta, viríli eárum étiam constántia vincerétur, statim
  jussæ sunt decollári.
\switchcolumn
\selectlanguage{english}
At Alexandria, the holy women 
 Ammonaria, virgin, Mercuria, Dionysia, and another Ammonaria. The 
 first named, after having triumphed over unheard-of kinds of torments, in 
 the persecution of Decius, ended her blessed life by beheading. As to 
 the three others, the judge, being ashamed to be overcome by women, and 
 fearing that by resorting to tortures he would be vanquished by their 
 constancy, ordered them to be beheaded immediately.
\switchcolumn*
\selectlanguage{latin}
\end{paracol}


% ---- martyrology/mart12/mart1213.htm
\needspace{10\baselineskip}
\begin{paracol}{2}
\selectlanguage{latin}
\begin{center}{\color{gregoriocolor} Idibus Decémbris. 
 Luna\dots\ }\end{center}
\switchcolumn
\selectlanguage{english}
\begin{center}{\color{gregoriocolor} The 
 13\textsuperscript{th} Day of 
 December. The\dots\ Day of the Moon.}\end{center}
\end{paracol}

\noindent\begin{tabularx}{\linewidth}{*{19}{>{\centering\arraybackslash}X}}
 \textcolor{gregoriocolor}{a} & \textcolor{gregoriocolor}{b} & \textcolor{gregoriocolor}{c} & \textcolor{gregoriocolor}{d} & \textcolor{gregoriocolor}{e} & \textcolor{gregoriocolor}{f} & \textcolor{gregoriocolor}{g} & \textcolor{gregoriocolor}{h} & \textcolor{gregoriocolor}{i} & \textcolor{gregoriocolor}{k} & \textcolor{gregoriocolor}{l} & \textcolor{gregoriocolor}{m} & \textcolor{gregoriocolor}{n} & \textcolor{gregoriocolor}{p} & \textcolor{gregoriocolor}{q} & \textcolor{gregoriocolor}{r} & \textcolor{gregoriocolor}{s} & \textcolor{gregoriocolor}{t} & \textcolor{gregoriocolor}{u} \\
 23 & 24 & 25 & 26 & 27 & 28 & 29 & 1 & 2 & 3 & 4 & 5 & 6 & 7 & 8 & 9 & 10 & 11 & 12 \\
\end{tabularx}
\vspace{0.5\baselineskip}
\noindent\begin{tabularx}{\linewidth}{*{12}{>{\centering\arraybackslash}X}}
 \textcolor{gregoriocolor}{A} & \textcolor{gregoriocolor}{B} & \textcolor{gregoriocolor}{C} & \textcolor{gregoriocolor}{D} & \textcolor{gregoriocolor}{E} & F & \textcolor{gregoriocolor}{F} & \textcolor{gregoriocolor}{G} & \textcolor{gregoriocolor}{H} & \textcolor{gregoriocolor}{M} & \textcolor{gregoriocolor}{N} & \textcolor{gregoriocolor}{P} \\
 13 & 14 & 15 & 16 & 17 & 18 & 17 & 18 & 19 & 20 & 21 & 22 \\
\end{tabularx}

\begin{paracol}{2}
\selectlanguage{latin}
\lettrine[lines=2]{S}{yracúsis,} in Sicília, 
 natális sanctæ Lúciæ, Vírginis et Mártyris, in persecutióne Diocletiáni. 
 Hæc nóbilis Virgo, cum eam lenónes, quibus, jubénte, Paschásio Consulári, 
 trádita erat ut a pópulo castitáti ejus illuderétur, dúcere vellent, 
 nullátenus per illos movéri pótuit, nec fúnibus ádditis, nec boum jugis 
 plúrimis; deínde vero, picem, resínam ac fervens óleum nil læsa súperans, 
 tandem, gládio in gútture percússa, martyrium consummávit.
\switchcolumn
\selectlanguage{english}
\lettrine[lines=2]{A}{t} Syracuse in Sicily, the birthday 
 of St. Lucy, virgin and martyr, in the persecution of Diocletian. By 
 order of the proconsul Paschasius, she was delivered to profligates, that 
 her chastity might be insulted by the people; but when they attempted to 
 lead her away they were not able to move her, either with ropes or by means 
 of many yoke of oxen. Then having hot pitch, resin, and burning oil 
 applied to her body without being injured, she finally had a sword driven 
 through her throat, and thus completed her martyrdom.
\switchcolumn*
\selectlanguage{latin}
Molíni, in Gállia, item natális sanctæ Joánnæ-Francíscæ Frémiot de Chantal
  Víduæ, quæ Ordinis Sanctimoniálium Visitatiónis sanctæ Maríæ fuit
  Institútrix; ac, nobilitáte géneris, vitæ sanctimónia, quam in quadrúplici
  statu constánter duxit, et miraculórum dono illústris, a Cleménte Décimo
  tértio, Pontífice Máximo, in Sanctárum númerum reláta est. Sacrum ejus corpus
  Annésium, in Sabáudia, translátum fuit, et solémni pompa in prima sui Ordinis
  Ecclésia tumulátum. Ipsíus autem festum duodécimo Kaléndas Septémbris ab
  univérsa Ecclésia Clemens Papa Décimus quartus celebrári mandávit.
\switchcolumn
\selectlanguage{english}
At Moulins in France, the birthday 
 of St. Jane Frances Fremiot de Chantal, widow, foundress of the Nuns of the 
 Visitation of St. Mary, distinguished by the nobility of her birth, by the 
 holiness she constantly displayed in four different states of life, and by 
 the gift of miracles. She was placed among the saints by Clement XIII. 
 Her holy body was taken to Annecy in Savoy and buried with great pomp in the 
 first church of her order. by order of Clement XIV, her feast is kept 
 by the whole Church on the 21st of August.
\switchcolumn*
\selectlanguage{latin}
In Arménia pássio 
 sanctórum Mártyrum Eustrátii, Auxéntii, Eugénii, Mardárii et Oréstis, in 
 persecutióne Diocletiáni. Ex ipsis Eustrátius, primum sub Lysia, 
 deínde Sebáste, sub Agricoláo Præside, una cum Oréste exquisítis torméntis 
 cruciátus, in fornácem missus réddidit spíritum; Oréstes autem, in stratum 
 férreum ignítum pósitus, migrávit ad Dóminum; céteri apud Arábracos, 
 sævíssimis agitáti supplíciis, sub Lysia Præside, diversímode martyrium 
 consummárunt. Eórum córpora, póstea Romam transláta, in Ecclésia 
 sancti Apollináris honorífice collocáta sunt.
\switchcolumn
\selectlanguage{english}
In Armenia, the martyrdom of the 
 holy martyrs Eustratius, Auxentius, Eugene, Mardarius, and Orestes, in the 
 persecution of Diocletian. Eustratius was the first subjected alone to 
 barbarous torments under Lysias. Then he was conducted to Sebaste, 
 where he was tortured together with Orestes under the governor Agricolaus, 
 and being cast into a furnace, yielded up his soul; but Orestes being laid 
 on a bed of heated iron, rendered his soul unto God. The others were 
 made to endure most grievous torments among the Arabraci, under the governor 
 Lysias, and fulfilled their martyrdom in different ways. Their relics 
 were afterwards carried to Rome and placed with due honours in the church of 
 St. Apollinaris.
\switchcolumn*
\selectlanguage{latin}
In Sulcitána ínsula, 
 apud Sardíniam, pássio sancti Antíochi, sub Hadriáno Imperatóre.
\switchcolumn
\selectlanguage{english}
At Sardinia, in the island of Sulci, 
 the martyrdom of St. Antiochus, under Emperor Hadrian.
\switchcolumn*
\selectlanguage{latin}
Cameráci, in Gállia, 
 sancti Authbérti, Epíscopi et Confessóris.
\switchcolumn
\selectlanguage{english}
At Cambrai in France, St. Aubert, 
 bishop and confessor.
\switchcolumn*
\selectlanguage{latin}
In pago Pontívo, in 
 Gállia, sancti Judóci, Presbyteri et Confessóris.
\switchcolumn
\selectlanguage{english}
In the parts of Ponthieu in France, 
 St. Judoc, priest and confessor.
\switchcolumn*
\selectlanguage{latin}
In território 
 Argentoraténsi sanctæ Othíliæ Vírginis.
\switchcolumn
\selectlanguage{english}
In the territory of Strasbourg, St. 
 Otilie, virgin.
\switchcolumn*
\selectlanguage{latin}
\end{paracol}


% ---- martyrology/mart12/mart1214.htm
\needspace{10\baselineskip}
\begin{paracol}{2}
\selectlanguage{latin}
\begin{center}{\color{gregoriocolor} Décimo nono Kaléndas Januárii. 
 Luna\dots\ }\end{center}
\switchcolumn
\selectlanguage{english}
\begin{center}{\color{gregoriocolor} The 
 14\textsuperscript{th} Day of 
 December. The\dots\ Day of the Moon.}\end{center}
\end{paracol}

\noindent\begin{tabularx}{\linewidth}{*{19}{>{\centering\arraybackslash}X}}
 \textcolor{gregoriocolor}{a} & \textcolor{gregoriocolor}{b} & \textcolor{gregoriocolor}{c} & \textcolor{gregoriocolor}{d} & \textcolor{gregoriocolor}{e} & \textcolor{gregoriocolor}{f} & \textcolor{gregoriocolor}{g} & \textcolor{gregoriocolor}{h} & \textcolor{gregoriocolor}{i} & \textcolor{gregoriocolor}{k} & \textcolor{gregoriocolor}{l} & \textcolor{gregoriocolor}{m} & \textcolor{gregoriocolor}{n} & \textcolor{gregoriocolor}{p} & \textcolor{gregoriocolor}{q} & \textcolor{gregoriocolor}{r} & \textcolor{gregoriocolor}{s} & \textcolor{gregoriocolor}{t} & \textcolor{gregoriocolor}{u} \\
 24 & 25 & 26 & 27 & 28 & 29 & 1 & 2 & 3 & 4 & 5 & 6 & 7 & 8 & 9 & 10 & 11 & 12 & 13 \\
\end{tabularx}
\vspace{0.5\baselineskip}
\noindent\begin{tabularx}{\linewidth}{*{12}{>{\centering\arraybackslash}X}}
 \textcolor{gregoriocolor}{A} & \textcolor{gregoriocolor}{B} & \textcolor{gregoriocolor}{C} & \textcolor{gregoriocolor}{D} & \textcolor{gregoriocolor}{E} & F & \textcolor{gregoriocolor}{F} & \textcolor{gregoriocolor}{G} & \textcolor{gregoriocolor}{H} & \textcolor{gregoriocolor}{M} & \textcolor{gregoriocolor}{N} & \textcolor{gregoriocolor}{P} \\
 14 & 15 & 16 & 17 & 18 & 19 & 18 & 19 & 20 & 21 & 22 & 23 \\
\end{tabularx}

\begin{paracol}{2}
\selectlanguage{latin}
\lettrine[lines=2]{U}{bédæ,} in Hispánia, 
 natális sancti Joánnis a Cruce, Presbyteri et Confessóris, sanctæ Terésiæ in 
 Carmelitárum reformatióne sócii; quem, a Summo Pontífice Benedícto Décimo 
 tértio Sanctis adscríptum, Pius Papa Undécimus Doctórem universális Ecclésiæ 
 declarávit. Ipsíus tamen festívitas ágitur octávo Kaléndas Decémbris.
\switchcolumn
\selectlanguage{english}
\lettrine[lines=2]{A}{t} Ubeda in Spain, the birthday of 
 St. John of the Cross, priest and confessor, and the companion of St. Teresa 
 in the reform of the Carmelites. Pope Benedict XIII placed him on the 
 list of the saints, and Pope Pius XI declared him a doctor of the universal 
 Church. His feast, however, is observed on the 24th of November.
\switchcolumn*
\selectlanguage{latin}
Rhemis, in Gállia, 
 pássio sanctórum Nicásii Epíscopi, ac soróris Eutrópiæ Vírginis, et Sociórum 
 Mártyrum; qui a bárbaris Ecclésiæ hóstibus cæsi sunt.
\switchcolumn
\selectlanguage{english}
At Rheims in France, holy Bishop 
 Nicasius, his sister, the virgin Eutropia, and their companions, martyrs, 
 who were put to death by barbarians hostile to the Church.
\switchcolumn*
\selectlanguage{latin}
Alexandríæ sanctórum 
 Mártyrum Herónis, Arsénii, Isidóri, et Dióscori púeri. Horum tres 
 primos Judex, in persecutióne Deciána, cum eos, váriis torméntis dilánians, 
 pari armátos constántia vidéret, tradi ígnibus jubet; Dióscorus vero, 
 multiplíciter flagellátus, divíno nutu ad consolatiónem fidélium dimíssus 
 est.
\switchcolumn
\selectlanguage{english}
At Alexandria, the holy martyrs 
 Heron, Arsenius, Isidore, and the boy Dioscorus. In the persecution of 
 Decius, the first three were subjected to all the refinements of cruelty by 
 the judge, who, seeing them displaying the same constancy, ordered that they 
 should be cast into the fire. But Dioscorus, after repeated scourgings, 
 was set free by the intervention of Providence to the great consolation of 
 the faithful.
\switchcolumn*
\selectlanguage{latin}
Antiochíæ natális 
 sanctórum Mártyrum Drusi, Zósimi et Theodóri.
\switchcolumn
\selectlanguage{english}
At Antioch, the birthday of the holy 
 martyrs Drusus, Zosimus, and Theodore.
\switchcolumn*
\selectlanguage{latin}
Eódem die pássio 
 sanctórum Justi et Abúndii, qui, sub Numeriáno Imperatóre et Olybrio Præside, 
 conjécti in ignem, et, cum inde evasíssent illæsi, gládio percússi sunt.
\switchcolumn
\selectlanguage{english}
On the same day, the martyrdom of 
 Saints Justus and Abundius, who were cast into the flames in the time of 
 Emperor Numerian and the governor Olybrius, but escaping all injury, they 
 were smitten with the sword.
\switchcolumn*
\selectlanguage{latin}
In Cypro natális beáti 
 Spiridiónis Epíscopi, qui unus fuit ex illis Confessóribus, quos Galérius 
 Maximiánus, dextro óculo effósso et sinístro póplite succíso, ad metálla damnáverat. Hic prophetíæ dono et signórum glória 
 ínclitus fuit, et in 
 Nicæno Concílio philósophum éthnicum, Christiánæ religióni insultántem, 
 devícit et ad fidem perdúxit.
\switchcolumn
\selectlanguage{english}
In the island of Cyprus, the 
 birthday of blessed Spiridion, bishop. He was one of those confessors 
 who were condemned by Galerius Maximian to labour in the mines, after 
 suffering the loss of his right eye and cutting of the sinews of his left 
 knee. This prelate was renowned for the gift of prophecy and glorious 
 miracles, and in the Council of Nicea he confounded a heathen philosopher, 
 who had insulted the Christian religion, and brought him to the faith.
\switchcolumn*
\selectlanguage{latin}
Bérgomi sancti Viatóris, 
 Epíscopi et Confessóris.
\switchcolumn
\selectlanguage{english}
At Bergamo, St. Viator, bishop and 
 confessor.
\switchcolumn*
\selectlanguage{latin}
Papíæ sancti Pompéji 
 Epíscopi.
\switchcolumn
\selectlanguage{english}
At Pavia, St. Pompey, bishop.
\switchcolumn*
\selectlanguage{latin}
Neápoli, in Campánia, 
 sancti Agnélli Abbátis, virtúte miraculórum illústris, qui obséssam urbem 
 sæpe visus est Crucis vexíllo ab hóstibus liberáre.
\switchcolumn
\selectlanguage{english}
At Naples in Campania, St. Agnellus, 
 abbot. Illustrious for the gift of miracles, he was often seen with 
 the standard of the Cross, delivering the city besieged by enemies.
\switchcolumn*
\selectlanguage{latin}
Medioláni sancti 
 Matroniáni Eremítæ.
\switchcolumn
\selectlanguage{english}
At Milan, St. Matronian, hermit.
\switchcolumn*
\selectlanguage{latin}
\end{paracol}


% ---- martyrology/mart12/mart1215.htm
\needspace{10\baselineskip}
\begin{paracol}{2}
\selectlanguage{latin}
\begin{center}{\color{gregoriocolor} Décimo octávo Kaléndas Januárii. 
 Luna\dots\ }\end{center}
\switchcolumn
\selectlanguage{english}
\begin{center}{\color{gregoriocolor} The 
 15\textsuperscript{th} Day of 
 December. The\dots\ Day of the Moon.}\end{center}
\end{paracol}

\noindent\begin{tabularx}{\linewidth}{*{19}{>{\centering\arraybackslash}X}}
 \textcolor{gregoriocolor}{a} & \textcolor{gregoriocolor}{b} & \textcolor{gregoriocolor}{c} & \textcolor{gregoriocolor}{d} & \textcolor{gregoriocolor}{e} & \textcolor{gregoriocolor}{f} & \textcolor{gregoriocolor}{g} & \textcolor{gregoriocolor}{h} & \textcolor{gregoriocolor}{i} & \textcolor{gregoriocolor}{k} & \textcolor{gregoriocolor}{l} & \textcolor{gregoriocolor}{m} & \textcolor{gregoriocolor}{n} & \textcolor{gregoriocolor}{p} & \textcolor{gregoriocolor}{q} & \textcolor{gregoriocolor}{r} & \textcolor{gregoriocolor}{s} & \textcolor{gregoriocolor}{t} & \textcolor{gregoriocolor}{u} \\
 25 & 26 & 27 & 28 & 29 & 1 & 2 & 3 & 4 & 5 & 6 & 7 & 8 & 9 & 10 & 11 & 12 & 13 & 14 \\
\end{tabularx}
\vspace{0.5\baselineskip}
\noindent\begin{tabularx}{\linewidth}{*{12}{>{\centering\arraybackslash}X}}
 \textcolor{gregoriocolor}{A} & \textcolor{gregoriocolor}{B} & \textcolor{gregoriocolor}{C} & \textcolor{gregoriocolor}{D} & \textcolor{gregoriocolor}{E} & F & \textcolor{gregoriocolor}{F} & \textcolor{gregoriocolor}{G} & \textcolor{gregoriocolor}{H} & \textcolor{gregoriocolor}{M} & \textcolor{gregoriocolor}{N} & \textcolor{gregoriocolor}{P} \\
 15 & 16 & 17 & 18 & 19 & 20 & 19 & 20 & 21 & 22 & 23 & 24 \\
\end{tabularx}

\begin{paracol}{2}
\selectlanguage{latin}
\lettrine[lines=2]{O}{ctáva} Conceptiónis 
 Immaculátæ beátæ Maríæ Vírginis.
\switchcolumn
\selectlanguage{english}
\lettrine[lines=2]{T}{he} Octave of the Immaculate 
 Conception of the Blessed Virgin Mary.
\switchcolumn*
\selectlanguage{latin}
Romæ sanctórum Mártyrum 
 Irenæi, Antónii, Theodóri, Saturníni, Victóris, et aliórum decem et septem, 
 qui, in persecutióne Valeriáni, pro Christo passi sunt.
\switchcolumn
\selectlanguage{english}
At Rome, the holy martyrs Irenaeus, 
 Anthony, Theodore, Saturninus, Victor, and seventeen others who suffered for 
 Christ in the persecution of Valerian.
\switchcolumn*
\selectlanguage{latin}
In Africa, pássio 
 sanctórum Faustíni, Lúcii, Cándidi, Cæliáni, Marci, Januárii et Fortunáti.
\switchcolumn
\selectlanguage{english}
In Africa, the martyrdom of Saints 
 Faustinus, Lucius, Candidus, Cælian, Mark, Januarius, and Fortunatus.
\switchcolumn*
\selectlanguage{latin}
Ibídem sancti Valeriáni 
 Epíscopi, qui, cum esset annórum plus octogínta, in persecutióne Wandálica, 
 sub Rege Ariáno Genseríco, convéntus ab eo ut tráderet Ecclésiæ utensília, 
 idque constánter renuísset, extra civitátem singuláris jussus est pelli; cumque præcéptum esset ut nullus eum neque in domo neque in agro dimítteret 
 habitáre, multo témpore in strata pública nudo sub áere jácuit, et, in 
 confessióne et defensióne cathólicæ veritátis, cursum beátæ vitæ complévit.
\switchcolumn
\selectlanguage{english}
In the same country, the holy bishop 
 Valerian, who, being upwards of eighty years of age, in the persecution of 
 the Vandals, under the Arian king Genseric, was asked to deliver the vessels 
 of the Church, and as he constantly refused, an order was issued to drive 
 him all alone out of the city, and all persons were forbidden to allow him 
 to stay in their houses or on their land. For a long time he remained 
 lying on the public road, in the open air, and thus in the confession and 
 defence of Catholic truth he ended his blessed life.
\switchcolumn*
\selectlanguage{latin}
In territorio 
 Aurelianénsi sancti Maximíni Confessóris.
\switchcolumn
\selectlanguage{english}
In the territory of Orleans, St. 
 Maximin, confessor.
\switchcolumn*
\selectlanguage{latin}
Apud Ibéros, trans 
 Pontum Euxínum, sanctæ Christiánæ ancíllæ, quæ miraculórum virtúte gentem 
 illam, témpore Constantíni, ad Christi fidem perdúxit.
\switchcolumn
\selectlanguage{english}
Among the Iberians across the Euxine 
 Sea, St. Christiana, a maidservant, who by virtue of her miracles led that 
 people to the faith of Christ, in the time of Constantine.
\switchcolumn*
\selectlanguage{latin}
Vercéllis Ordinátio 
 sancti Eusébii, Epíscopi et Mártyris.
\switchcolumn
\selectlanguage{english}
At Vercelli, the ordination of St. 
 Eusebius, bishop and martyr.
\switchcolumn*
\selectlanguage{latin}
\end{paracol}


% ---- martyrology/mart12/mart1216.htm
\needspace{10\baselineskip}
\begin{paracol}{2}
\selectlanguage{latin}
\begin{center}{\color{gregoriocolor} Décimo séptimo Kaléndas Januárii. 
 Luna\dots\ }\end{center}
\switchcolumn
\selectlanguage{english}
\begin{center}{\color{gregoriocolor} The 
 16\textsuperscript{th} Day of 
 December. The\dots\ Day of the Moon.}\end{center}
\end{paracol}

\noindent\begin{tabularx}{\linewidth}{*{19}{>{\centering\arraybackslash}X}}
 \textcolor{gregoriocolor}{a} & \textcolor{gregoriocolor}{b} & \textcolor{gregoriocolor}{c} & \textcolor{gregoriocolor}{d} & \textcolor{gregoriocolor}{e} & \textcolor{gregoriocolor}{f} & \textcolor{gregoriocolor}{g} & \textcolor{gregoriocolor}{h} & \textcolor{gregoriocolor}{i} & \textcolor{gregoriocolor}{k} & \textcolor{gregoriocolor}{l} & \textcolor{gregoriocolor}{m} & \textcolor{gregoriocolor}{n} & \textcolor{gregoriocolor}{p} & \textcolor{gregoriocolor}{q} & \textcolor{gregoriocolor}{r} & \textcolor{gregoriocolor}{s} & \textcolor{gregoriocolor}{t} & \textcolor{gregoriocolor}{u} \\
 26 & 27 & 28 & 29 & 1 & 2 & 3 & 4 & 5 & 6 & 7 & 8 & 9 & 10 & 11 & 12 & 13 & 14 & 15 \\
\end{tabularx}
\vspace{0.5\baselineskip}
\noindent\begin{tabularx}{\linewidth}{*{12}{>{\centering\arraybackslash}X}}
 \textcolor{gregoriocolor}{A} & \textcolor{gregoriocolor}{B} & \textcolor{gregoriocolor}{C} & \textcolor{gregoriocolor}{D} & \textcolor{gregoriocolor}{E} & F & \textcolor{gregoriocolor}{F} & \textcolor{gregoriocolor}{G} & \textcolor{gregoriocolor}{H} & \textcolor{gregoriocolor}{M} & \textcolor{gregoriocolor}{N} & \textcolor{gregoriocolor}{P} \\
 16 & 17 & 18 & 19 & 20 & 21 & 20 & 21 & 22 & 23 & 24 & 25 \\
\end{tabularx}

\begin{paracol}{2}
\selectlanguage{latin}
\lettrine[lines=2]{S}{ancti} Eusébii, Epíscopi Vercellénsis et Mártyris; cujus dies natális 
 Kaléndis Augústi, et Ordinátio décimo octávo Kaléndas Januárii refértur.
\switchcolumn
\selectlanguage{english}
\lettrine[lines=2]{S}{t.} Eusebius, bishop of Vercelli and martyr. His birthday is 
 commemorated on the 1st of August and his ordination on the 15th of 
 December.
\switchcolumn*
\selectlanguage{latin}
Sanctórum Trium Puerórum, id est Ananíæ, Azaríæ 
 et Misaélis; quorum córpora apud Babylóniam, sub quodam specu, sunt pósita.
\switchcolumn
\selectlanguage{english}
The three young men, Ananias, Azarias, and Misael, whose bodies are buried 
 in a cave near Babylon.
\switchcolumn*
\selectlanguage{latin}
Ravénnæ sanctórum Mártyrum Valentíni, magístri mílitum, ejúsque fílii
  Concórdii, atque Navális et Agrícolæ; qui, in persecutióne Maximiáni, pro
  Christo passi sunt.
\switchcolumn
\selectlanguage{english}
At 
 Ravenna, the holy martyrs Valentine, an officer of the army, Concordius, his 
 son, Navalis, and Agricola, who suffered for Christ in the persecution of 
 Maximian.
\switchcolumn*
\selectlanguage{latin}
Fórmiis, in Campánia, sanctæ Albínæ, Vírginis 
 et Mártyris, sub Décio Imperatóre.
\switchcolumn
\selectlanguage{english}
At 
 Mola di Gaeta in Campania, St. Albina, virgin and martyr, under Emperor 
 Decius.
\switchcolumn*
\selectlanguage{latin}
In 
 Africa pássio plurimárum sanctárum Vírginum, quæ, 
 in persecutióne Wandálica, sub Ariáno Rege Hunneríco, suspéndia, póndera 
 laminásque ignítas perpéssæ, martyrii agónem felíciter consummárunt.
\switchcolumn
\selectlanguage{english}
In 
 Africa, many holy virgins who reached a happy end of their martyrdom in the 
 persecution of the Vandals under the Arian king Hunneric by having heavy 
 weights tied to them and burning plates of metal applied to their bodies.
\switchcolumn*
\selectlanguage{latin}
Viénnæ, in Gállia, beáti Adónis, Epíscopi et 
 Confessóris.
\switchcolumn
\selectlanguage{english}
At 
 Vienne in France, blessed Ado, bishop and confessor.
\switchcolumn*
\selectlanguage{latin}
In 
 Hibérnia sancti Beáni Epíscopi.
\switchcolumn
\selectlanguage{english}
In 
 Ireland, St. Bean, bishop.
\switchcolumn*
\selectlanguage{latin}
Gazæ, in Palæstína, sancti Ireniónis Epíscopi.
\switchcolumn
\selectlanguage{english}
At Gaza in Palestine, St. Irenion, bishop.
\switchcolumn*
\selectlanguage{latin}
\end{paracol}


% ---- martyrology/mart12/mart1217.htm
\needspace{10\baselineskip}
\begin{paracol}{2}
\selectlanguage{latin}
\begin{center}{\color{gregoriocolor} Sextodécimo Kaléndas Januárii. 
 Luna\dots\ }\end{center}
\switchcolumn
\selectlanguage{english}
\begin{center}{\color{gregoriocolor} The 
 17\textsuperscript{th} Day of 
 December. The\dots\ Day of the Moon.}\end{center}
\end{paracol}

\noindent\begin{tabularx}{\linewidth}{*{19}{>{\centering\arraybackslash}X}}
 \textcolor{gregoriocolor}{a} & \textcolor{gregoriocolor}{b} & \textcolor{gregoriocolor}{c} & \textcolor{gregoriocolor}{d} & \textcolor{gregoriocolor}{e} & \textcolor{gregoriocolor}{f} & \textcolor{gregoriocolor}{g} & \textcolor{gregoriocolor}{h} & \textcolor{gregoriocolor}{i} & \textcolor{gregoriocolor}{k} & \textcolor{gregoriocolor}{l} & \textcolor{gregoriocolor}{m} & \textcolor{gregoriocolor}{n} & \textcolor{gregoriocolor}{p} & \textcolor{gregoriocolor}{q} & \textcolor{gregoriocolor}{r} & \textcolor{gregoriocolor}{s} & \textcolor{gregoriocolor}{t} & \textcolor{gregoriocolor}{u} \\
 27 & 28 & 29 & 1 & 2 & 3 & 4 & 5 & 6 & 7 & 8 & 9 & 10 & 11 & 12 & 13 & 14 & 15 & 16 \\
\end{tabularx}
\vspace{0.5\baselineskip}
\noindent\begin{tabularx}{\linewidth}{*{12}{>{\centering\arraybackslash}X}}
 \textcolor{gregoriocolor}{A} & \textcolor{gregoriocolor}{B} & \textcolor{gregoriocolor}{C} & \textcolor{gregoriocolor}{D} & \textcolor{gregoriocolor}{E} & F & \textcolor{gregoriocolor}{F} & \textcolor{gregoriocolor}{G} & \textcolor{gregoriocolor}{H} & \textcolor{gregoriocolor}{M} & \textcolor{gregoriocolor}{N} & \textcolor{gregoriocolor}{P} \\
 17 & 18 & 19 & 20 & 21 & 22 & 21 & 22 & 23 & 24 & 25 & 26 \\
\end{tabularx}

\begin{paracol}{2}
\selectlanguage{latin}
\lettrine[lines=2]{R}{omæ} natális sancti Joánnis de Matha, 
 Presbyteri et Confessóris, qui Ordinis sanctíssimæ Trinitátis redemptiónis 
 captivórum Fundátor éxstitit. Ipsíus tamen festívitas, ex dispositióne 
 Innocéntii Papæ Undécimi, ágitur sexto Idus Februárii.
\switchcolumn
\selectlanguage{english}
\lettrine[lines=2]{A}{t} Rome, the birthday of St. John of Matha, priest and confessor, founder of 
 the Order of the Most Holy Trinity for the Redemption of Captives, whose 
 feast, by decree of Pope Innocent XI, is observed on the 8th of 
 February.
\switchcolumn*
\selectlanguage{latin}
Massíliæ, in Gállia, beáti Lázari Epíscopi, 
 sanctárum Maríæ Magdalénæ ac Marthæ fratris, quem Dóminus in Evangélio 
 appellásse amícum et a mórtuis excitásse légitur.
\switchcolumn
\selectlanguage{english}
At Marseilles in France, blessed Lazarus, brother of the Saints Mary 
 Magdalene and Martha, of whom we read in the Gospel that our Lord called him 
 his friend and raised him from the dead.
\switchcolumn*
\selectlanguage{latin}
Eleutherópoli, in Palæstína, sanctórum Mártyrum 
 Floriáni, Calaníci, et Sociórum quinquagínta et octo; qui, témpore Heraclíi 
 Imperatóris, a Saracénis ob Christi fidem occísi sunt.
\switchcolumn
\selectlanguage{english}
At 
 Eleutheropolis, the holy martyrs Florian, Calanicus, and their fifty-eight 
 companions, who were slain by the Saracens because of the faith of Christ, 
 during the reign of Emperor Haraclius.
\switchcolumn*
\selectlanguage{latin}
In monastério Fuldénsi sancti Stúrmii, Abbátis et Saxóniæ 
 Apóstoli; quem Innocéntius Papa Secúndus, in Concílio secúndo Lateranénsi, 
 in Sanctórum númerum rétulit.
\switchcolumn
\selectlanguage{english}
In 
 the monastery of Fulda, the holy abbot Sturmius, apostle of Saxony, who was 
 ranked among the saints by Innocent II, in the second Lateran Council.
\switchcolumn*
\selectlanguage{latin}
Bigárdis, prope Bruxéllas, in Brabántia, sanctæ 
 Wivínæ Vírginis, cujus egrégiam sanctitátem mirácula crebra testántur.
\switchcolumn
\selectlanguage{english}
At 
 Bigarden, near Brussels, St. Wivina, virgin, whose eminent sanctity is 
 attested to by frequent miracles.
\switchcolumn*
\selectlanguage{latin}
Constantinópoli sanctæ Olympíadis Víduæ.
\switchcolumn
\selectlanguage{english}
At 
 Constantinople, St. Olympias, widow.
\switchcolumn*
\selectlanguage{latin}
Andániæ, apud Septem Ecclésias, in Bélgio, 
 beátæ Beggæ Víduæ, quæ fuit soror sanctæ Gertrúdis.
\switchcolumn
\selectlanguage{english}
At 
 Andenne, at the Seven Churches, blessed Begga, widow, the sister of St. 
 Gertrude.
\switchcolumn*
\selectlanguage{latin}
Eódem die Translátio sancti Ignátii, Epíscopi et Mártyris; qui, tértius post 
 beátum Petrum Apóstolum, Antiochénam rexit Ecclésiam. Ejus corpus ab 
 urbe Roma, ubi ipse, sub Trajáno, glorióse martyrium tertiodécimo Kaléndas 
 Januárii consummáverat, Antiochíam delátum, ibídem, in cœmetério 
 Ecclésiæ, extra portam Daphníticam, pósitum fuit; in qua celebritáte sanctus 
 Joánnes Chrysóstomus conciónem ad pópulum hábuit. Póstmodum vero ejus 
 relíquiæ rursus Romam translátæ sunt, et in Ecclésia sancti Cleméntis, una 
 cum córpore ejúsdem beatíssimi Papæ et Mártyris, summa veneratióne recónditæ.
\switchcolumn
\selectlanguage{english}
Also, the translation of St. Ignatius, bishop and martyr, who, the third 
 after the blessed Apostle Peter, governed the Church of Antioch. His 
 body was taken from Rome, where he had suffered martyrdom under Trajan on 
 the 20th of December, and deposited in the church cemetery near the Gate of 
 Daphne at Antioch. St. John Chrysostom, on that solemn occasion, 
 preached the sermon to the people. Afterwards his relics were carried 
 back to Rome and placed with the highest reverence in the church of St. 
 Clement, together with the body of that blessed pope and martyr.
\switchcolumn*
\selectlanguage{latin}
\end{paracol}


% ---- martyrology/mart12/mart1218.htm
\needspace{10\baselineskip}
\begin{paracol}{2}
\selectlanguage{latin}
\begin{center}{\color{gregoriocolor} Quintodécimo Kaléndas Januárii. 
 Luna\dots\ }\end{center}
\switchcolumn
\selectlanguage{english}
\begin{center}{\color{gregoriocolor} The 
 18\textsuperscript{th} Day of 
 December. The\dots\ Day of the Moon.}\end{center}
\end{paracol}

\noindent\begin{tabularx}{\linewidth}{*{19}{>{\centering\arraybackslash}X}}
 \textcolor{gregoriocolor}{a} & \textcolor{gregoriocolor}{b} & \textcolor{gregoriocolor}{c} & \textcolor{gregoriocolor}{d} & \textcolor{gregoriocolor}{e} & \textcolor{gregoriocolor}{f} & \textcolor{gregoriocolor}{g} & \textcolor{gregoriocolor}{h} & \textcolor{gregoriocolor}{i} & \textcolor{gregoriocolor}{k} & \textcolor{gregoriocolor}{l} & \textcolor{gregoriocolor}{m} & \textcolor{gregoriocolor}{n} & \textcolor{gregoriocolor}{p} & \textcolor{gregoriocolor}{q} & \textcolor{gregoriocolor}{r} & \textcolor{gregoriocolor}{s} & \textcolor{gregoriocolor}{t} & \textcolor{gregoriocolor}{u} \\
 28 & 29 & 1 & 2 & 3 & 4 & 5 & 6 & 7 & 8 & 9 & 10 & 11 & 12 & 13 & 14 & 15 & 16 & 17 \\
\end{tabularx}
\vspace{0.5\baselineskip}
\noindent\begin{tabularx}{\linewidth}{*{12}{>{\centering\arraybackslash}X}}
 \textcolor{gregoriocolor}{A} & \textcolor{gregoriocolor}{B} & \textcolor{gregoriocolor}{C} & \textcolor{gregoriocolor}{D} & \textcolor{gregoriocolor}{E} & F & \textcolor{gregoriocolor}{F} & \textcolor{gregoriocolor}{G} & \textcolor{gregoriocolor}{H} & \textcolor{gregoriocolor}{M} & \textcolor{gregoriocolor}{N} & \textcolor{gregoriocolor}{P} \\
 18 & 19 & 20 & 21 & 22 & 23 & 22 & 23 & 24 & 25 & 26 & 27 \\
\end{tabularx}

\begin{paracol}{2}
\selectlanguage{latin}
\lettrine[lines=2]{P}{hilíppis,} in Macedónia, natális sanctórum Mártyrum Rufi et Zósimi, 
 qui ex illórum número discipulórum fuérunt, per quos primitíva Ecclésia in 
 Judæis et Græcis fundáta est; de quorum étiam felíci agóne scribit sanctus 
 Polycárpus in epístola ad Philippénses.
\switchcolumn
\selectlanguage{english}
\lettrine[lines=2]{A}{t} Philippi in Macedonia, the birthday of the holy martyrs Rufus and Zosimus, 
 who were of the number of disciples by whom the primitive church was founded 
 among the Jews and the Greeks. Their happy martyrdom is mentioned by 
 St. Polycarp in his Epistle to the Philippians.
\switchcolumn*
\selectlanguage{latin}
Laodicéæ, in Syria, pássio sanctórum Theotími 
 et Basiliáni.
\switchcolumn
\selectlanguage{english}
At Laodicea in Syria, the martyrdom of the Saints Theotimus and Basilian.
\switchcolumn*
\selectlanguage{latin}
In Africa sanctórum Mártyrum Quincti, Simplícii et aliórum; qui sub Décii et 
 Valeriáni persecutióne passi sunt.
\switchcolumn
\selectlanguage{english}
In 
 Africa, the holy martyrs Quinctus, Simplicius, and others who suffered in 
 the persecution of Decius and Valerian.
\switchcolumn*
\selectlanguage{latin}
Ibídem sancti Moysétis Mártyris.
\switchcolumn
\selectlanguage{english}
In 
 the same country, St. Moses, martyr.
\switchcolumn*
\selectlanguage{latin}
Item in Africa sanctórum Mártyrum Victúri, Victóris, Victoríni, Adjutóris, Quarti et 
 aliórum trigínta.
\switchcolumn
\selectlanguage{english}
Also in Africa, the holy martyrs Victurus, Victor, Victorinus, Adjutor, 
 Quartus, and thirty others.
\switchcolumn*
\selectlanguage{latin}
Mopsuéstiæ, in Cilícia, sancti Auxéntii 
 Epíscopi, qui, olim sub Licínio miles, pótius elégit cíngulum exúere quam 
 uvas Baccho offérre; factúsque Epíscopus, præclárus méritis quiévit in pace.
\switchcolumn
\selectlanguage{english}
At 
 Mopsuestia in Cilicia, St. Auxentius, bishop, who, being at first a soldier 
 under Licinius, preferred to surrender his military insignia rather than 
 offer grapes to Bacchus. Having been made a bishop, he was renowned 
 for his merit, and died in peace.
\switchcolumn*
\selectlanguage{latin}
Turónis, in Gállia, sancti Gratiáni Epíscopi, qui, a sancto Fabiáno Papa 
 primus ejúsdem civitátis Epíscopus ordinátus est, et multis clarus miráculis 
 obdormívit in Dómino.
\switchcolumn
\selectlanguage{english}
At 
 Tours in France, St. Gratian, appointed first bishop of that city by Pope 
 St. Fabian. Celebrated for many miracles, he calmly went to his repose 
 in the Lord.
\switchcolumn*
\selectlanguage{latin}
\end{paracol}


% ---- martyrology/mart12/mart1219.htm
\needspace{10\baselineskip}
\begin{paracol}{2}
\selectlanguage{latin}
\begin{center}{\color{gregoriocolor} Quartodécimo Kaléndas Januárii. 
 Luna\dots\ }\end{center}
\switchcolumn
\selectlanguage{english}
	\begin{center}{\color{gregoriocolor} The 19\textsuperscript{th} Day of 
 December. The\dots\ Day of the Moon.}\end{center}
\end{paracol}

\noindent\begin{tabularx}{\linewidth}{*{19}{>{\centering\arraybackslash}X}}
 \textcolor{gregoriocolor}{a} & \textcolor{gregoriocolor}{b} & \textcolor{gregoriocolor}{c} & \textcolor{gregoriocolor}{d} & \textcolor{gregoriocolor}{e} & \textcolor{gregoriocolor}{f} & \textcolor{gregoriocolor}{g} & \textcolor{gregoriocolor}{h} & \textcolor{gregoriocolor}{i} & \textcolor{gregoriocolor}{k} & \textcolor{gregoriocolor}{l} & \textcolor{gregoriocolor}{m} & \textcolor{gregoriocolor}{n} & \textcolor{gregoriocolor}{p} & \textcolor{gregoriocolor}{q} & \textcolor{gregoriocolor}{r} & \textcolor{gregoriocolor}{s} & \textcolor{gregoriocolor}{t} & \textcolor{gregoriocolor}{u} \\
 29 & 1 & 2 & 3 & 4 & 5 & 6 & 7 & 8 & 9 & 10 & 11 & 12 & 13 & 14 & 15 & 16 & 17 & 18 \\
\end{tabularx}
\vspace{0.5\baselineskip}
\noindent\begin{tabularx}{\linewidth}{*{12}{>{\centering\arraybackslash}X}}
 \textcolor{gregoriocolor}{A} & \textcolor{gregoriocolor}{B} & \textcolor{gregoriocolor}{C} & \textcolor{gregoriocolor}{D} & \textcolor{gregoriocolor}{E} & F & \textcolor{gregoriocolor}{F} & \textcolor{gregoriocolor}{G} & \textcolor{gregoriocolor}{H} & \textcolor{gregoriocolor}{M} & \textcolor{gregoriocolor}{N} & \textcolor{gregoriocolor}{P} \\
 19 & 20 & 21 & 22 & 23 & 24 & 23 & 24 & 25 & 26 & 27 & 28 \\
\end{tabularx}

\begin{paracol}{2}
\selectlanguage{latin}
\lettrine[lines=2]{I}{n} Mauritánia sancti Timóthei Diáconi, qui, ob Christi fidem, post diros cárceres, in ignem conjéctus, martyrium consummávit.
\switchcolumn
\selectlanguage{english}
\lettrine[lines=2]{I}{n} Morocco, St. Timothy, deacon, who after severe imprisonment for the sake 
 of Christ was cast into the fire and achieved martyrdom.
\switchcolumn*
\selectlanguage{latin}
Alexandríæ beáti Nemésii Mártyris, qui, primo 
 per calúmniam quasi latro Júdici delátus, eóque crímine absolútus, mox, in 
 persecutióne Décii, Christiánæ religiónis nómine accusátus, cum latrónibus 
 jussus est incéndi, Salvatóris déferens similitúdinem, qui una cum latrónibus 
 pértulit crucem.
\switchcolumn
\selectlanguage{english}
At Alexandria in Egypt, blessed Nemesius, martyr, who first was denounced 
 before the judge as a robber, and being freed from that charge, soon after, 
 in the persecution of Decius, was accused before the judge Emilian of being 
 a Christian. He was twice subjected to torture and condemned to be 
 burned alive with robbers, thus bearing a resemblance to our Saviour, who 
 was crucified with thieves.
\switchcolumn*
\selectlanguage{latin}
Nicææ, in Bithynia, sanctórum Mártyrum Daríi, 
 Zósimi, Pauli et Secúndi.
\switchcolumn
\selectlanguage{english}
At Nicaea, the Saints Darius, Zosimus, Paul, and Secundus, martyrs.
\switchcolumn*
\selectlanguage{latin}
Nicomedíæ sanctórum Mártyrum Cyríaci, Paulílli, 
 Secúndi, Anastásii, Syndímii et Sociórum.
\switchcolumn
\selectlanguage{english}
At Nicomedia, the holy martyrs Cyriac, Paulillus, Secundus, Anastasius, 
 Sindimius, and their companions.
\switchcolumn*
\selectlanguage{latin}
Gazæ, in Palæstína, pássio sanctárum Meuris et 
 Theæ.
\switchcolumn
\selectlanguage{english}
At Gaza in Palestine, the martyrdom of Saints Meuris and Thea.
\switchcolumn*
\selectlanguage{latin}
Romæ deposítio sancti Anastásii Papæ Primi, viri ditíssimæ paupertátis et
  apostólicæ sollicitúdinis, quem (ut ait sanctus Hierónymus) diu Roma habére
  non méruit, ne orbis caput sub tali Epíscopo truncarétur; nam, haud multo
  post ejus óbitum, Roma a Gothis capta et dirépta fuit.
\switchcolumn
\selectlanguage{english}
At Rome, the death of Pope St. Anastasius I, a man who was rich in his 
 poverty and filled with apostolic zeal. St. Jerome says that Rome did 
 not deserve to possess him long, lest the capital of the world should be 
 devastated under so fine a bishop, for shortly after his death Rome was 
 taken and sacked by the Goths.
\switchcolumn*
\selectlanguage{latin}
Antisiodóri sancti Gregórii, Epíscopi et Confessóris.
\switchcolumn
\selectlanguage{english}
At Auxerre, St. Gregory, bishop and confessor.
\switchcolumn*
\selectlanguage{latin}
Aureliánis, in Gállia, sancti Adjúti Abbátis, prophético spíritu illústris.
\switchcolumn
\selectlanguage{english}
At Orleans in France, St. Adjutus, abbot, famous for the spirit of prophecy.
\switchcolumn*
\selectlanguage{latin}
Romæ sanctæ Faustæ, quæ fuit mater sanctæ 
 Anastásiæ, ac nobilitáte et pietáte éxstitit insígnis.
\switchcolumn
\selectlanguage{english}
At Rome, St. Fausta, mother of St. Anastasia, renowned for her noble birth 
 and her holiness.
\switchcolumn*
\selectlanguage{latin}
Avenióne beáti Urbáni Papæ Quinti, qui, Sede 
 Apostólica Romæ restitúta, Græcórum cum Latínis conjunctióne perfécta, 
 infidélibus coércitis, de Ecclésia óptime méritus est. Ejus cultum 
 pervetústum Pius Nonus, Póntifex Máximus, ratum hábuit et confirmávit.
\switchcolumn
\selectlanguage{english}
At 
 Avignon, blessed Urban V, who deserved well of the Church by restoring the 
 Apostolic See to Rome, by bringing about a reunion of the Latins and the 
 Greeks, and by suppressing heretics. Pius IX approved and confirmed 
 the veneration which had long been paid to him.
\switchcolumn*
\selectlanguage{latin}
\end{paracol}


% ---- martyrology/mart12/mart1220.htm
\needspace{10\baselineskip}
\begin{paracol}{2}
\selectlanguage{latin}
\begin{center}{\color{gregoriocolor} Tertiodécimo Kaléndas Januárii. 
 Luna\dots\ }\end{center}
\switchcolumn
\selectlanguage{english}
\begin{center}{\color{gregoriocolor} The 
 20\textsuperscript{th} Day of 
 December. The\dots\ Day of the Moon.}\end{center}
\end{paracol}

\noindent\begin{tabularx}{\linewidth}{*{19}{>{\centering\arraybackslash}X}}
 \textcolor{gregoriocolor}{a} & \textcolor{gregoriocolor}{b} & \textcolor{gregoriocolor}{c} & \textcolor{gregoriocolor}{d} & \textcolor{gregoriocolor}{e} & \textcolor{gregoriocolor}{f} & \textcolor{gregoriocolor}{g} & \textcolor{gregoriocolor}{h} & \textcolor{gregoriocolor}{i} & \textcolor{gregoriocolor}{k} & \textcolor{gregoriocolor}{l} & \textcolor{gregoriocolor}{m} & \textcolor{gregoriocolor}{n} & \textcolor{gregoriocolor}{p} & \textcolor{gregoriocolor}{q} & \textcolor{gregoriocolor}{r} & \textcolor{gregoriocolor}{s} & \textcolor{gregoriocolor}{t} & \textcolor{gregoriocolor}{u} \\
 1 & 2 & 3 & 4 & 5 & 6 & 7 & 8 & 9 & 10 & 11 & 12 & 13 & 14 & 15 & 16 & 17 & 18 & 19 \\
\end{tabularx}
\vspace{0.5\baselineskip}
\noindent\begin{tabularx}{\linewidth}{*{12}{>{\centering\arraybackslash}X}}
 \textcolor{gregoriocolor}{A} & \textcolor{gregoriocolor}{B} & \textcolor{gregoriocolor}{C} & \textcolor{gregoriocolor}{D} & \textcolor{gregoriocolor}{E} & F & \textcolor{gregoriocolor}{F} & \textcolor{gregoriocolor}{G} & \textcolor{gregoriocolor}{H} & \textcolor{gregoriocolor}{M} & \textcolor{gregoriocolor}{N} & \textcolor{gregoriocolor}{P} \\
 20 & 21 & 22 & 23 & 24 & 25 & 24 & 25 & 26 & 27 & 28 & 29 \\
\end{tabularx}

\begin{paracol}{2}
\selectlanguage{latin}
\lettrine[lines=1]{V}{igília} sancti Thomæ Apóstoli.
\switchcolumn
\selectlanguage{english}
\lettrine[lines=1]{T}{he} Vigil of St. Thomas, Apostle.
\switchcolumn*
\selectlanguage{latin}
Romæ natális sancti Zephyríni, Papæ et Mártyris. 
 Ipsíus tamen festum recólitur séptimo Kaléndas Septémbris.
\switchcolumn
\selectlanguage{english}
At Rome, the birthday of St. Zephyrinus, pope and martyr. His feast is 
 celebrated on the 26th of August.
\switchcolumn*
\selectlanguage{latin}
Ibídem pássio sancti Ignátii, Epíscopi et Mártyris; qui, tértius post beátum 
 Petrum Apóstolum, Antiochénam rexit Ecclésiam. Hic, in persecutióne 
 Trajáni, damnátus ad béstias, Romam vinctus míttitur; ibíque, circumsedénte 
 Senátu, immaníssimis pœnárum supplíciis primo 
 est afféctus, dehinc objícitur leónibus, quorum déntibus præfocátus, hóstia 
 Christi effícitur. Ejus vero festívitas Kaléndis Februárii celebrátur.
\switchcolumn
\selectlanguage{english}
In the same city, the martyrdom of St. Ignatius, bishop and martyr. He 
 was the third after St. Peter the Apostle to rule the church of Antioch, and 
 in the persecution of Trajan was condemned to the beasts. By order of 
 Trajan he was sent to Rome in fetters, and there tortured and afflicted with 
 the most cruel torments in the midst of the assembled Senate. Finally 
 he was cast to the lions, and being ground by their teeth became a sacrifice 
 for Christ. His feast is observed on the 1st of February.
\switchcolumn*
\selectlanguage{latin}
Item Romæ sanctórum Mártyrum Liberáti et Bájuli.
\switchcolumn
\selectlanguage{english}
At Rome, the holy martyrs Liberatus and Bajulus.
\switchcolumn*
\selectlanguage{latin}
In Arábia sanctórum Mártyrum Eugénii et Macárii Presbyterórum, qui a Juliáno Apóstata, cum ipsíus impietátem arguíssent, sævíssimis 
 plagis affécti sunt, atque in vastíssimam erémum relegáti, et gládio cæsi.
\switchcolumn
\selectlanguage{english}
In Arabia, the holy martyrs Eugene and Macarius, priests. For 
 reproving Julian the Apostate for his impiety, they received severe stripes, 
 were banished to a vast desert, and finally were put to the sword.
\switchcolumn*
\selectlanguage{latin}
Alexandríæ sanctórum mílitum et Mártyrum Ammónis, Zenónis, Ptolemæi, Ingenis et
  Theóphili; qui, tribunálibus astántes, cum quidam Christiánus, in supplíciis
  pósitus, trepidáret et jam prope ad negándum declináret, vultu, óculis ac
  nútibus illum conabántur erígere. Cumque hac de causa clamor totíus pópuli in
  eos prosilíret, prorumpéntes in médium se Christiános esse testáti sunt; per
  quorum victóriam Christus, qui suis eam ánimi constántiam déderat,
  gloriosíssime triumphávit.
\switchcolumn
\selectlanguage{english}
At Alexandria, the holy martyrs Ammon, Zeno, Ptolemy, Ingen, and Theophilus, 
 soldiers. Standing near the tribunals, and seeing a Christian under 
 torture and almost ready to apostatize, they endeavoured to encourage him by 
 their looks and by signs. When on account of this the crowd raised an 
 outcry against them, they stepped forward and declared themselves 
 Christians. In their victory, Christ also who had given them fortitude 
 triumphed.
\switchcolumn*
\selectlanguage{latin}
Géldubæ, in 
 Germánia, sancti Júlii Mártyris.
\switchcolumn
\selectlanguage{english}
At Gelduba in Germany, St. Julius, martyr.
\switchcolumn*
\selectlanguage{latin}
Antiochíæ natális sancti Philogónii Epíscopi, 
 qui, Dei nutu ex causídico ad eam Ecclésiam regéndam accersítus, advérsus 
 Aríum, una cum sancto Alexándro Epíscopo et Sóciis, primum pro fide 
 cathólica certámen íniit, clarúsque méritis quiévit in Dómino; cujus ánnuam 
 festivitátem sanctus Joánnes Chrysóstomus præcláro encómio celebrávit.
\switchcolumn
\selectlanguage{english}
At Antioch, the birthday of St. Philogonius, bishop, who was called by the 
 will of God from the office of lawyer to the government of that church. 
 With the saintly bishop Alexander and his companions, he engaged in the 
 first contest for the Catholic faith against Arius. Renowned for 
 merits he rested in the Lord, and his feast was commemorated by St. John 
 Chrysostom with an excellent eulogy.
\switchcolumn*
\selectlanguage{latin}
Bríxiæ sancti Domínici, Epíscopi et Confessóris.
\switchcolumn
\selectlanguage{english}
At Brescia, St. Dominic, bishop and confessor.
\switchcolumn*
\selectlanguage{latin}
In Hispánia deposítio sancti Domínici de Silos Abbátis, e sancti Benedícti 
 Ordine, miráculis in captivórum liberatióne celebérrimi.
\switchcolumn
\selectlanguage{english}
In 
 Spain, the death of St. Dominic of Silos, abbot of the Order of St. 
 Benedict, renowned for the miracles which he had wrought for the liberation 
 of captives.
\switchcolumn*
\selectlanguage{latin}
\end{paracol}


% ---- martyrology/mart12/mart1221.htm
\needspace{10\baselineskip}
\begin{paracol}{2}
\selectlanguage{latin}
\begin{center}{\color{gregoriocolor} Duodécimo Kaléndas Januárii. 
 Luna\dots\ }\end{center}
\switchcolumn
\selectlanguage{english}
\begin{center}{\color{gregoriocolor} The 
 21\textsuperscript{st} Day of 
 December. The\dots\ Day of the Moon.}\end{center}
\end{paracol}

\noindent\begin{tabularx}{\linewidth}{*{19}{>{\centering\arraybackslash}X}}
 \textcolor{gregoriocolor}{a} & \textcolor{gregoriocolor}{b} & \textcolor{gregoriocolor}{c} & \textcolor{gregoriocolor}{d} & \textcolor{gregoriocolor}{e} & \textcolor{gregoriocolor}{f} & \textcolor{gregoriocolor}{g} & \textcolor{gregoriocolor}{h} & \textcolor{gregoriocolor}{i} & \textcolor{gregoriocolor}{k} & \textcolor{gregoriocolor}{l} & \textcolor{gregoriocolor}{m} & \textcolor{gregoriocolor}{n} & \textcolor{gregoriocolor}{p} & \textcolor{gregoriocolor}{q} & \textcolor{gregoriocolor}{r} & \textcolor{gregoriocolor}{s} & \textcolor{gregoriocolor}{t} & \textcolor{gregoriocolor}{u} \\
 2 & 3 & 4 & 5 & 6 & 7 & 8 & 9 & 10 & 11 & 12 & 13 & 14 & 15 & 16 & 17 & 18 & 19 & 20 \\
\end{tabularx}
\vspace{0.5\baselineskip}
\noindent\begin{tabularx}{\linewidth}{*{12}{>{\centering\arraybackslash}X}}
 \textcolor{gregoriocolor}{A} & \textcolor{gregoriocolor}{B} & \textcolor{gregoriocolor}{C} & \textcolor{gregoriocolor}{D} & \textcolor{gregoriocolor}{E} & F & \textcolor{gregoriocolor}{F} & \textcolor{gregoriocolor}{G} & \textcolor{gregoriocolor}{H} & \textcolor{gregoriocolor}{M} & \textcolor{gregoriocolor}{N} & \textcolor{gregoriocolor}{P} \\
 21 & 22 & 23 & 24 & 25 & 26 & 25 & 26 & 27 & 28 & 29 & 1 \\
\end{tabularx}

\begin{paracol}{2}
\selectlanguage{latin}
\lettrine[lines=2]{C}{alamínæ} natális beáti Thomæ Apóstoli, qui 
 Parthis, Medis, Persis et Hyrcánis Evangélium prædicávit; ac demum in Indiam 
 pervénit, ibíque, cum eos pópulos in Christiána religióne instituísset, 
 Regis jussu lánceis transfíxus occúbuit. Ipsíus relíquiæ primo ad 
 urbem Edéssam, in Mesopotámia, deínde Ortónam, apud Frentános, translátæ 
 sunt.
\switchcolumn
\selectlanguage{english}
\lettrine[lines=2]{A}{t} Mylapore, the birthday of the blessed Apostle Thomas, who preached the 
 Gospel to the Parthians, Medes, Persians, and Hyrcanians. Having 
 finally penetrated into India, and instructed those nations in the Christian 
 religion, he died pierced with lances at the order of the king. His 
 remains were first taken to the city of Edessa in Mesopotamia, and then to 
 Ortona.
\switchcolumn*
\selectlanguage{latin}
Fribúrgi Helvetiórum item natális sancti Petri Canísii, Sacerdótis e 
 Societáte Jesu et Confessóris, doctrína et sanctitáte præclári; 
 qui, difficíllimis Germániæ tempóribus, fidem cathólicam strénue deféndit ac 
 propagávit. Eum vero Pius Undécimus, Póntifex Máximus, Sanctórum 
 catálogo adscrípsit, simúlque Doctórem universális Ecclésiæ declarávit, et 
 ipsíus festum quinto Kaléndas Maji agéndum esse decrévit.
\switchcolumn
\selectlanguage{english}
At Fribourg in Switzerland, the birthday also of St. Peter Canisius, priest 
 of the Society of Jesus, a confessor famed for his sanctity and learning. 
 He defended and spread the Catholic faith with the utmost zeal in Germany 
 during its most difficult times. Pope Pius XI added him to the list of 
 the saints, and at the same time declared him to be a doctor of the 
 universal Church, appointing his feast to be observed on the 27th of April.
\switchcolumn*
\selectlanguage{latin}
Antiochíæ sancti Anastásii, Epíscopi et 
 Mártyris; qui, Phocæ Imperatóris témpore, a Judæis, in seditióne ab ipsis 
 contra Christiános facta, sævíssime necátus est.
\switchcolumn
\selectlanguage{english}
At Antioch, St. Anastasius, bishop and martyr. During the reign of 
 Emperor Phocas he was cruelly murdered by Jews in a riot which they had 
 instigated against the Christians.
\switchcolumn*
\selectlanguage{latin}
Nicomedíæ sancti Glycérii Presbyteri, qui, in 
 persecutióne Diocletiáni, multis torméntis vexátus, demum, in ignem 
 conjéctus, martyrium consummávit.
\switchcolumn
\selectlanguage{english}
At Nicomedia, St. Glycerius, priest. During the persecution of 
 Diocletian he was subjected to many torments, and finally fulfilled his 
 martyrdom by being cast into the flames.
\switchcolumn*
\selectlanguage{latin}
In Túscia sanctórum Mártyrum Joánnis et Festi.
\switchcolumn
\selectlanguage{english}
In Tuscany, the holy martyrs John and Festus.
\switchcolumn*
\selectlanguage{latin}
In Lycia sancti Themístoclis Mártyris, qui, sub 
 Décio Imperatóre, pro sancto Dióscoro, qui quærebátur ad necem, se óbtulit, 
 et, equúleo tortus, raptátus ac fústibus cæsus, martyrii corónam adéptus 
 est.
\switchcolumn
\selectlanguage{english}
In Lycia, St. Themistocles, martyr. In the reign of Emperor Decius, he 
 offered himself to take the place of Dioscorus, whom they were seeking to 
 slay. He was tortured on the rack, dragged over rough ways and 
 scourged, and thus obtained the crown of martyrdom.
\switchcolumn*
\selectlanguage{latin}
Tréviris sancti Severíni, Epíscopi et 
 Confessóris.
\switchcolumn
\selectlanguage{english}
At Treves, St. Severinus, bishop and confessor.
\switchcolumn*
\selectlanguage{latin}
\end{paracol}


% ---- martyrology/mart12/mart1222.htm
\needspace{10\baselineskip}
\begin{paracol}{2}
\selectlanguage{latin}
\begin{center}{\color{gregoriocolor} Undécimo Kaléndas Januárii. 
 Luna\dots\ }\end{center}
\switchcolumn
\selectlanguage{english}
\begin{center}{\color{gregoriocolor} The 22\textsuperscript{nd} Day of December. The\dots\ Day of the Moon.}\end{center}
\end{paracol}

\noindent\begin{tabularx}{\linewidth}{*{19}{>{\centering\arraybackslash}X}}
 \textcolor{gregoriocolor}{a} & \textcolor{gregoriocolor}{b} & \textcolor{gregoriocolor}{c} & \textcolor{gregoriocolor}{d} & \textcolor{gregoriocolor}{e} & \textcolor{gregoriocolor}{f} & \textcolor{gregoriocolor}{g} & \textcolor{gregoriocolor}{h} & \textcolor{gregoriocolor}{i} & \textcolor{gregoriocolor}{k} & \textcolor{gregoriocolor}{l} & \textcolor{gregoriocolor}{m} & \textcolor{gregoriocolor}{n} & \textcolor{gregoriocolor}{p} & \textcolor{gregoriocolor}{q} & \textcolor{gregoriocolor}{r} & \textcolor{gregoriocolor}{s} & \textcolor{gregoriocolor}{t} & \textcolor{gregoriocolor}{u} \\
 3 & 4 & 5 & 6 & 7 & 8 & 9 & 10 & 11 & 12 & 13 & 14 & 15 & 16 & 17 & 18 & 19 & 20 & 21 \\
\end{tabularx}
\vspace{0.5\baselineskip}
\noindent\begin{tabularx}{\linewidth}{*{12}{>{\centering\arraybackslash}X}}
 \textcolor{gregoriocolor}{A} & \textcolor{gregoriocolor}{B} & \textcolor{gregoriocolor}{C} & \textcolor{gregoriocolor}{D} & \textcolor{gregoriocolor}{E} & F & \textcolor{gregoriocolor}{F} & \textcolor{gregoriocolor}{G} & \textcolor{gregoriocolor}{H} & \textcolor{gregoriocolor}{M} & \textcolor{gregoriocolor}{N} & \textcolor{gregoriocolor}{P} \\
 22 & 23 & 24 & 25 & 26 & 27 & 26 & 27 & 28 & 29 & 1 & 2 \\
\end{tabularx}

\begin{paracol}{2}

\begin{american}

\selectlanguage{latin}
\lettrine[lines=2]{C}{hicágiæ,} Sanctæ Francíscæ Xavériæ
Cabríni, Vírginis, Institúti Missionariárum a Sacratíssimo Corde Jesu
fundatrícis, exímia caritáte, invícta ánimi fortitúdine et humilitáte
insígnis, quamn Pius Papa Duodécimus sanctárum catálogo adscrípsit.
\switchcolumn \selectlanguage{english}
% TODO: find the enlish version...
\switchcolumn* 

\selectlanguage{latin}
Romæ, via Lavicána, inter duas Lauros, natális 
sanctórum trigínta Mártyrum, qui omnes una die, in persecutióne Diocletiáni, 
martyrio coronáti sunt.
\switchcolumn
\selectlanguage{english}
At Rome, on the Lavican Way, between the two laurels, the birthday of thirty 
holy martyrs who were all crowned with martyrdom on the one day in the 
persecution of Diocletian.

\end{american}

\begin{roman}
	\selectlanguage{latin}
	\lettrine[lines=2]{R}{omæ,} via Lavicána, inter duas Lauros, natális 
	 sanctórum trigínta Mártyrum, qui omnes una die, in persecutióne Diocletiáni, 
	 martyrio coronáti sunt.
	\switchcolumn
	\selectlanguage{english}
	\lettrine[lines=2]{A}{t} Rome, on the Lavican Way, between the two laurels, the birthday of thirty 
	 holy martyrs who were all crowned with martyrdom on the one day in the 
	 persecution of Diocletian.
\end{roman}

\switchcolumn*
\selectlanguage{latin}
Item Romæ sancti Flaviáni Expræfécti, viri 
 beátæ Mártyris Dafrósæ atque patris beatárum Vírginum et Mártyrum Bibiánæ ac 
 Demétriæ; qui, sub Juliáno Apóstata, pro Christo inscriptióne damnátus, et 
 ad Aquas Taurínas, in Etrúria, in exsílium missus, illic in oratióne 
 spíritum Deo réddidit.
\switchcolumn
\selectlanguage{english}
In the same city, St. Flavian, an ex-prefect, the husband of the blessed 
 martyr Dafrosa, and the father of the holy virgin martyrs, Bibiana and 
 Demetria. He was condemned under Julian the Apostate to be branded for 
 Christ, and was exiled to Aquæ Taurinæ, where he gave up his soul to God 
 in prayer.
\switchcolumn*
\selectlanguage{latin}
In Ægypto sanctórum Chærémonis, Epíscopi 
 Nilópolis, et aliórum plurimórum Mártyrum. Horum álii, sæviénte Décii 
 persecutióne, fuga dispérsi, in solitúdinis errántes, a béstiis interémpti 
 sunt; álii fame, frígore ac languóre consúmpti; álii a bárbaris et 
 latrónibus necáti; atque ita omnes, divérso mortis génere, eádem martyrii 
 glória coronáti sunt.
\switchcolumn
\selectlanguage{english}
In Egypt, St. Chaeremon, bishop of Nilopolis, and many other martyrs. 
 While the persecution of Decius was raging, some of them were dispersed in 
 flight, and wandering through deserts were killed by wild beasts; others 
 perished by famine, cold, and sickness; others again were murdered by 
 barbarians and robbers, and thus all were crowned with a glorious martyrdom.
\switchcolumn*
\selectlanguage{latin}
Apud Ostia Tiberína sanctórum Mártyrum Demétrii, Honoráti et Flori.
\switchcolumn
\selectlanguage{english}
At Ostia, the holy martyrs Demetrius, Honoratus, and Florus.
\switchcolumn*
\selectlanguage{latin}
Alexandríæ sancti Ischyriónis Mártyris, qui, 
 cum ad sacrificándum convíciis et injúriis cogerétur atque contémneret, ídeo, 
 præacúta sude per média víscera transverberátus, neci tráditur.
\switchcolumn
\selectlanguage{english}
At Alexandria, St. Ischyrion, martyr. Because he despised all the 
 injuries he was made to suffer in attempts to force him to sacrifice to 
 idols, his bowels were pierced with a sharp stake, bringing his death.
\switchcolumn*
\selectlanguage{latin}
Nicomedíæ sancti Zenónis mílitis, qui, cum 
 Diocletiánum Céreri immolántem derisísset, proptérea, maxíllis confráctis 
 dentibúsque excússis, cápite truncátus est.
\switchcolumn
\selectlanguage{english}
At Nicomedia, St. Zeno, a soldier who mocked Diocletian for sacrificing to 
 Ceres, wherefore his jawbones were broken, his teeth knocked out, and his 
 head struck off.
\switchcolumn*
\selectlanguage{latin}
Chicágiæ 
 sanctæ Francíscæ Xavériæ Cabríni, Vírginis, Institúti Missionariárum a 
 Sacratíssimo Corde Jesu Fundatrícis, exímia caritáte, invícta ánimi 
 fortitúdine et humilitáte insígnis, quam Pius Papa Duodécimus, Sanctárum 
 catálogo adscrípsit, et ómnium emigrántium cæléstem apud Deum Patrónam 
 constítuit.
\switchcolumn
\selectlanguage{english}
At Chicago, St. Frances Xavier Cabrini, virgin, foundress of the 
 Congregation of Missionaries of the Sacred Heart of Jesus, distinguished for 
 charity, humility, and invincible fortitude. Pope Pius XII added her to the 
 catalogue of saints, and named her as the heavenly patroness of all 
 emigrants.
\switchcolumn*
\selectlanguage{latin}
\end{paracol}


% ---- martyrology/mart12/mart1223.htm
\needspace{10\baselineskip}
\begin{paracol}{2}
\selectlanguage{latin}
\begin{center}{\color{gregoriocolor} Décimo Kaléndas Januárii. 
 Luna\dots\ }\end{center}
\switchcolumn
\selectlanguage{english}
\begin{center}{\color{gregoriocolor} The 
 23\textsuperscript{rd} Day of 
 December. The\dots\ Day of the Moon.}\end{center}
\end{paracol}

\noindent\begin{tabularx}{\linewidth}{*{19}{>{\centering\arraybackslash}X}}
 \textcolor{gregoriocolor}{a} & \textcolor{gregoriocolor}{b} & \textcolor{gregoriocolor}{c} & \textcolor{gregoriocolor}{d} & \textcolor{gregoriocolor}{e} & \textcolor{gregoriocolor}{f} & \textcolor{gregoriocolor}{g} & \textcolor{gregoriocolor}{h} & \textcolor{gregoriocolor}{i} & \textcolor{gregoriocolor}{k} & \textcolor{gregoriocolor}{l} & \textcolor{gregoriocolor}{m} & \textcolor{gregoriocolor}{n} & \textcolor{gregoriocolor}{p} & \textcolor{gregoriocolor}{q} & \textcolor{gregoriocolor}{r} & \textcolor{gregoriocolor}{s} & \textcolor{gregoriocolor}{t} & \textcolor{gregoriocolor}{u} \\
 4 & 5 & 6 & 7 & 8 & 9 & 10 & 11 & 12 & 13 & 14 & 15 & 16 & 17 & 18 & 19 & 20 & 21 & 22 \\
\end{tabularx}
\vspace{0.5\baselineskip}
\noindent\begin{tabularx}{\linewidth}{*{12}{>{\centering\arraybackslash}X}}
 \textcolor{gregoriocolor}{A} & \textcolor{gregoriocolor}{B} & \textcolor{gregoriocolor}{C} & \textcolor{gregoriocolor}{D} & \textcolor{gregoriocolor}{E} & F & \textcolor{gregoriocolor}{F} & \textcolor{gregoriocolor}{G} & \textcolor{gregoriocolor}{H} & \textcolor{gregoriocolor}{M} & \textcolor{gregoriocolor}{N} & \textcolor{gregoriocolor}{P} \\
 23 & 24 & 25 & 26 & 27 & 28 & 27 & 28 & 29 & 1 & 2 & 3 \\
\end{tabularx}

\begin{paracol}{2}
\selectlanguage{latin}
\lettrine[lines=2]{R}{omæ} sanctæ Victóriæ, Vírginis et Mártyris, quæ, 
 in persecutióne Décii Imperatóris, cum esset desponsáta Eugénio pagáno et 
 nec núbere vellet neque sacrificáre, ídeo, post multa facta mirácula, quibus 
 plúrimas Deo Vírgines aggregáverat, a carnífice percússa est gládio in corde, 
 rogátu sui sponsi.
\switchcolumn
\selectlanguage{english}
\lettrine[lines=2]{A}{t} Rome, St. Victoria, virgin and martyr, during the persecution of Emperor 
 Decius. She had been promised in marriage to a pagan named Eugene, but 
 because she had refused to marry him and to offer sacrifice to idols, and 
 because by working many miracles she had brought many virgins to the service 
 of God, the executioner thrust a sword into her heart at the request of her 
 spouse.
\switchcolumn*
\selectlanguage{latin}
Nicomedíæ pássio sanctórum Migdónii et Mardónii, 
 quorum alter, in Diocletiáni persecutióne, igne cremátus, alter in fossam 
 projéctus occúbuit. Tunc étiam Diáconus sancti Anthimi, Epíscopi 
 Nicomediénsis, passus est; qui, cum lítteras perférret ad Mártyres, a 
 Gentílibus tentus est, atque, lapídibus óbrutus, migrávit ad Dóminum.
\switchcolumn
\selectlanguage{english}
At Nicomedia, the passion of Saints Migdonius and Mardonius, one of whom was 
 burned alive in the same persecution of Diocletian, and the other died in a 
 pit where he had been thrown. A deacon of St. Anthimus, bishop of 
 Nicomedia, suffered at the same time. He had been arrested by the 
 heathen when he was carrying letters to the martyrs, and being overwhelmed 
 with stones, went to our Lord.
\switchcolumn*
\selectlanguage{latin}
Ibídem natális sanctórum vigínti Mártyrum, quos 
 ípsamet Diocletiána persecútio, gravíssimis torméntis cruciátos, 
 Mártyres Christi fecit.
\switchcolumn
\selectlanguage{english}
Likewise, the birthday of twenty holy martyrs, whom the persecution of 
 Diocletian made martyrs for the faith of Christ, after subjecting them to 
 the most painful torments.
\switchcolumn*
\selectlanguage{latin}
In Creta sanctórum Mártyrum Theodúli, Saturníni, 
 Eupori, Gelásii, Euniciáni, Zétici, Leóminis, Agathópodis, Basílidis et 
 Evarísti; qui, in Décii persecutióne, crudélia passi sunt et cápite cæsi.
\switchcolumn
\selectlanguage{english}
In Crete, the holy martyrs Theodulus, Saturninus, Euporus, Gelasius, 
 Eunicianus, Zeticus, Leomines, Agathopodes, Basilides, and Everistus, who 
 were beheaded after suffering cruel torments in the persecution of Decius.
\switchcolumn*
\selectlanguage{latin}
Romæ beáti Sérvuli, qui (ut sanctus Gregórius 
 Papa scribit), a primæva sui ætáte usque ad finem vitæ, paralyticus jácuit 
 in pórticu prope Ecclésiam sancti Cleméntis, et demum, Angelórum cántibus 
 invitátus, ad paradísi glóriam transívit; ad cujus túmulum Deus mirácula 
 crebérrime osténdit.
\switchcolumn
\selectlanguage{english}
At Rome, blessed Servulus of whom St. Gregory writes that from his early 
 years to the end of his life he was a paralytic and had remained lying in a 
 porch near St. Clement's Church, and being invited by the chant of angels, 
 he went to enjoy the glory of Paradise. At his tomb frequent miracles 
 are wrought by God.
\switchcolumn*
\selectlanguage{latin}
\end{paracol}


% ---- martyrology/mart12/mart1224.htm
\needspace{10\baselineskip}
\begin{paracol}{2}
\selectlanguage{latin}
\begin{center}{\color{gregoriocolor} Nono Kaléndas Januárii. 
 Luna\dots\ }\end{center}
\switchcolumn
\selectlanguage{english}
\begin{center}{\color{gregoriocolor} The 
 24\textsuperscript{th} Day of 
 December. The\dots\ Day of the Moon.}\end{center}
\end{paracol}

\noindent\begin{tabularx}{\linewidth}{*{19}{>{\centering\arraybackslash}X}}
 \textcolor{gregoriocolor}{a} & \textcolor{gregoriocolor}{b} & \textcolor{gregoriocolor}{c} & \textcolor{gregoriocolor}{d} & \textcolor{gregoriocolor}{e} & \textcolor{gregoriocolor}{f} & \textcolor{gregoriocolor}{g} & \textcolor{gregoriocolor}{h} & \textcolor{gregoriocolor}{i} & \textcolor{gregoriocolor}{k} & \textcolor{gregoriocolor}{l} & \textcolor{gregoriocolor}{m} & \textcolor{gregoriocolor}{n} & \textcolor{gregoriocolor}{p} & \textcolor{gregoriocolor}{q} & \textcolor{gregoriocolor}{r} & \textcolor{gregoriocolor}{s} & \textcolor{gregoriocolor}{t} & \textcolor{gregoriocolor}{u} \\
 5 & 6 & 7 & 8 & 9 & 10 & 11 & 12 & 13 & 14 & 15 & 16 & 17 & 18 & 19 & 20 & 21 & 22 & 23 \\
\end{tabularx}
\vspace{0.5\baselineskip}
\noindent\begin{tabularx}{\linewidth}{*{12}{>{\centering\arraybackslash}X}}
 \textcolor{gregoriocolor}{A} & \textcolor{gregoriocolor}{B} & \textcolor{gregoriocolor}{C} & \textcolor{gregoriocolor}{D} & \textcolor{gregoriocolor}{E} & F & \textcolor{gregoriocolor}{F} & \textcolor{gregoriocolor}{G} & \textcolor{gregoriocolor}{H} & \textcolor{gregoriocolor}{M} & \textcolor{gregoriocolor}{N} & \textcolor{gregoriocolor}{P} \\
 24 & 25 & 26 & 27 & 28 & 29 & 28 & 29 & 1 & 2 & 3 & 4 \\
\end{tabularx}

\begin{paracol}{2}
\selectlanguage{latin}
\lettrine[lines=1]{V}{igília} Nativitátis Dómini nostri Jesu Christi.
\switchcolumn
\selectlanguage{english}
\lettrine[lines=1]{T}{he} Vigil of the Nativity of our Lord Jesus Christ.
\switchcolumn*
\selectlanguage{latin}
Cracóviæ, in Polónia, natális sancti Joánnis 
 Cántii, Presbyteri et Confessóris, quem, doctrína, propagándæ fídei zelo, 
 virtútibus ac miráculis clarum, Clemens Décimus tértius, Póntifex Máximus, 
 Sanctórum número adscrípsit. Ejus autem festívitas tertiodécimo 
 Kaléndas Novémbris celebrátur.
\switchcolumn
\selectlanguage{english}
At Cracow in Poland, the birthday of St. John Cantius, priest and confessor, 
 celebrated for his learning, for his zeal in propagating the faith, and for 
 his virtues and miracles, for which Pope Clement XIII added him to the 
 number of the saints. His feast is observed on the 20th of October.
\switchcolumn*
\selectlanguage{latin}
Apud Spolétum sancti Gregórii, Presbyteri et Mártyris; qui, tempóribus 
 Diocletiáni et Maximiáni Imperatórum, primo nodósis fústibus cæsus, 
 ac deínde, post cratículam et cárcerem, cárduis férreis in génibus percússus, 
 sed et ardéntibus lampádibus per látera incénsus, tandem est decollátus.
\switchcolumn
\selectlanguage{english}
At Spoleto, St. Gregory, priest and martyr. In the time of Emperors 
 Diocletian and Maximian, he was first beaten with rough clubs, exposed on 
 the gridiron and imprisoned, struck on the knees with iron carding 
 instruments, burned on the sides with firebrands, and finally beheaded.
\switchcolumn*
\selectlanguage{latin}
Trípoli, in Phœnícia, sanctórum Mártyrum Luciáni, Metróbii, Pauli, Zenóbii,
  Theótimi et Drusi.
\switchcolumn
\selectlanguage{english}
At Tripoli in Phoenicia, the holy martyrs Leucian, Metrobius, Paul, Zenobius, 
 Theotimus, and Drusus.
\switchcolumn*
\selectlanguage{latin}
Nicomedíæ sancti Euthymii Mártyris, qui, in 
 persecutióne Diocletiáni, cum multos ad martyrium præmisísset, et ipse, ense 
 tranfíxus, eos secútus est ad corónam.
\switchcolumn
\selectlanguage{english}
At Nicomedia, during the persecution of Diocletian, St. Euthymius, martyr, 
 who sent many before him to martyrdom, and being pierced with a sword, 
 followed them to share their crown.
\switchcolumn*
\selectlanguage{latin}
Antiochíæ natális sanctárum Vírginum 
 quadragínta, quæ, in Deciána persecutióne, per divérsa torménta martyrium 
 consummárunt.
\switchcolumn
\selectlanguage{english}
At Antioch, the birthday of forty holy virgins who suffered martyrdom by 
 divers torments in the Decian persecution.
\switchcolumn*
\selectlanguage{latin}
Burdígalæ sancti Delphíni Epíscopi, qui, 
 Theodósii témpore, cláruit sanctitáte.
\switchcolumn
\selectlanguage{english}
At Bordeaux, St. Delphinus, bishop, who was renowned for holiness in the 
 time of Theodosius.
\switchcolumn*
\selectlanguage{latin}
Romæ natális sanctæ Tharsíllæ Vírginis, ámitæ 
 sancti Gregórii Papæ, de qua ipse testátur quod in hora éxitus sui Jesum ad 
 se veniéntem víderit.
\switchcolumn
\selectlanguage{english}
At Rome, the birthday of the holy virgin Tharsilla, aunt of Pope St. 
 Gregory, who writes of her that at the hour of her death she saw Jesus 
 coming to her.
\switchcolumn*
\selectlanguage{latin}
Tréviris sanctæ Irmínæ Vírginis, fíliæ 
 Dagobérti Regis.
\switchcolumn
\selectlanguage{english}
At Treves, St. Irmina, virgin, daughter of King Dagobert.
\switchcolumn*
\selectlanguage{latin}
\end{paracol}


% ---- martyrology/mart12/mart1225.htm
\needspace{10\baselineskip}
\begin{paracol}{2}
\selectlanguage{latin}
\begin{center}{\color{gregoriocolor} Octávo Kaléndas Januárii. 
 Luna\dots\ }\end{center}
\switchcolumn
\selectlanguage{english}
\begin{center}{\color{gregoriocolor} The 
 25\textsuperscript{th} Day of 
 December. The\dots\ Day of the Moon.}\end{center}
\end{paracol}

\noindent\begin{tabularx}{\linewidth}{*{19}{>{\centering\arraybackslash}X}}
 \textcolor{gregoriocolor}{a} & \textcolor{gregoriocolor}{b} & \textcolor{gregoriocolor}{c} & \textcolor{gregoriocolor}{d} & \textcolor{gregoriocolor}{e} & \textcolor{gregoriocolor}{f} & \textcolor{gregoriocolor}{g} & \textcolor{gregoriocolor}{h} & \textcolor{gregoriocolor}{i} & \textcolor{gregoriocolor}{k} & \textcolor{gregoriocolor}{l} & \textcolor{gregoriocolor}{m} & \textcolor{gregoriocolor}{n} & \textcolor{gregoriocolor}{p} & \textcolor{gregoriocolor}{q} & \textcolor{gregoriocolor}{r} & \textcolor{gregoriocolor}{s} & \textcolor{gregoriocolor}{t} & \textcolor{gregoriocolor}{u} \\
 6 & 7 & 8 & 9 & 10 & 11 & 12 & 13 & 14 & 15 & 16 & 17 & 18 & 19 & 20 & 21 & 22 & 23 & 24 \\
\end{tabularx}
\vspace{0.5\baselineskip}
\noindent\begin{tabularx}{\linewidth}{*{12}{>{\centering\arraybackslash}X}}
 \textcolor{gregoriocolor}{A} & \textcolor{gregoriocolor}{B} & \textcolor{gregoriocolor}{C} & \textcolor{gregoriocolor}{D} & \textcolor{gregoriocolor}{E} & F & \textcolor{gregoriocolor}{F} & \textcolor{gregoriocolor}{G} & \textcolor{gregoriocolor}{H} & \textcolor{gregoriocolor}{M} & \textcolor{gregoriocolor}{N} & \textcolor{gregoriocolor}{P} \\
 25 & 26 & 27 & 28 & 29 & 30 & 29 & 1 & 2 & 3 & 4 & 5 \\
\end{tabularx}

\begin{paracol}{2}
\selectlanguage{latin}
\lettrine[lines=2]{A}{nno } a creatióne mundi, quando 
 in princípio Deus creávit cœlum et terram, quínquies millésimo centésimo 
 nonagésimo nono: A dilúvio autem, anno bis millésimo nongentésimo 
 quinquagésimo séptimo: A nativitáte Abrahæ, anno bis millésimo quintodécimo: 
 A Moyse et egréssu pópuli Israël de Ægypto, anno millésimo quingentésimo 
 décimo: Ab unctióne David in Regem, anno millésimo trigésimo secúndo; 
 Hebdómada sexagésima quinta, juxta Daniélis prophetíam: Olympíade centésima 
 nonagésima quarta: Ab urbe Roma cóndita, anno septingentésimo quinquagésimo 
 secúndo: Anno Impérii Octaviáni Augústi quadragésimo secúndo, toto Orbe in 
 pace compósito, sexta mundi ætáte, Jesus Christus ætérnus Deus, æterníque 
 Patris Fílius, mundum volens advéntu suo piíssimo consecráre, de Spíritu 
 Sancto concéptus, novémque post conceptiónem decúrsis ménsibus
	\rubric{(Hic vox elevatur, et omnes genua flectunt)}, in Béthlehem Judæ náscitur ex María Vírgine factus Homo.
\switchcolumn
\selectlanguage{english}
\lettrine[lines=2]{I}{n } the 5199th year 
 of the creation of the world, from the time when in the beginning God 
 created heaven and earth; from the flood, the 2957th year; from 
 the birth of Abraham, the 2015th year; from Moses and the 
 going-out of the people of Israel from Egypt, the 1510th year; from 
 the anointing of David as king, the 1032nd year; in the 65th 
 week according to the prophecy of Daniel; in the 194th Olympiad; 
 from the founding of the city of Rome, the 752nd year; in the 42nd 
 year of the rule of Octavian Augustus, when the whole world was at peace, in 
 the sixth age of the world: Jesus Christ, the eternal God and Son of the 
 eternal Father, desiring to sanctify the world by His most merciful coming, 
 having been conceived by the Holy Ghost, and nine months having passed since 
	His conception \rubric{(A higher tone of voice is now used, and all kneel)} was born in 
 Bethlehem of Juda of the Virgin Mary, having become man.
\switchcolumn*
\selectlanguage{latin}
	\rubric{Hic autem in priori voce dicitur, et in tono passionis:} NATÍVITAS Dómini nostri Jesu Christi secúndum carnem. \rubric{Quod sequitur, legitur in tono Lectionis consueto; et surgunt omnes.}
\switchcolumn
\selectlanguage{english}
	\rubric{In the same higher tone of voice and in the tone of the Passion:} THE NATIVITY of our Lord Jesus Christ according to the flesh. \rubric{That which follows is said in the customary tone of the Martyrology, and all arise.}
\switchcolumn*
\selectlanguage{latin}
Eódem die natális sanctæ Anastásiæ, quæ, 
 témpore Diocletiáni, primo diram et immítem custódiam a viro suo Públio 
 perpéssa est, in qua tamen a Confessóre Christi Chrysógono multum consoláta 
 et confortáta fuit; deínde a Floro, Præfécto Illyrici, per diútinam 
 custódiam maceráta, ad últimum, mánibus et pédibus exténsis, ligáta est ad 
 palos, et circa eam ignis accénsus, in quo martyrium consummávit in ínsula 
 Palmária, ad quam una cum ducéntis viris et septuagínta féminis deportáta 
 fúerat, qui váriis interfectiónibus martyrium celebrárunt.
\switchcolumn
\selectlanguage{english}
The same day, the birthday of St. Anastasia, who, in the time of Diocletian, 
 first suffered a severe and harsh imprisonment on the part of her husband 
 Publius, in which, however, she was much consoled and encouraged by the 
 confessor of Christ, Chrysogonus. Afterwards she was thrown into 
 prison again by order of Florus, prefect of Illyria; and finally, having her 
 hands and feet stretched, she was tied to stakes with a fire kindled about 
 her, in the midst of which she ended her martyrdom on the island of Palmaria, 
 whither she had been brought with two hundred men and seventy women, who 
 have made martyrdom a glorious thing by the various kinds of death they so 
 valiantly endured.
\switchcolumn*
\selectlanguage{latin}
Barcinóne, in Hispánia, item natális sancti Petri Nolásci Confessóris, qui 
 Fundátor éxstitit Ordinis beátæ Maríæ de 
 Mercéde redemptiónis captivórum, ac virtúte et miráculis cláruit. 
 Ipsíus autem festum cólitur quinto Kaléndas Februárii.
\switchcolumn
\selectlanguage{english}
At Barcelona in Spain, St. Peter Nolasco, confessor and founder of the Order 
 of our Lady of Ransom for the Redemption of Captives, renowned for virtue 
 and miracles. His feast is celebrated on the 28th of January.
\switchcolumn*
\selectlanguage{latin}
Romæ, in cœmetério Aproniáni, sanctæ Eugéniæ 
 Vírginis, beáti Mártyris Philíppi fíliæ, quæ, témpore Galliéni Imperatóris, 
 post plúrima virtútum insígnia, post sacros Vírginum choros Christo 
 aggregátos, sub Præfécto Urbis Nicétio diu agonizávit, ac novíssime gládio 
 juguláta est.
\switchcolumn
\selectlanguage{english}
At Rome, in the cemetery of Apronian, St. Eugenia, virgin, the daughter of 
 blessed Philip, martyr. In the time of Emperor Gallienus, after 
 displaying many signs and virtues, gathering to Christ holy choirs of 
 virgins, and after long trials under Nicetius, prefect of the city, she was 
 finally put to the sword.
\switchcolumn*
\selectlanguage{latin}
Nicomediæ pássio multórum míllium Mártyrum, qui 
 cum in Christi Natáli ad Domínicum conveníssent, Diocletiánus Imperátor 
 jánuas Ecclésiæ claudi jussit, et ignem circumcírca parári, tripodémque cum 
 thure præ fóribus poni, ac præcónem magna voce clamáre ut qui incéndium 
 vellent effúgere, foras exírent et Jovi thus adolérent; cumque omnes una 
 voce respondíssent se pro Christo libéntius mori, incénso igne consúmpti 
 sunt, atque ita eo die nasci meruérunt in cælis, quo Christus in terris pro 
 salúte mundi olim nasci dignátus est.
\switchcolumn
\selectlanguage{english}
At Nicomedia, many thousand martyrs, who had assembled for divine service on 
 our Lord's Nativity. When Emperor Diocletian ordered the doors of the 
 church to be closed, fire to kindled here and there, a vessel with incense 
 to be put before the entrance, and a man to cry out that those who wished to 
 escape from the fire should come out and burn incense to Jupiter, all with 
 one voice answered that they preferred to die for Christ. They were 
 consumed in the fire, and thus merited to be born in heaven on the day on 
 which Christ vouchsafed to be born on earth for the salvation of the world.
\switchcolumn*
\selectlanguage{latin}
\end{paracol}


% ---- martyrology/mart12/mart1226.htm
\needspace{10\baselineskip}
\begin{paracol}{2}
\selectlanguage{latin}
\begin{center}{\color{gregoriocolor} Séptimo Kaléndas Januárii. 
 Luna\dots\ }\end{center}
\switchcolumn
\selectlanguage{english}
\begin{center}{\color{gregoriocolor} The 
 26\textsuperscript{th} Day of 
 December. The\dots\ Day of the Moon.}\end{center}
\end{paracol}

\noindent\begin{tabularx}{\linewidth}{*{19}{>{\centering\arraybackslash}X}}
 \textcolor{gregoriocolor}{a} & \textcolor{gregoriocolor}{b} & \textcolor{gregoriocolor}{c} & \textcolor{gregoriocolor}{d} & \textcolor{gregoriocolor}{e} & \textcolor{gregoriocolor}{f} & \textcolor{gregoriocolor}{g} & \textcolor{gregoriocolor}{h} & \textcolor{gregoriocolor}{i} & \textcolor{gregoriocolor}{k} & \textcolor{gregoriocolor}{l} & \textcolor{gregoriocolor}{m} & \textcolor{gregoriocolor}{n} & \textcolor{gregoriocolor}{p} & \textcolor{gregoriocolor}{q} & \textcolor{gregoriocolor}{r} & \textcolor{gregoriocolor}{s} & \textcolor{gregoriocolor}{t} & \textcolor{gregoriocolor}{u} \\
 7 & 8 & 9 & 10 & 11 & 12 & 13 & 14 & 15 & 16 & 17 & 18 & 19 & 20 & 21 & 22 & 23 & 24 & 25 \\
\end{tabularx}
\vspace{0.5\baselineskip}
\noindent\begin{tabularx}{\linewidth}{*{12}{>{\centering\arraybackslash}X}}
 \textcolor{gregoriocolor}{A} & \textcolor{gregoriocolor}{B} & \textcolor{gregoriocolor}{C} & \textcolor{gregoriocolor}{D} & \textcolor{gregoriocolor}{E} & F & \textcolor{gregoriocolor}{F} & \textcolor{gregoriocolor}{G} & \textcolor{gregoriocolor}{H} & \textcolor{gregoriocolor}{M} & \textcolor{gregoriocolor}{N} & \textcolor{gregoriocolor}{P} \\
 26 & 27 & 28 & 29 & 30 & 1 & 1 & 2 & 3 & 4 & 5 & 6 \\
\end{tabularx}

\begin{paracol}{2}
\selectlanguage{latin}
\lettrine[lines=2]{H}{ierosólymis} natális sancti Stéphani Protomártyris, qui a Judæis, 
 non longe post Ascensiónem Dómini, lapidátus est.
\switchcolumn
\selectlanguage{english}
\lettrine[lines=2]{A}{t} Jerusalem, the birthday of St. Stephen, the first martyr, who was stoned 
 to death by the Jews shortly after the Ascension of our Lord.
\switchcolumn*
\selectlanguage{latin}
Romæ sancti Maríni, ex órdine Senatório viri, 
 qui, sub Numeriáno Imperatóre et Marciáno Præfécto, Christiánæ religiónis 
 causa comprehénsus, equúleo et úngulis servíli more punítus, in sartáginem 
 deínde conjéctus, sed, igne in rorem convérso, liberátus, objéctus quoque 
 feris et ab illis nullátenus læsus; tandem, ad aram íterum ductus, et, cum 
 idóla oratióne ejus corruíssent, percússus gládio, martyrii triúmphum 
 adéptus est.
\switchcolumn
\selectlanguage{english}
At Rome, St. Marinus, a man of senatorial rank. In the time of Emperor 
 Numerian and the prefect Marcian, he was arrested for the Christian 
 religion, racked and torn with iron claws like a slave, then thrown into a 
 boiling cauldron; but being delivered because the fire became like a dew, he 
 was exposed to the beasts without being injured by them, and finally being 
 led to the altar, the idols of which toppled over at his prayer, he was 
 struck with the sword, and thus obtained the triumph of martyrs.
\switchcolumn*
\selectlanguage{latin}
Ibídem, via Appia, deposítio sancti Dionysii Papæ, 
 qui, multis pro Ecclésia impénsis labóribus, fídei documéntis clarus 
 effúlsit.
\switchcolumn
\selectlanguage{english}
Likewise at Rome, on the Appian Way, the death of Pope St. Dionysius, who 
 sustained many labours for the Church, and was renowned for his doctrinal 
 writings.
\switchcolumn*
\selectlanguage{latin}
Item Romæ sancti Zósimi, Papæ et Confessóris.
\switchcolumn
\selectlanguage{english}
In the same city, St. Zosimus, pope and confessor.
\switchcolumn*
\selectlanguage{latin}
In Mesopotámia sancti Archelái Epíscopi, doctrína et sanctitáte célebris.
\switchcolumn
\selectlanguage{english}
In Mesopotamia, St. Archelaus, bishop, famous for learning and holiness.
\switchcolumn*
\selectlanguage{latin}
Majúmæ, in Palæstína, sancti Zenónis Epíscopi.
\switchcolumn
\selectlanguage{english}
At Majuma, in Palestine, St. Zeno, bishop.
\switchcolumn*
\selectlanguage{latin}
Romæ sancti Theodóri, qui Mansionárius Ecclésiæ 
 sancti Petri fuit, cujus et méminit beátus Gregórius Papa.
\switchcolumn
\selectlanguage{english}
At Rome, St. Theodore, sacristan of the church of St. Peter, who is 
 mentioned by blessed Pope Gregory.
\switchcolumn*
\selectlanguage{latin}
\end{paracol}


% ---- martyrology/mart12/mart1227.htm
\needspace{10\baselineskip}
\begin{paracol}{2}
\selectlanguage{latin}
\begin{center}{\color{gregoriocolor} Sexto Kaléndas Januárii. 
 Luna\dots\ }\end{center}
\switchcolumn
\selectlanguage{english}
\begin{center}{\color{gregoriocolor} The 
 27\textsuperscript{th} Day of 
 December. The\dots\ Day of the Moon.}\end{center}
\end{paracol}

\noindent\begin{tabularx}{\linewidth}{*{19}{>{\centering\arraybackslash}X}}
 \textcolor{gregoriocolor}{a} & \textcolor{gregoriocolor}{b} & \textcolor{gregoriocolor}{c} & \textcolor{gregoriocolor}{d} & \textcolor{gregoriocolor}{e} & \textcolor{gregoriocolor}{f} & \textcolor{gregoriocolor}{g} & \textcolor{gregoriocolor}{h} & \textcolor{gregoriocolor}{i} & \textcolor{gregoriocolor}{k} & \textcolor{gregoriocolor}{l} & \textcolor{gregoriocolor}{m} & \textcolor{gregoriocolor}{n} & \textcolor{gregoriocolor}{p} & \textcolor{gregoriocolor}{q} & \textcolor{gregoriocolor}{r} & \textcolor{gregoriocolor}{s} & \textcolor{gregoriocolor}{t} & \textcolor{gregoriocolor}{u} \\
 8 & 9 & 10 & 11 & 12 & 13 & 14 & 15 & 16 & 17 & 18 & 19 & 20 & 21 & 22 & 23 & 24 & 25 & 26 \\
\end{tabularx}
\vspace{0.5\baselineskip}
\noindent\begin{tabularx}{\linewidth}{*{12}{>{\centering\arraybackslash}X}}
 \textcolor{gregoriocolor}{A} & \textcolor{gregoriocolor}{B} & \textcolor{gregoriocolor}{C} & \textcolor{gregoriocolor}{D} & \textcolor{gregoriocolor}{E} & F & \textcolor{gregoriocolor}{F} & \textcolor{gregoriocolor}{G} & \textcolor{gregoriocolor}{H} & \textcolor{gregoriocolor}{M} & \textcolor{gregoriocolor}{N} & \textcolor{gregoriocolor}{P} \\
 27 & 28 & 29 & 30 & 1 & 2 & 2 & 3 & 4 & 5 & 6 & 7 \\
\end{tabularx}

\begin{paracol}{2}
\selectlanguage{latin}
\lettrine[lines=2]{A}{pud} Ephesum natális sancti Joánnis, Apóstoli et Evangelístæ, 
 qui, post Evangélii scriptiónem, post exsílii relegatiónem et Apocalypsim 
 divínam, usque ad Trajáni Príncipis témpora persevérans, totíus Asiæ 
 fundávit rexítque Ecclésias, ac tandem, conféctus sénio, sexagésimo octávo 
 post passiónem Dómini anno mórtuus est, et juxta eándem urbem sepúltus.
\switchcolumn
\selectlanguage{english}
\lettrine[lines=2]{A}{t} Ephesus, the birthday of St. John, apostle and evangelist. After 
 writing his gospel, and after enduring exile and writing the divine 
 Apocalypse, he lived until the time of Emperor Trajan and founded and 
 governed the churches of all Asia. Worn out with age, he died in the 
 sixty-eighth year after the passion of our Lord and was buried near Ephesus.
\switchcolumn*
\selectlanguage{latin}
Constantinópoli sanctórum Confessórum Theodóri et Theóphanis fratrum, qui, a 
 puerítia in Palæstinénsi sancti Sabbæ 
 monastério nutríti, cum póstea pro sanctárum Imáginum cultu advérsus Leónem 
 Arménum strénue decertárent, ejus jussu verbéribus affécti sunt et exsílio 
 relegáti. Sed, eódem Leóne mórtuo, rursus Theóphilo Imperatóri, qui 
 eádem impietáte detinebátur, constánter resisténtes, verbéribus íterum cæsi 
 et in exsílium pulsi sunt, ubi Theodórus in cárcere exspirávit. 
 Theóphanes vero, pace demum Ecclésiæ réddita, factus est Nicænæ civitátis 
 Epíscopus, et confessiónis glória præclárus quiévit in Dómino.
\switchcolumn
\selectlanguage{english}
At Constantinople, the holy confessors Theodore and Theophanes, brothers, 
 who were brought up from their childhood in the monastery of St. Sabas. 
 Afterwards, they strove zealously for the veneration of holy images against 
 Leo the Armenian, and at his command they were scourged and banished. 
 After his death they again firmly opposed Emperor Theophilus, who was imbued 
 with the same impiety, and were scourged a second time and driven into exile, where Theodore died in prison. Theophanes, after peace had at 
 length been restored to the Church, was made bishop of Nicaea, and there, 
 famous for his glorious witness of the faith, rested in the Lord.
\switchcolumn*
\selectlanguage{latin}
Alexandríæ sancti Máximi Epíscopi, qui satis 
 clarus et insígnis título confessiónis efféctus est.
\switchcolumn
\selectlanguage{english}
At Alexandria, St. Maximus, bishop, well known and renowned by reason of his 
 confession.
\switchcolumn*
\selectlanguage{latin}
Constantinópoli sanctæ Nicáretes Vírginis, quæ, 
 sub Arcádio Imperatóre, cláruit sanctitáte.
\switchcolumn
\selectlanguage{english}
At Constantinople, St. Niceras, virgin, who was renowned for sanctity in the 
 time of Emperor Arcadius.
\switchcolumn*
\selectlanguage{latin}
\end{paracol}


% ---- martyrology/mart12/mart1228.htm
\needspace{10\baselineskip}
\begin{paracol}{2}
\selectlanguage{latin}
\begin{center}{\color{gregoriocolor} Quinto Kaléndas Januárii. 
 Luna\dots\ }\end{center}
\switchcolumn
\selectlanguage{english}
\begin{center}{\color{gregoriocolor} The 
 28\textsuperscript{th} Day of 
 December. The\dots\ Day of the Moon.}\end{center}
\end{paracol}

\noindent\begin{tabularx}{\linewidth}{*{19}{>{\centering\arraybackslash}X}}
 \textcolor{gregoriocolor}{a} & \textcolor{gregoriocolor}{b} & \textcolor{gregoriocolor}{c} & \textcolor{gregoriocolor}{d} & \textcolor{gregoriocolor}{e} & \textcolor{gregoriocolor}{f} & \textcolor{gregoriocolor}{g} & \textcolor{gregoriocolor}{h} & \textcolor{gregoriocolor}{i} & \textcolor{gregoriocolor}{k} & \textcolor{gregoriocolor}{l} & \textcolor{gregoriocolor}{m} & \textcolor{gregoriocolor}{n} & \textcolor{gregoriocolor}{p} & \textcolor{gregoriocolor}{q} & \textcolor{gregoriocolor}{r} & \textcolor{gregoriocolor}{s} & \textcolor{gregoriocolor}{t} & \textcolor{gregoriocolor}{u} \\
 9 & 10 & 11 & 12 & 13 & 14 & 15 & 16 & 17 & 18 & 19 & 20 & 21 & 22 & 23 & 24 & 25 & 26 & 27 \\
\end{tabularx}
\vspace{0.5\baselineskip}
\noindent\begin{tabularx}{\linewidth}{*{12}{>{\centering\arraybackslash}X}}
 \textcolor{gregoriocolor}{A} & \textcolor{gregoriocolor}{B} & \textcolor{gregoriocolor}{C} & \textcolor{gregoriocolor}{D} & \textcolor{gregoriocolor}{E} & F & \textcolor{gregoriocolor}{F} & \textcolor{gregoriocolor}{G} & \textcolor{gregoriocolor}{H} & \textcolor{gregoriocolor}{M} & \textcolor{gregoriocolor}{N} & \textcolor{gregoriocolor}{P} \\
 28 & 29 & 30 & 1 & 2 & 3 & 3 & 4 & 5 & 6 & 7 & 8 \\
\end{tabularx}

\begin{paracol}{2}
\selectlanguage{latin}
\lettrine[lines=2]{I}{n} Béthlehem Judæ natális sanctórum Innocéntium 
 Mártyrum, qui pro Christo ab Heróde Rege interfécti sunt.
\switchcolumn
\selectlanguage{english}
\lettrine[lines=2]{I}{n} Bethlehem of Juda, the birthday of the Holy Innocents, who were slain for 
 Christ by Herod the king.
\switchcolumn*
\selectlanguage{latin}
Lugdúni, in Gállia, item natális sancti Francísci Salésii, Epíscopi Gebennénsis
  et Confessóris; quem, doctrína et flagrantíssimo in converténdis hæréticis
  zelo præclárum, Alexánder Papa Séptimus in Sanctórum númerum rétulit, et
  ipsíus festivitátem quarto Kaléndas Februárii, quo die sacrum illíus corpus
  Lugdúno Annésium, in Sabáudia, fuit translátum, agéndam esse constítuit. Eum
  Pius Nonus, Póntifex Máximus, Doctórem universális Ecclésiæ declarávit; et
  Pius Papa Undécimus ómnibus Scriptóribus cathólicis, qui diáriis aliísve
  scriptis in vulgus edéndis Christiánam sapiéntiam illústrant ac próvehunt et
  tuéntur, cæléstem Patrónum dedit seu confirmávit.
\switchcolumn
\selectlanguage{english}
At Lyons in France, the birthday also of St. Francis de Sales, bishop of 
 Geneva and confessor. Because of his burning zeal for the conversion 
 of heretics and his learning, Pope Alexander VII placed him among the number 
 of the saints, and his feast is observed on the 29th of January, on which 
 day his holy body was translated from Lyons to Annecy in Savoy. Pope 
 Pius IX decreed him a doctor of the universal Church, and Pope Pius XI 
 constituted him the heavenly patron of all Catholic writers who explain, 
 promote, or defend Christian doctrine by publishing journals or other 
 writings in the vernacular.
\switchcolumn*
\selectlanguage{latin}
Ancyræ, in Galátia, sanctórum Mártyrum Eutychii 
 Presbyteri, et Domitiáni Diáconi.
\switchcolumn
\selectlanguage{english}
At Ancyra in Galatia, the holy martyrs Eutychius, priest, and Domitian, 
 deacon.
\switchcolumn*
\selectlanguage{latin}
In Africa natális sanctórum Mártyrum Cástoris, 
 Victóris et Rogatiáni.
\switchcolumn
\selectlanguage{english}
In Africa, the birthday of the holy martyrs Castor, Victor, and Rogatian.
\switchcolumn*
\selectlanguage{latin}
Nicomedíæ sanctórum Mártyrum Indis eunúchi, Domnæ et Agapis ac Theóphilæ
  Vírginum, et Sociórum; qui, in persecutióne Diocletiáni, post longa
  certámina, divérso mortis génere corónam martyrii sunt assecúti.
\switchcolumn
\selectlanguage{english}
At Nicomedia, the holy martyrs Indes, a eunuch, Domna, Agapes, and Theophila, 
 virgins, and their companions, who, after long trials, attained to the crown 
 of martyrdom by various kinds of death, during the persecution of 
 Diocletian.
\switchcolumn*
\selectlanguage{latin}
Neocæsaréæ, in Ponto, sancti Troádii Mártyris, 
 in persecutióne Décii; cui quidem Troádio agonizánti sanctus Gregórius 
 Thaumatúrgus in spíritu ádfuit, eúmque ad subeúndum martyrium 
 roborávit.
\switchcolumn
\selectlanguage{english}
At Neocaesarea in Pontus, St. Troadius, martyr, in the persecution of Decius. 
 During his trial St. Gregory Thaumaturgus appeared to him in spirit and 
 encouraged him to undergo martyrdom.
\switchcolumn*
\selectlanguage{latin}
Arabíssi, in Arménia inferióre, sancti Cæsárii 
 Mártyris, qui sub Galério Maximiáno passus est.
\switchcolumn
\selectlanguage{english}
At Arabissus in Lower Armenia, St. Caesarius, martyr, who suffered under 
 Galerius Maximian.
\switchcolumn*
\selectlanguage{latin}
Romæ sancti Domniónis Presbyteri.
\switchcolumn
\selectlanguage{english}
At Rome, St. Domnio, priest.
\switchcolumn*
\selectlanguage{latin}
In monastério Lirinénsi, in Gállia, sancti Antónii Mónachi, miráculis clari.
\switchcolumn
\selectlanguage{english}
In the monastery of Lerins in France, St. Anthony, a monk famed for his 
 miracles.
\switchcolumn*
\selectlanguage{latin}
\end{paracol}


% ---- martyrology/mart12/mart1229.htm
\needspace{10\baselineskip}
\begin{paracol}{2}
\selectlanguage{latin}
\begin{center}{\color{gregoriocolor} Quarto Kaléndas Januárii. 
 Luna\dots\ }\end{center}
\switchcolumn
\selectlanguage{english}
\begin{center}{\color{gregoriocolor} The 
 29\textsuperscript{th} Day of 
 December. The\dots\ Day of the Moon.}\end{center}
\end{paracol}

\noindent\begin{tabularx}{\linewidth}{*{19}{>{\centering\arraybackslash}X}}
 \textcolor{gregoriocolor}{a} & \textcolor{gregoriocolor}{b} & \textcolor{gregoriocolor}{c} & \textcolor{gregoriocolor}{d} & \textcolor{gregoriocolor}{e} & \textcolor{gregoriocolor}{f} & \textcolor{gregoriocolor}{g} & \textcolor{gregoriocolor}{h} & \textcolor{gregoriocolor}{i} & \textcolor{gregoriocolor}{k} & \textcolor{gregoriocolor}{l} & \textcolor{gregoriocolor}{m} & \textcolor{gregoriocolor}{n} & \textcolor{gregoriocolor}{p} & \textcolor{gregoriocolor}{q} & \textcolor{gregoriocolor}{r} & \textcolor{gregoriocolor}{s} & \textcolor{gregoriocolor}{t} & \textcolor{gregoriocolor}{u} \\
 10 & 11 & 12 & 13 & 14 & 15 & 16 & 17 & 18 & 19 & 20 & 21 & 22 & 23 & 24 & 25 & 26 & 27 & 28 \\
\end{tabularx}
\vspace{0.5\baselineskip}
\noindent\begin{tabularx}{\linewidth}{*{12}{>{\centering\arraybackslash}X}}
 \textcolor{gregoriocolor}{A} & \textcolor{gregoriocolor}{B} & \textcolor{gregoriocolor}{C} & \textcolor{gregoriocolor}{D} & \textcolor{gregoriocolor}{E} & F & \textcolor{gregoriocolor}{F} & \textcolor{gregoriocolor}{G} & \textcolor{gregoriocolor}{H} & \textcolor{gregoriocolor}{M} & \textcolor{gregoriocolor}{N} & \textcolor{gregoriocolor}{P} \\
 29 & 30 & 1 & 2 & 3 & 4 & 4 & 5 & 6 & 7 & 8 & 9 \\
\end{tabularx}

\begin{paracol}{2}
\selectlanguage{latin}
\lettrine[lines=2]{C}{antuáriæ,} in Anglia, natális sancti Thomæ, 
 Epíscopi et Mártyris, qui, ob defensiónem justítiæ et ecclesiásticæ 
 immunitátis, in Basílica sua, ab impiórum hóminum factióne percússus gládio, 
 Martyr migrávit ad Christum.
\switchcolumn
\selectlanguage{english}
\lettrine[lines=2]{A}{t} Canterbury in England, the birthday of St. Thomas, bishop and martyr, 
 who, for the defence of justice and ecclesiastical immunity, was struck with 
 the sword in his own basilica by a faction of wicked men, and thus went to 
 Christ as martyr.
\switchcolumn*
\selectlanguage{latin}
Hierosólymis sancti David, Regis et Prophétæ.
\switchcolumn
\selectlanguage{english}
At Jerusalem, holy David, king and prophet.
\switchcolumn*
\selectlanguage{latin}
Areláte, in Gállia, natális sancti Tróphimi, 
 cujus méminit sanctus Paulus ad Timótheum scribens. Ipse autem 
 Tróphimus, ab eódem Apóstolo Epíscopus ordinátus, præfátæ urbi primus ad 
 Christi Evangélium prædicándum diréctus est; ex cujus prædicatiónis fonte (ut 
 sanctus Zósimus Papa scribit) tota Gállia rívulos fídei recépit.
\switchcolumn
\selectlanguage{english}
At Arles in France, the birthday of St. Trophimus, mentioned by St. Paul in 
 his Epistle to Timothy. Being ordained bishop by that apostle, he was 
 the first sent to preach the gospel of Christ in that city. From his 
 preaching, as from a fountain, according to the expression of Pope St. 
 Zosimus, all France received the waters of salvation.
\switchcolumn*
\selectlanguage{latin}
Romæ sanctórum Mártyrum Callísti, Felícis et Bonifátii.
\switchcolumn
\selectlanguage{english}
At Rome, the holy martyrs Callistus, Felix, and Boniface.
\switchcolumn*
\selectlanguage{latin}
In Africa pássio sanctórum Mártyrum Domínici, Victóris, Primiáni, Lybósi, 
 Saturníni, Crescéntii, Secúndi et Honoráti.
\switchcolumn
\selectlanguage{english}
In Africa, the passion of the holy martyrs Dominic, Victor, Primian, Lybosus, 
 Saturninus, Crescentius, Secundus, and Honoratus.
\switchcolumn*
\selectlanguage{latin}
Constantinópoli sancti Marcélli Abbátis.
\switchcolumn
\selectlanguage{english}
At Constantinople, St. Marcellus, abbot.
\switchcolumn*
\selectlanguage{latin}
In pago Oxyménsi, in Gállia, sancti Ebrúlphi, Abbátis et Confessóris, 
 témpore Childebérti Regis.
\switchcolumn
\selectlanguage{english}
In the country of Hiesmes in France, St. Ebrulf, abbot and confessor, in the 
 time of King Childebert.
\switchcolumn*
\selectlanguage{latin}
Viénnæ, in Gállia, Commemorátio sancti 
 Crescéntis, Epíscopi et Mártyris, qui fuit discípulus beáti Pauli Apóstoli 
 ac primus ejúsdem civitátis Epíscopus, et cujus dies natális quinto Kaléndas 
 Júlii celebrátur.
\switchcolumn
\selectlanguage{english}
At Vienne in France, the commemoration of St. Crescens, bishop and martyr. 
 He was a disciple of St. Paul the Apostle and was the first bishop of that 
 city. His birthday is mentioned on the 27th of June.
\switchcolumn*
\selectlanguage{latin}
\end{paracol}


% ---- martyrology/mart12/mart1230.htm
\needspace{10\baselineskip}
\begin{paracol}{2}
\selectlanguage{latin}
\begin{center}{\color{gregoriocolor} Tértio Kaléndas Januárii. 
 Luna\dots\ }\end{center}
\switchcolumn
\selectlanguage{english}
\begin{center}{\color{gregoriocolor} The 
 30\textsuperscript{th} Day of 
 December. The\dots\ Day of the Moon.}\end{center}
\end{paracol}

\noindent\begin{tabularx}{\linewidth}{*{19}{>{\centering\arraybackslash}X}}
 \textcolor{gregoriocolor}{a} & \textcolor{gregoriocolor}{b} & \textcolor{gregoriocolor}{c} & \textcolor{gregoriocolor}{d} & \textcolor{gregoriocolor}{e} & \textcolor{gregoriocolor}{f} & \textcolor{gregoriocolor}{g} & \textcolor{gregoriocolor}{h} & \textcolor{gregoriocolor}{i} & \textcolor{gregoriocolor}{k} & \textcolor{gregoriocolor}{l} & \textcolor{gregoriocolor}{m} & \textcolor{gregoriocolor}{n} & \textcolor{gregoriocolor}{p} & \textcolor{gregoriocolor}{q} & \textcolor{gregoriocolor}{r} & \textcolor{gregoriocolor}{s} & \textcolor{gregoriocolor}{t} & \textcolor{gregoriocolor}{u} \\
 11 & 12 & 13 & 14 & 15 & 16 & 17 & 18 & 19 & 20 & 21 & 22 & 23 & 24 & 25 & 26 & 27 & 28 & 29 \\
\end{tabularx}
\vspace{0.5\baselineskip}
\noindent\begin{tabularx}{\linewidth}{*{12}{>{\centering\arraybackslash}X}}
 \textcolor{gregoriocolor}{A} & \textcolor{gregoriocolor}{B} & \textcolor{gregoriocolor}{C} & \textcolor{gregoriocolor}{D} & \textcolor{gregoriocolor}{E} & F & \textcolor{gregoriocolor}{F} & \textcolor{gregoriocolor}{G} & \textcolor{gregoriocolor}{H} & \textcolor{gregoriocolor}{M} & \textcolor{gregoriocolor}{N} & \textcolor{gregoriocolor}{P} \\
 30 & 1 & 2 & 3 & 4 & 5 & 5 & 6 & 7 & 8 & 9 & 10 \\
\end{tabularx}

\begin{paracol}{2}
\selectlanguage{latin}
\lettrine[lines=2]{R}{omæ} natális sancti Felícis Primi, Papæ et 
 Mártyris, qui sub Aureliáno Príncipe Ecclésiam rexit. Ipsíus tamen 
 festum tértio Kaléndas Júnii celebrátur.
\switchcolumn
\selectlanguage{english}
\lettrine[lines=2]{A}{t} Rome, the birthday of St. Felix I, pope and martyr, who governed the 
 Church during the reign of Emperor Aurelian. His feast day is 
 celebrated on the 30th of May.
\switchcolumn*
\selectlanguage{latin}
Spoléti item natális sanctórum Mártyrum Sabíni, Assisiénsis Epíscopi, atque 
 Exsuperántii et Marcélli Diaconórum, ac Venustiáni Præsidis 
 cum uxóre et fíliis, sub Maximiáno Imperatóre. Ex ipsis Marcéllus et Exsuperántius, primum equúleo suspénsi, deínde fústibus gráviter mactáti, 
 postrémum, abrási úngulis et láterum exustióne assáti, martyrium 
 complevérunt; Venustiánus autem non multo post, simul cum uxóre et fíliis, 
 est gládio necátus; sanctus vero Sabínus, post detruncatiónem mánuum et 
 diútinam cárceris maceratiónem, ad mortem usque cæsus est. Horum 
 martyrium, licet divérso exstíterit témpore, una tamen die recólitur.
\switchcolumn
\selectlanguage{english}
At Spoleto, the birthday also of the holy martyrs Sabinus, bishop, 
 Exuperantius and Marcellus, deacons, and also Venustian, governor, along 
 with his wife and sons, under Emperor Maximian. Marcellus and 
 Exuperantius were first racked, then severely beaten with rods; afterwards 
 being torn with iron hooks, and burned in the sides, they fulfilled their 
 martyrdom. Not long after, Venustian was put to the sword with his 
 wife and sons. St. Sabinus, after having his hands cut off, and being 
 a long time confined in prison, was scourged to death. The martyrdom 
 of these saints is commemorated on the same day, although it occurred at 
 different times.
\switchcolumn*
\selectlanguage{latin}
Alexandríæ sanctórum Mansuéti, Sevéri, Appiáni, 
 Donáti, Honórii et Sociórum Mártyrum.
\switchcolumn
\selectlanguage{english}
At Alexandria, the Saints Mansuetus, Severus, Appian, Donatus, Honorius, and 
 their martyr companions.
\switchcolumn*
\selectlanguage{latin}
Thessalonícæ sanctæ Anysiæ Mártyris.
\switchcolumn
\selectlanguage{english}
At Thessalonica, St. Anysia, martyr.
\switchcolumn*
\selectlanguage{latin}
Ibídem sancti Anysii, ejúsdem civitátis Epíscopi.
\switchcolumn
\selectlanguage{english}
Likewise, St. Anysius, bishop of the same city.
\switchcolumn*
\selectlanguage{latin}
Medioláni sancti Eugénii, Epíscopi et Confessóris.
\switchcolumn
\selectlanguage{english}
At Milan, St. Eugene, bishop and confessor.
\switchcolumn*
\selectlanguage{latin}
Ravénnæ sancti Libérii Epíscopi.
\switchcolumn
\selectlanguage{english}
At Ravenna, St. Liberius, bishop.
\switchcolumn*
\selectlanguage{latin}
Aquilæ, in Vestínis, sancti Rainérii Epíscopi.
\switchcolumn
\selectlanguage{english}
At Aquila, in Abruzzi, St. Rainer, bishop.
\switchcolumn*
\selectlanguage{latin}
\end{paracol}


% ---- martyrology/mart12/mart1231.htm
\needspace{10\baselineskip}
\begin{paracol}{2}
\selectlanguage{latin}
\begin{center}{\color{gregoriocolor} Prídie Kaléndas Januárii. 
 Luna\dots\ }\end{center}
\switchcolumn
\selectlanguage{english}
\begin{center}{\color{gregoriocolor} The 31\textsuperscript{st} Day of 
 December. The\dots\ Day of the Moon.}\end{center}
\end{paracol}

\noindent\begin{tabularx}{\linewidth}{*{19}{>{\centering\arraybackslash}X}}
 \textcolor{gregoriocolor}{a} & \textcolor{gregoriocolor}{b} & \textcolor{gregoriocolor}{c} & \textcolor{gregoriocolor}{d} & \textcolor{gregoriocolor}{e} & \textcolor{gregoriocolor}{f} & \textcolor{gregoriocolor}{g} & \textcolor{gregoriocolor}{h} & \textcolor{gregoriocolor}{i} & \textcolor{gregoriocolor}{k} & \textcolor{gregoriocolor}{l} & \textcolor{gregoriocolor}{m} & \textcolor{gregoriocolor}{n} & \textcolor{gregoriocolor}{p} & \textcolor{gregoriocolor}{q} & \textcolor{gregoriocolor}{r} & \textcolor{gregoriocolor}{s} & \textcolor{gregoriocolor}{t} & \textcolor{gregoriocolor}{u} \\
 12 & 13 & 14 & 15 & 16 & 17 & 18 & 19 & 20 & 21 & 22 & 23 & 24 & 25 & 26 & 27 & 28 & 29 & 30 \\
\end{tabularx}
\vspace{0.5\baselineskip}
\noindent\begin{tabularx}{\linewidth}{*{12}{>{\centering\arraybackslash}X}}
 \textcolor{gregoriocolor}{A} & \textcolor{gregoriocolor}{B} & \textcolor{gregoriocolor}{C} & \textcolor{gregoriocolor}{D} & \textcolor{gregoriocolor}{E} & F & \textcolor{gregoriocolor}{F} & \textcolor{gregoriocolor}{G} & \textcolor{gregoriocolor}{H} & \textcolor{gregoriocolor}{M} & \textcolor{gregoriocolor}{N} & \textcolor{gregoriocolor}{P} \\
 1 & 2 & 3 & 4 & 5 & 6 & 6 & 7 & 8 & 9 & 10 & 11 \\
\end{tabularx}

\begin{paracol}{2}
\selectlanguage{latin}
\lettrine[lines=2]{R}{omæ} natális sancti Silvéstri Primi, Papæ et 
 Confessóris; qui Magnum Constantínum Imperatórem baptizávit, et Nicænam 
 confirmávit Synodum, ac, multis áliis rebus sanctíssime gestis, quiévit in 
 pace.
\switchcolumn
\selectlanguage{english}
\lettrine[lines=2]{A}{t} Rome, the birthday of Pope St. Sylvester I, confessor, who baptized 
 Emperor Constantine the Great, and confirmed the council of Nicaea. 
 After performing many other holy deeds, he rested in peace.
\switchcolumn*
\selectlanguage{latin}
Item Romæ, via Salária, in cœmetério Priscíllæ, 
 sanctárum Mártyrum Donátæ, Paulínæ, Rústicæ, Nominándæ, Serótinæ, Hiláriæ et 
 Sociárum.
\switchcolumn
\selectlanguage{english}
At Rome, on the Salarian Way, in the cemetery of Priscilla, the holy martyrs 
 Donata, Paulina, Rustica, Nominanda, Serotina, Hilaria, and their 
 companions.
\switchcolumn*
\selectlanguage{latin}
Apud Sénonas beatórum Sabiniáni Epíscopi, et 
 Potentiáni; qui, a Pontífice Románo illuc ad prædicándum dirécti, 
 confessiónis suæ martyrio eándem metrópolim illustrárunt.
\switchcolumn
\selectlanguage{english}
At Sens, the blessed Sabinian, bishop, and Potentian. They had been 
 sent there to preach by the Roman Pontiff, and that metropolitan church was 
 illustrated by their confession and martyrdom.
\switchcolumn*
\selectlanguage{latin}
Cátanæ, in Sicília, pássio sanctórum Stéphani, 
 Pontiáni, Attali, Fabiáni, Cornélii, Sexti, Floris, Quinctiáni, Minervíni et 
 Simpliciáni.
\switchcolumn
\selectlanguage{english}
At Catania in Sicily, the passion of the Saints Stephen, Pontian, Attalus, 
 Fabian, Cornelius, Sextus, Flos, Quinctian, Minervinus, and Simplician.
\switchcolumn*
\selectlanguage{latin}
Apud Sénonas sanctæ Colúmbæ, Vírginis et 
 Mártyris; quæ, igne superáto, in persecutióne Aureliáni Imperatóris, gládio 
 cæsa est.
\switchcolumn
\selectlanguage{english}
At Sens, St. Columba, virgin and martyr, who, after having triumphed over 
 fire, was beheaded during the persecution of Emperor Aurelian.
\switchcolumn*
\selectlanguage{latin}
Eódem die sancti Zótici, Presbyteri Románi; 
 qui, Constantinópolim proféctus, illic alendórum orphanórum curam suscépit.
\switchcolumn
\selectlanguage{english}
On the same day, St. Zoticus, a Roman priest who went to Constantinople and 
 undertook the work of caring for orphans.
\switchcolumn*
\selectlanguage{latin}
Ravénnæ sancti Barbatiáni Presbyteri et 
 Confessóris.
\switchcolumn
\selectlanguage{english}
At Ravenna, St. Barbatian, priest and confessor.
\switchcolumn*
\selectlanguage{latin}
In pago Lalovésci, diœcésis Viennénsis, in Delphinátu, deposítio sancti
Joánnis-Francísci Regis, Sacerdótis e Societáte Jesu et Confessóris, exímiæ in
salúte animárum procuránda caritátis et patiéntiæ viri; quem Clemens Papa
Duodécimus in Sanctórum cánonem rétulit.
\switchcolumn
\selectlanguage{english}
At La Louvesc, in the diocese of Vienne in Dauphine, the death of St. John 
 Francis Regis, priest of the Society of Jesus and confessor. He was a 
 man of great love and patience in securing the salvation of souls.
\switchcolumn*
\selectlanguage{latin}

\selectlanguage{latin}
\rubric{\P} Et álibi aliórum plurimórum sanctórum Mártyrum et Confessórum, atque sanctárum Vírginum.

\textcolor{gregoriocolor}{\Rbar.}~Deo grátias.
\switchcolumn
\selectlanguage{english}
\rubric{\P} And elsewhere in divers places, many other holy martyrs, confessors, and holy virgins.

\textcolor{gregoriocolor}{\Rbar.}~Thanks be to God.
\switchcolumn*

\end{paracol}
